\chapter{多项式系数微分算子环上的模}

\section{Hilbert 多项式}

设 $A=\bigoplus_{n\in\mathbb{Z}} A^n$ 是一个分次诺特交换环，带有单位元 $1$，且 $1$ 含于 $A^0$ 中。于是 $A^0$ 是带单位元 $1$ 的交换环。假设当 $n<0$ 时 $A^n=0$。

\begin{lemma}\ 
\begin{enumerate}[label=(\roman*)]
  \item $A^0$ 是诺特环；
  \item $A$ 是有限生成的 $A^0$-代数。
\end{enumerate}
\end{lemma}

\begin{proof}
(i) 令 $A_+=\bigoplus_{n=1}^{\infty}A^n$。则 $A_+$ 是 $A$ 的理想且 $A^0=A/A_+$。

(ii) $A_+$ 是有限生成的。设 $x_1,x_2,\dots,x_s$ 是 $A_+$ 的一组齐次生成元，记 $d_i=\deg x_i$，$1\le i\le s$。令 $B$ 为 $x_1,x_2,\dots,x_s$ 生成的 $A^0$-子代数。我们断言对所有 $n\in\mathbb{Z}_+$ 有 $A^n\subseteq B$。显然 $A^0\subseteq B$。设 $n>0$ 且 $y\in A^n$。则 $y\in A_+$，故 $y=\sum_{i=1}^s y_i x_i$，其中 $y_i\in A^{n-d_i}$。由归纳假设可用于 $y_i$，$1\le i\le s$。这说明 $y\in B$。
\end{proof}

1.1 的逆命题由 Hilbert 定理得出，该定理指出：若环 $A^0$ 是诺特的，则多项式环 $A^0[X_1,X_2,\dots,X_n]$ 是诺特的。

设 $M=\bigoplus_{n\in\mathbb{Z}} M^n$ 是一个有限生成的分次 $A$-模。则每个 $M^n$（$n\in\mathbb{Z}$）都是 $A^0$-模。并且当 $n$ 充分小时 $M^n=0$。

\begin{lemma}
$A^0$-模 $M^n$（$n\in\mathbb{Z}$）是有限生成的。
\end{lemma}

\begin{proof}
设 $m_i$（$1\le i\le k$）是 $M$ 的齐次生成元，$\deg m_i=r_i$。对 $j\in\mathbb{Z}_+$，记 $z_i(j)$，$1\le i\le \ell(j)$，为 $x_1,x_2,\dots,x_s$ 中次数为 $j$ 的所有齐次单项式。设 $m\in M^n$。则 $m=\sum_{i=1}^k y_i m_i$，其中 $y_i\in A^{n-r_i}$。由 1.1，$y_i=\sum_j a_{ij} z_j(n-r_i)$，其中 $a_{ij}\in A^0$。这意味着 $m=\sum_{i,j}a_{ij}z_j(n-r_i)m_i$，故 $M^n$ 由 $(z_j(n-r_i)m_i;\ 1\le j\le \ell(n-r_i),\ 1\le i\le k)$ 生成。
\end{proof}

令 $\mathcal{M}_{fg}(A^0)$ 为有限生成 $A^0$-模的范畴。设 $\lambda$ 是 $\mathcal{M}_{fg}(A^0)$ 上取值于 $\mathbb{Z}$ 的函数。若对任意短正合列
$$0\longrightarrow M'\longrightarrow M\longrightarrow M''\longrightarrow 0$$
都有
$$\lambda(M)=\lambda(M')+\lambda(M''),$$
则称 $\lambda$ 是\emph{可加的}。显然，可加性蕴含 $\lambda(0)=0$。

\begin{lemma}
设
$$0\to M_0\to M_1\to M_2\to\cdots\to M_n\to 0$$
是 $\mathcal{M}_{fg}(A^0)$ 中的正合列。则
$$\sum_{i=0}^n (-1)^i\lambda(M_i)=0.$$
\end{lemma}

\begin{proof}
显然。
\end{proof}

设 $\mathbb{Z}[[t]]$ 为系数在 $\mathbb{Z}$ 中的形式幂级数环。记 $\mathbb{Z}((t))$ 为 $\mathbb{Z}[[t]]$ 关于乘法系 $\{t^n\mid n\in\mathbb{Z}_+\}$ 的局部化。

设 $M$ 是有限生成分次 $A$-模。则 $M$ 关于 $\lambda$ 的 Poincaré 级数为
$$P(M,t)=\sum_{n\in\mathbb{Z}}\lambda(M^n)\,t^n\ \in \mathbb{Z}((t)).$$

例如，设 $A=k[X_1,X_2,\dots,X_s]$ 是 $s$ 个变量、系数在域 $k$ 中、按全次数分次的多项式代数。则 $A^0=k$，且对每个有限生成分次 $A$-模 $M$ 有 $\dim_k M^n<\infty$。因此可对 $\lambda=\dim_k$ 定义 Poincaré 级数。特别地，对 $A$ 本身有
$$P(A,t)=\sum_{n\in\mathbb{Z}}\dim_k A^n t^n=\sum_{n=0}^{\infty}\binom{s+n-1}{s-1}t^n=\frac{1}{(1-t)^s}.$$

下一个结果表明一般情形的 Poincaré 级数具有类似形式。

\begin{theorem}[Hilbert, Serre]
对任意有限生成分次 $A$-模 $M$ 有
$$P(M,t)=\frac{f(t)}{\prod_{i=1}^s (1-t^{d_i})},$$
其中 $f(t)\in\mathbb{Z}[t,t^{-1}]$。
\end{theorem}

\begin{proof}
对 $s$ 归纳证明。若 $s=0$，则 $A=A^0$，$M$ 是有限生成 $A^0$-模。这意味着当 $n$ 充分大时 $M^n=0$。因此除有限多个 $n\in\mathbb{Z}$ 外 $\lambda(M^n)=0$，且 $P(M,t)\in\mathbb{Z}[t,t^{-1}]$。

现设 $s>0$。乘以 $x_s$ 定义了 $M$ 的一个 $A$-模自同态 $f$。令 $K=\ker f$，$I=\operatorname{im} f$，$L=M/I$。则 $K,I,L$ 是分次 $A$-模，且有正合列
$$0\longrightarrow K\longrightarrow M\xrightarrow{\ f\ } M\longrightarrow L\longrightarrow 0.$$
这意味着
$$0\longrightarrow K^n\longrightarrow M^n\xrightarrow{\ x_s\ } M^{n+d_s}\longrightarrow L^{n+d_s}\longrightarrow 0$$
对所有 $n\in\mathbb{Z}$ 都是 $A^0$-模的正合列。特别地，由 1.3，
$$\lambda(K^n)-\lambda(M^n)+\lambda(M^{n+d_s})-\lambda(L^{n+d_s})=0$$
对所有 $n\in\mathbb{Z}$ 成立。这意味着
\begin{align*}
(1-t^{d_s})P(M,t)&=\sum_{n\in\mathbb{Z}}\lambda(M^n)t^n-\sum_{n\in\mathbb{Z}}\lambda(M^n)t^{n+d_s}\\
&=\sum_{n\in\mathbb{Z}}\bigl(\lambda(M^{n+d_s})-\lambda(M^n)\bigr)t^{n+d_s}\\
&=\sum_{n\in\mathbb{Z}}\bigl(\lambda(L^{n+d_s})-\lambda(K^n)\bigr)t^{n+d_s}\\
&=P(L,t)-P(K,t)\,t^{d_s},
\end{align*}
即
$$(1-t^{d_s})P(M,t)=P(L,t)-t^{d_s}P(K,t).$$
由构造可知 $x_s$ 在 $L$ 和 $K$ 上的作用都是乘以 $0$，即可以把它们看作 $A/(x_s)$-模。因此归纳假设适用，这就立即得到结论。
\end{proof}

由于 Poincaré 级数 $P(M,t)$ 是有理函数，可以讨论它在一点的极点阶数。记 $d_\lambda(M)$ 为 $P(M,t)$ 在 $1$ 处极点的阶数。

由定理，$f(t)=\sum_{k\in\mathbb{Z}}a_k t^k$，其中 $a_k\in\mathbb{Z}$ 且除有限多个 $k$ 外 $a_k=0$。设 $p$ 为 $f$ 在 $1$ 处零点的阶数。假设 $p>0$。则 $f(t)=(1-t)g(t)$，其中 $g(t)=\sum_{k\in\mathbb{Z}}b_k t^k$，$b_k\in\mathbb{Q}$ 且除有限多个 $k$ 外 $b_k=0$。此外对所有 $k\in\mathbb{Z}$ 有 $a_k=b_k-b_{k-1}$。对 $k$ 归纳可得 $b_k\in\mathbb{Z}$。必要时重复此过程，得到 $f(t)=(1-t)^p g(t)$，其中 $g(t)=\sum_{k\in\mathbb{Z}}b_k t^k$，$b_k\in\mathbb{Z}$ 且除有限多个 $k$ 外 $b_k=0$，并且 $g(1)\neq 0$。

\begin{corollary}
若对 $1\le i\le s$ 有 $d_i=1$，则函数 $n\longmapsto \lambda(M^n)$ 对充分大的 $n\in\mathbb{Z}$ 等于一个有理系数多项式，次数为 $d_\lambda(M)-1$。
\end{corollary}

\begin{proof}
设 $p$ 为 $f$ 在 $1$ 处零点的阶数。则可写 $f(t)=(1-t)^p g(t)$，其中 $g(1)\neq 0$。此外令 $d=d_\lambda(M)=s-p$，于是
$$P(M,t)=\frac{g(t)}{(1-t)^d}.$$
另一方面
$$(1-t)^{-d}=\sum_{k=0}^{\infty}\frac{d(d+1)\cdots(d+k-1)}{k!}t^k=\sum_{k=0}^{\infty}\binom{d+k-1}{d-1}t^k,$$
若写 $g(t)=\sum_{k=-N}^{N}a_k t^k$，则对所有 $n\ge N$，
$$\lambda(M^n)=\sum_{k=-N}^{N}a_k\binom{d+n-k-1}{d-1}.$$
这等于
$$\sum_{k=-N}^{N}a_k\frac{(d+n-k-1)!}{(d-1)!(n-k)!}
=\sum_{k=-N}^{N}a_k\frac{(n-k+1)(n-k+2)\cdots(n-k+d-1)}{(d-1)!},$$
故 $\lambda(M^n)$ 是 $n$ 的多项式，首项为
$$\Bigl(\sum_{k=-N}^{N}a_k\Bigr)\frac{n^{d-1}}{(d-1)!}=g(1)\frac{n^{d-1}}{(d-1)!}\neq 0.$$
\end{proof}

我们把对充分大的 $n\in\mathbb{Z}$ 给出 $\lambda(M^n)$ 的多项式称为 $M$ 关于 $\lambda$ 的 \emph{Hilbert 多项式}。由证明可见 $M$ 的 Hilbert 多项式的首项系数为 $\dfrac{g(1)}{(d-1)!}$。

回到 $A=k[X_1,X_2,\dots,X_s]$ 的例子，我们看到
$$\dim_k A^n=\binom{n+s-1}{s-1}=\frac{n^{s-1}}{(s-1)!}+\cdots,$$
故对 $A=k[X_1,X_2,\dots,X_s]$，Hilbert 多项式的次数等于 $s-1$。

现在证明关于（系数在特征 $0$ 的域中、对充分大正整数取整数值的）多项式的一个刻画。首先注意，对任意 $s\in\mathbb{Z}_+$ 和 $q\ge s$ 有
$$q^s=s!\binom{q}{s}+Q(q),$$
其中 $Q$ 是次数为 $s-1$ 的多项式。因此任意 $d$ 次多项式 $P$ 对充分大的 $q$ 可唯一写成
$$P(q)=c_0\binom{q}{d}+c_1\binom{q}{d-1}+\cdots+c_{d-1}\binom{q}{1}+c_d,$$
其中系数 $c_i$（$0\le i\le d$）适当选取。由于二项式系数是整数，若所有 $c_i$（$0\le i\le d$）是整数，则对整数 $n\ge d$，$P$ 取整数值。下一个结果是这一观察的逆命题。

\begin{lemma}
若多项式
$$q\longmapsto P(q)=c_0\binom{q}{d}+c_1\binom{q}{d-1}+\cdots+c_{d-1}\binom{q}{1}+c_d$$
对充分大的 $n\in\mathbb{Z}$ 取整数值 $P(n)$，则所有系数 $c_i$（$0\le i\le d$）都是整数。
\end{lemma}

\begin{proof}
对 $d$ 归纳证明。若 $d=0$ 结论显然。此外
\begin{align*}
P(q+1)-P(q)&=\sum_{i=0}^d c_i\binom{q+1}{d-i}-\sum_{i=0}^d c_i\binom{q}{d-i}\\
&=\sum_{i=0}^d c_i\left(\binom{q+1}{d-i}-\binom{q}{d-i}\right)=\sum_{i=0}^{d-1}c_i\binom{q}{d-i-1},
\end{align*}
这里用到恒等式（对 $q\ge s\ge 1$）
$$\binom{q+1}{s}=\binom{q}{s}+\binom{q}{s-1}.$$
因此 $q\longmapsto P(q+1)-P(q)$ 是系数为 $c_0,c_1,\dots,c_{d-1}$ 的多项式，且对充分大的 $n\in\mathbb{Z}$ 有 $P(n)\in\mathbb{Z}$。由归纳假设，所有 $c_i$（$0\le i\le d-1$）是整数。这立即蕴含 $c_d$ 也是整数。
\end{proof}

还需要另一个相关的注记。若 $F$ 是首项系数为 $a_0$ 的 $d$ 次多项式，则
\begin{align*}
G(n)&=F(n)-F(n-1)\\
&=(a_0 n^d+a_1 n^{d-1}+\cdots)-(a_0(n-1)^d+a_1(n-1)^{d-1}+\cdots)\\
&=a_0 d\, n^{d-1}+\cdots
\end{align*}
是 $n$ 的 $d-1$ 次多项式，首项系数为 $d a_0$。下一个结果是这一事实的逆命题。

\begin{lemma}
设 $F$ 是 $\mathbb{Z}$ 上的函数，使得
$$G(n)=F(n)-F(n-1)$$
对充分大的 $n\in\mathbb{Z}$ 等于 $n$ 的 $d-1$ 次多项式。则 $F$ 对充分大的 $n\in\mathbb{Z}$ 等于 $n$ 的 $d$ 次多项式。
\end{lemma}

\begin{proof}
假设对 $n\ge N\ge d$ 有 $G(n)=P(n-1)$，其中 $P$ 是 $n$ 的 $d-1$ 次多项式。则由 1.6 有
$$P(n)=\sum_{i=0}^{d-1}c_i\binom{n}{d-i-1}.$$
于是对 $n\ge N+1$，
$$F(n)=\sum_{k=N+1}^{n}\bigl(F(k)-F(k-1)\bigr)+F(N)=\sum_{k=N+1}^{n}G(k)+F(N)=\sum_{k=d}^{n}P(k-1)+C,$$
其中 $C$ 是常数。此外，由上一定理中用到的恒等式，对 $q>s\ge 1$，
$$\binom{q}{s}=\sum_{j=s+1}^{q}\left(\binom{j}{s}-\binom{j-1}{s}\right)+1=\sum_{j=s+1}^{q}\binom{j-1}{s-1}+1=\sum_{j=s}^{q}\binom{j-1}{s-1}.$$
这意味着
\begin{align*}
\sum_{k=d}^{n}P(k-1)&=\sum_{k=d}^{n}\sum_{i=0}^{d-1}c_i\binom{k-1}{d-i-1}=\sum_{i=0}^{d-1}c_i\sum_{k=d}^{n}\binom{k-1}{d-i-1}\\
&=\sum_{i=0}^{d-1}c_i\sum_{k=d-i}^{n}\binom{k-1}{d-i-1}-\sum_{i=1}^{d-1}c_i\sum_{k=d-i}^{d-i-1}\binom{k-1}{d-i-1}\\
&=\sum_{i=0}^{d-1}c_i\binom{n}{d-i}+C_0
\end{align*}
对某个常数 $C_0$ 成立。
\end{proof}

特别地，由此可知 $\sum_{n\le N}\lambda(M^n)$ 对充分大的 $N\in\mathbb{Z}$ 等于次数为 $d_\lambda(M)$ 的多项式。此外，若对充分大的 $N\in\mathbb{Z}$ 写
$$\sum_{n\le N}\lambda(M^n)=a_0 N^d+a_1 N^{d-1}+\cdots+a_{d-1}N+a_d,$$
则 $d!\,a_0$ 是整数。

例如，若 $A=k[X_1,X_2,\dots,X_s]$，则次数 $\le N$ 的所有多项式空间的维数等于
$$\sum_{n=0}^{N}\dim_k(A^n)=\sum_{n=0}^{N}\binom{s+n-1}{s-1}=\binom{s+N}{s}=\frac{N^s}{s!}+\cdots.$$
\begin{annote}
讲义用 Hilbert 多项式 $P(M,t)$ 的次数与首项系数来定义 $\dim M$ 与重数 $e(M)$。这与交换代数里按素理想链定义的 Krull 维数一致：若 $R$ 是诺特局部环、$M$ 有限生成，则 $\dim M=\dim R/\mathrm{Ann}(M)$。文中反复使用的“$\dim M=\deg P(M,t)$”正是 Hilbert–Serre 定理的推论，它把“模的大小”换算成一条多项式的次数，后文特征簇的维数比较全靠这一点。[来源：Stacks Project Tag 00KE（Definition 10.60.2, Krull dimension）；原书 §1, p1–6]
\end{annote}



\section{局部环上模的维数}

\begin{lemma}[Nakayama]
设 $A$ 是以 $\mathfrak{m}$ 为极大理想的局部环。设 $V$ 是有限生成 $A$-模且满足 $\mathfrak{m}V=V$。则 $V=0$。
\end{lemma}

\begin{proof}
假设 $V\neq 0$。则可取 $V$ 作为 $A$-模的一组最小生成元 $v_1,\dots,v_s$。由假设，$v_s=\sum_{i=1}^s m_i v_i$，其中 $m_i\in\mathfrak{m}$。因此 $(1-m_s)v_s=\sum_{i=1}^{s-1}m_i v_i$。由于 $1-m_s$ 可逆，这说明 $v_1,\dots,v_{s-1}$ 生成 $V$，与假设矛盾。
\end{proof}

以下假设 $A$ 是诺特局部环，$\mathfrak{m}$ 为其极大理想，$k=A/\mathfrak{m}$ 为 $A$ 的剩余域。

\begin{lemma}
$\dim_k(\mathfrak{m}/\mathfrak{m}^2)<+\infty$。
\end{lemma}

\begin{proof}
由诺特性假设 $\mathfrak{m}$ 有限生成。若 $a_1,\dots,a_p$ 是 $\mathfrak{m}$ 的生成元，则它们在 $\mathfrak{m}/\mathfrak{m}^2$ 中的像 $\bar a_1,\dots,\bar a_p$ 张成它作为 $k$-向量空间。
\end{proof}

令 $s=\dim_k(\mathfrak{m}/\mathfrak{m}^2)$。则可找到 $a_1,\dots,a_s\in\mathfrak{m}$，使得 $\bar a_1,\dots,\bar a_s$ 构成 $\mathfrak{m}/\mathfrak{m}^2$ 的一组基。我们断言它们生成 $\mathfrak{m}$。设 $I$ 为由 $a_1,\dots,a_s$ 生成的理想。则 $I+\mathfrak{m}^2=\mathfrak{m}$ 且 $\mathfrak{m}(\mathfrak{m}/I)=\mathfrak{m}/I$。故由 2.1 得 $\mathfrak{m}/I=0$。因此我们证明了：

\begin{lemma}
正整数 $\dim_k(\mathfrak{m}/\mathfrak{m}^2)$ 等于 $\mathfrak{m}$ 的最小生成元个数。
\end{lemma}

$\mathfrak{m}$ 中满足 $(\bar a_1,\dots,\bar a_s)$ 构成 $\mathfrak{m}/\mathfrak{m}^2$ 的一组基的任意 $s$ 元组 $(a_1,\dots,a_s)$ 称为 $A$ 中的一个\emph{坐标系}。

显然，$(\mathfrak{m}^p;\ p\in\mathbb{Z}_+)$ 是 $A$ 的一个递减滤过。因此可以构成
$$\Gr A=\bigoplus_{p=0}^{\infty}\mathfrak{m}^p/\mathfrak{m}^{p+1}.$$
我们断言 $\Gr A$ 是 $k$ 上的有限生成代数，因而是诺特分次环。事实上，映射 $X_i\longmapsto \bar a_i\in\mathfrak{m}/\mathfrak{m}^2\subset\Gr A$ 扩充为 $k[X_1,\dots,X_s]$ 到 $\Gr A$ 上的满射。

设 $M$ 是有限生成 $A$-模。则可用 $\mathfrak{m}^p M$（$p\in\mathbb{Z}_+$）定义 $M$ 的递减滤过，并考虑分次 $\Gr A$-模
$$\Gr M=\bigoplus_{p=0}^{\infty}\mathfrak{m}^p M/\mathfrak{m}^{p+1}M.$$

\begin{lemma}
若 $M$ 是有限生成 $A$-模，则 $\Gr M$ 是有限生成 $\Gr A$-模。
\end{lemma}

\begin{proof}
由 $\Gr M$ 的定义可见对所有 $p\in\mathbb{Z}_+$ 有 $\mathfrak{m}\cdot \Gr_p M=\Gr_{p+1}M$，故 $\Gr_0 M=M/\mathfrak{m}M$ 生成 $\Gr M$。另一方面，$M/\mathfrak{m}M$ 是 $k$ 上的有限维线性空间。
\end{proof}

这意味着（由 1.2）$\dim_k(\mathfrak{m}^p M/\mathfrak{m}^{p+1}M)<+\infty$，特别地 $A$-模 $\mathfrak{m}^p M/\mathfrak{m}^{p+1}M$ 具有有限长度。由于长度显然是可加函数，由 1.5 可见 $p\longmapsto \operatorname{length}_A(\mathfrak{m}^p M/\mathfrak{m}^{p+1}M)=\dim_k(\mathfrak{m}^p M/\mathfrak{m}^{p+1}M)$ 对充分大的 $p\in\mathbb{Z}_+$ 等于有理系数多项式。此外，函数
$$p\longmapsto \operatorname{length}_A(M/\mathfrak{m}^p M)=\sum_{q=0}^{p-1}\operatorname{length}_A(\mathfrak{m}^q M/\mathfrak{m}^{q+1}M)$$
对充分大的 $p\in\mathbb{Z}_+$ 等于有理系数多项式，且其首项形如 $e\dfrac{p^d}{d!}$，其中 $e,d\in\mathbb{Z}_+$。我们记 $d(M)=d$，$e(M)=e$，并称这些数为 $M$ 的\emph{维数}和\emph{重数}。

现在讨论函数 $M\longmapsto d(M)$ 的一些性质。控制 $A$-模滤过的关键结果是 Artin--Rees 引理。

\begin{theorem}[Artin, Rees]
设 $M$ 是有限生成 $A$-模，$N$ 是其子模。则存在 $m_0\in\mathbb{Z}_+$ 使得对所有 $p\in\mathbb{Z}_+$，
$$\mathfrak{m}^{p+m_0}M\cap N=\mathfrak{m}^p(\mathfrak{m}^{m_0}M\cap N).$$
\end{theorem}

\begin{proof}
令 $A^*=\bigoplus_{n=0}^{\infty}\mathfrak{m}^n$。则 $A^*$ 具有天然的分次环结构。设 $(a_1,\dots,a_s)$ 是 $A$ 中的坐标系，则有自然的满射
$$A[a_1,\dots,a_s]\longrightarrow A^*,$$
且 $A^*$ 是分次诺特环。令 $M^*=\bigoplus_{n=0}^{\infty}\mathfrak{m}^n M$。则 $M^*$ 是分次 $A^*$-模。它显然作为 $A^*$-模由 $M^*_0=M$ 生成。由于 $M$ 是有限生成 $A$-模，我们断定 $M^*$ 是有限生成 $A^*$-模。

此外，令 $N^*=\bigoplus_{n=0}^{\infty}(N\cap \mathfrak{m}^n M)\subset M^*$。则由
$$\mathfrak{m}^p(N\cap \mathfrak{m}^n M)\subset \mathfrak{m}^p N\cap \mathfrak{m}^{n+p}M\subset N\cap \mathfrak{m}^{n+p}M$$
可知 $N^*$ 是 $M^*$ 的 $A^*$-子模。由于 $A^*$ 是诺特环，$N^*$ 有限生成。存在 $m_0\in\mathbb{Z}_+$ 使得 $\bigoplus_{n=0}^{m_0}(N\cap\mathfrak{m}^n M)$ 生成 $N^*$。于是对任意 $p\in\mathbb{Z}_+$，
$$N\cap\mathfrak{m}^{p+m_0}M=\sum_{s=0}^{m_0}\mathfrak{m}^{p+m_0-s}(N\cap\mathfrak{m}^s M)\subset \mathfrak{m}^p(N\cap\mathfrak{m}^{m_0}M)\subset N\cap\mathfrak{m}^{p+m_0}M.$$
\end{proof}

这一结果有如下推论——Krull 交定理。

\begin{theorem}[Krull]
设 $M$ 是有限生成 $A$-模。则
$$\bigcap_{p=0}^{\infty}\mathfrak{m}^p M=\{0\}.$$
\end{theorem}

\begin{proof}
令 $E=\bigcap_{p=0}^{\infty}\mathfrak{m}^p M$。则由 2.5，
$$E=\mathfrak{m}^{p+m_0}M\cap E=\mathfrak{m}^p(\mathfrak{m}^{m_0}M\cap E)=\mathfrak{m}^p E,$$
特别地 $\mathfrak{m}E=E$，由 Nakayama 引理得 $E=0$。
\end{proof}

\begin{lemma}
设
$$0\longrightarrow M'\longrightarrow M\longrightarrow M''\longrightarrow 0$$
是有限生成 $A$-模的正合列。则
\begin{enumerate}[label=(\roman*)]
  \item $d(M)=\max(d(M'),d(M''))$；
  \item 若 $d(M)=d(M')=d(M'')$，则 $e(M)=e(M')+e(M'')$。
\end{enumerate}
\end{lemma}

\begin{proof}
可把 $M'$ 看作 $M$ 的子模。若给 $M$ 配置滤过 $\mathfrak{m}^p M$（$p\in\mathbb{Z}_+$），给 $M'$ 和 $M''$ 配置诱导滤过 $M'\cap\mathfrak{m}^p M$（$p\in\mathbb{Z}_+$）和 $\mathfrak{m}^p M''$（$p\in\mathbb{Z}_+$），则得到正合列
$$0\longrightarrow \Gr M'\longrightarrow \Gr M\longrightarrow \Gr M''\longrightarrow 0.$$
这意味着对任意 $p\in\mathbb{Z}_+$，
$$\operatorname{length}_A(\mathfrak{m}^p M/\mathfrak{m}^{p+1}M)
=\operatorname{length}_A\bigl((M'\cap\mathfrak{m}^p M)/(M'\cap\mathfrak{m}^{p+1}M)\bigr)+\operatorname{length}_A(\mathfrak{m}^p M''/\mathfrak{m}^{p+1}M'')$$
并且求和得
$$\operatorname{length}_A(M/\mathfrak{m}^p M)=\operatorname{length}_A(M'/(M'\cap\mathfrak{m}^p M))+\operatorname{length}_A(M''/\mathfrak{m}^p M'').$$
因此函数 $p\longmapsto \operatorname{length}_A(M'/(M'\cap\mathfrak{m}^p M))$ 对充分大的 $p\in\mathbb{Z}_+$ 等于 $p$ 的多项式。另一方面，由 2.5，
$$\mathfrak{m}^{p+m_0}M'\subset \mathfrak{m}^{p+m_0}M\cap M'\subset \mathfrak{m}^p M';$$
故对充分大的 $p\in\mathbb{Z}_+$，函数 $p\longmapsto \operatorname{length}_A(M'/(M'\cap\mathfrak{m}^p M))$ 与 $p\longmapsto \operatorname{length}_A(M'/\mathfrak{m}^p M')$ 由首项相同的多项式给出。
\end{proof}

\begin{corollary}
设 $A$ 是诺特局部环，$s=\dim_k(\mathfrak{m}/\mathfrak{m}^2)$。则对任意有限生成 $A$-模 $M$ 有 $d(M)\le s$。
\end{corollary}

\begin{proof}
由 2.7，只需证明 $d(A)\le s$。这由 $k[X_1,\dots,X_s]$ 到 $\Gr A$ 的满同态的存在性，以及 $s$ 个变量中次数 $\le n$ 的多项式空间的维数是 $n$ 的 $s$ 次多项式这一事实立即得出。
\end{proof}

若诺特局部环满足 $d(A)=\dim_k(\mathfrak{m}/\mathfrak{m}^2)$，则称为\emph{正则的}。

\begin{theorem}
设 $A$ 是诺特局部环，$(a_1,a_2,\dots,a_s)$ 是 $A$ 中的坐标系。则下列条件等价：
\begin{enumerate}[label=(\roman*)]
  \item $A$ 是正则局部环；
  \item 由 $X_i\longmapsto \bar a_i$（$1\le i\le s$）定义的 $k[X_1,X_2,\dots,X_s]$ 到 $\Gr A$ 的典范态射是同构。
\end{enumerate}
\end{theorem}

\begin{proof}
由定义，$k[X_1,\dots,X_s]$ 到 $\Gr A$ 的典范态射是满射。设 $I$ 是 $k[X_1,\dots,X_s]$ 到 $\Gr A$ 的自然满射的核，它是分次理想。若 $I\neq 0$，它含有次数 $d>0$ 的齐次多项式 $P$。设 $J$ 是 $k[X_1,X_2,\dots,X_s]$ 中由 $P$ 生成的理想。则其 Poincaré 级数为 $P(J,t)=\dfrac{t^d}{(1-t)^s}$。显然，
$$P(k[X_1,\dots,X_s]/J,t)=P(k[X_1,\dots,X_s],t)-P(J,t)=\frac{1-t^d}{(1-t)^s}=\frac{1+t+\cdots+t^{d-1}}{(1-t)^{s-1}}.$$
Poincaré 级数 $P(k[X_1,\dots,X_s]/J,t)$ 在 $1$ 处的极点阶数为 $s-1$，由 1.5 函数 $\dim_k(k[X_1,\dots,X_s]/J)_n$ 对充分大的 $n\in\mathbb{Z}_+$ 由 $n$ 的 $s-2$ 次多项式给出。由此可知 $\dim_k(k[X_1,\dots,X_s]/I)_n=\dim_k\Gr_n A$ 对充分大的 $n\in\mathbb{Z}_+$ 由次数 $\le s-2$ 的多项式给出。这意味着 $d(A)\le s-1$。因此 $I=0$ 当且仅当 $d(A)=s$。
\end{proof}

\begin{theorem}
设 $A$ 是正则局部环。则 $A$ 是整环。
\end{theorem}

\begin{proof}
设 $a,b\in A$ 且 $a\neq 0$，$b\neq 0$。则由 2.6，可找到 $p,q\in\mathbb{Z}_+$ 使得 $a\in\mathfrak{m}^p$，$a\notin\mathfrak{m}^{p+1}$，且 $b\in\mathfrak{m}^q$，$b\notin\mathfrak{m}^{q+1}$。则它们在 $\Gr_p A$ 和 $\Gr_q A$ 中的像 $\bar a$ 与 $\bar b$ 均非零，且由于 $\Gr A$ 由 2.9 是整环，得 $\bar a\bar b\neq 0$。因此 $ab\neq 0$。
\end{proof}

最后讨论一个以后将起重要作用例子。设 $k$ 是域，$A=k[X_1,X_2,\dots,X_n]$ 是系数在 $k$ 中的 $n$ 元多项式环，$\widehat A=k[[X_1,X_2,\dots,X_n]]$ 是系数在 $k$ 中的 $n$ 元形式幂级数环。容易验证 $\widehat A$ 是以 $X_1,X_2,\dots,X_n$ 生成的极大理想 $\widehat{\mathfrak{m}}$ 的局部环。此外，$k[X_1,X_2,\dots,X_n]$ 到 $\Gr\widehat A$ 的典范态射显然是同构。

对任意 $x\in k^n$，记 $\mathfrak{m}_x$ 为 $A$ 中由 $X_i-x_i$（$1\le i\le n$）生成的极大理想。则它在 $A$ 中的补集是 $A$ 中的乘法系，记 $A_x$ 为 $A$ 相应的局部化。它同构于 $k^n$ 上在 $x$ 处正则的所有有理函数环。这显然是诺特局部环。$\mathfrak{m}_x$ 的局部化是极大理想 $\mathfrak{n}_x=(\mathfrak{m}_x)_x$，即所有在 $x$ 处消失的有理函数。由 $X_i\longmapsto X_i-x_i$（$1\le i\le n$）定义的 $A$ 的自同构给出对任意 $x\in k^n$，$A^0$ 与 $A_x$ 的一个同构。另一方面，$A$ 到 $\widehat A$ 的自然同态扩充为 $A^0$ 到 $\widehat A$ 的单射同态。此同态保持这些局部环上的滤过，并诱导 $\Gr A^0$ 到 $\Gr\widehat A$ 的典范同构。因此有如下结果。

\begin{proposition}
环 $A_x$（$x\in k^n$）是 $n$ 维正则局部环。
\end{proposition}
\begin{annote}
讲义 §2 证的正是 Artin–Rees 引理（Lemma 10.51.2）：$R$ 诺特、$I\subset R$ 理想、$N\subset M$ 有限生成，则存在 $c>0$ 使 $I^nM\cap N=I^{n-c}(I^cM\cap N)$ 对 $n\ge c$ 成立。文中引入的 $A^*=\bigoplus_{n\ge 0}\mathfrak m^n$ 就是标准证明里的 Rees 环 $S=R\oplus I\oplus I^2\oplus\cdots$，同一构造还给出 Krull 交定理 $\bigcap_{n\ge0}I^nM=0$（$R$ 诺特局部、$I$ 真理想）。这两条是 §3、§4 滤过论证的唯一外部输入。[来源：Stacks Project Tag 00IN（Lemma 10.51.2, Artin–Rees）、Tag 00IP（Lemma 10.51.4, Krull 交定理）]
\end{annote}



\section{滤过环上模的维数}

设 $D$ 是带单位的环，$(D_n;\ n\in\mathbb{Z})$ 是 $D$ 关于加法子群的递增滤过，满足
\begin{enumerate}[label=(\roman*)]
  \item $D_n=\{0\}$ 当 $n<0$；
  \item $\bigcup_{n\in\mathbb{Z}}D_n=D$；
  \item $1\in D_0$；
  \item 对任意 $n,m\in\mathbb{Z}$，$D_n\cdot D_m\subset D_{n+m}$；
  \item 对任意 $n,m\in\mathbb{Z}$，$[D_n,D_m]\subset D_{n+m-1}$。
\end{enumerate}
则 $\Gr D=\bigoplus_{n\in\mathbb{Z}}\Gr_n D=\bigoplus_{n\in\mathbb{Z}}D_n/D_{n-1}$ 是带单位的分次环。性质 (v) 蕴含它是交换的。特别地 $D_0=\Gr_0 D$ 是带单位的交换环。因此可以把 $\Gr D$ 看作 $D_0$-代数。进一步假设 $D$ 满足
\begin{enumerate}[label=(\roman*),resume]
  \item $\Gr D$ 是诺特环；
  \item $\Gr_1 D$ 作为 $D_0$-代数生成 $\Gr D$。
\end{enumerate}
则由 1.1，$D_0$ 是诺特环。此外，由 (vi)、(vii) 和 1.2 可知可取有限多个元素 $x_1,x_2,\dots,x_s\in\Gr_1 D$，使得 $\Gr D$ 作为 $D_0$-代数由它们生成。显然由 (vii)，对 $n\in\mathbb{Z}_+$ 有
$$\Gr_{n+1}D=\Gr_1 D\cdot\Gr_n D,$$
因此
$$D_{n+1}=D_n\cdot D_1\quad\text{对 }n\in\mathbb{Z}_+.$$

设 $D^\circ$ 是 $D$ 的反环。则滤过 $(D_n;\ n\in\mathbb{Z})$ 关于 $D^\circ$ 的乘法具有相同性质。此外，恒等映射 $D\longrightarrow D^\circ$ 诱导 $\Gr D$ 与 $\Gr D^\circ$ 的分次环同构。

设 $M$ 是 $D$-模。$M$ 关于加法子群的递增滤过 $F M=(F_n M;\ n\in\mathbb{Z})$ 称为 $D$-模滤过，若对 $n,m\in\mathbb{Z}$ 有 $D_n\cdot F_m M\subset F_{m+n}M$。特别地 $F_n M$ 是 $D_0$-模。

$D$-模滤过 $F M$ 称为\emph{Hausdorff 的}，若 $\bigcap_{n\in\mathbb{Z}}F_n M=\{0\}$。它是\emph{穷尽的}，若 $\bigcup_{n\in\mathbb{Z}}F_n M=M$。若存在 $m_0\in\mathbb{Z}$ 使得对所有 $n\in\mathbb{Z}_+$ 和 $m\ge m_0$ 有 $D_n\cdot F_m M=F_{m+n}M$，则称它是\emph{稳定的}。

$D$-模滤过称为\emph{好的}，若
\begin{enumerate}[label=(\roman*)]
  \item 对充分负的 $n\in\mathbb{Z}$，$F_n M=\{0\}$；
  \item 滤过 $F M$ 是穷尽的；
  \item $F_n M$（$n\in\mathbb{Z}$）是有限生成 $D_0$-模；
  \item 滤过 $F M$ 是稳定的。
\end{enumerate}
特别地，好滤过是 Hausdorff 的。

\begin{lemma}
设 $F M$ 是 $M$ 的穷尽 Hausdorff $D$-模滤过。则下列陈述等价：
\begin{enumerate}[label=(\roman*)]
  \item $F M$ 是好滤过；
  \item $\Gr D$-模 $\Gr M$ 是有限生成的。
\end{enumerate}
\end{lemma}

\begin{proof}
先证 (i)$\Rightarrow$(ii)。存在 $m_0\in\mathbb{Z}$ 使得对所有 $n\in\mathbb{Z}_+$ 有 $D_n\cdot F_{m_0}M=F_{n+m_0}M$。因此对 $n\in\mathbb{Z}_+$ 有 $\Gr_n D\cdot\Gr_{m_0}M=\Gr_{n+m_0}M$。由此可知 $\bigoplus_{n\le m_0}\Gr_n M$ 作为 $\Gr D$-模生成 $\Gr M$。由于 $F_n M$ 是有限生成 $D_0$-模，$\Gr_n M$ 也是有限生成 $D_0$-模。这说明（由于对充分负的 $n$，$F_n M=\{0\}$）$\bigoplus_{n\le m_0}\Gr_n M$ 是有限生成 $D_0$-模。

(ii)$\Rightarrow$(i)。显然对充分负的 $n\in\mathbb{Z}$ 有 $\Gr_n M=\{0\}$。此外由 1.2，所有 $\Gr_n M$ 是有限生成 $D_0$-模。正合列
$$0\longrightarrow F_{n-1}M\longrightarrow F_n M\longrightarrow \Gr_n M\longrightarrow 0$$
蕴含 $\bigcap_{n\in\mathbb{Z}}F_n M=F_{n_0}M$（对充分负的 $n$ 有 $F_n M=F_{n-1}M$，故存在 $n_0\in\mathbb{Z}$）。由于滤过 $F M$ 是 Hausdorff 的，$F_{n_0}M=\{0\}$。这由 $n$ 上的归纳蕴含所有 $F_n M$ 是有限生成 $D_0$-模。

设 $m_0\in\mathbb{Z}$ 使得 $\bigoplus_{n\le m_0}\Gr_n M$ 作为 $\Gr D$-模生成 $\Gr M$。设 $m\ge m_0$。则
\begin{align*}
\Gr_{m+1}M&=\bigoplus_{k\le m_0}\Gr_{m+1-k}D\cdot\Gr_k M
=\bigoplus_{k\le m_0}\Gr_1 D\cdot\Gr_{m-k}D\cdot\Gr_k M\\
&\subset\Gr_1 D\cdot\Gr_m M\subset\Gr_{m+1}M,
\end{align*}
即 $\Gr_1 D\cdot\Gr_m M=\Gr_{m+1}M$。这意味着
$$F_{m+1}M=D_1\cdot F_m M+F_m M=D_1\cdot F_m M,$$
并由 $n$ 上的归纳，
$$F_{m+n}M=D_1\cdot D_1\cdots D_1\cdot F_m M\subset D_n\cdot F_m M\subset F_{m+n}M.$$
因此对所有 $n\in\mathbb{Z}_+$ 有 $F_{m+n}M=D_n\cdot F_m M$。故 $F M$ 是好滤过。
\end{proof}

特别地，$(D_n;\ n\in\mathbb{Z})$ 是把 $D$ 看作左乘 $D$-模时的好滤过。

\begin{remark}
由证明可知，好滤过定义中的稳定性条件可替换为表面上较弱的条件：

(iv)$'$ 存在 $m_0\in\mathbb{Z}$ 使得对所有 $n\in\mathbb{Z}_+$ 有 $D_n\cdot F_{m_0}M=F_{m_0+n}M$。
\end{remark}

\begin{lemma}
设 $M$ 是带好滤过 $F M$ 的 $D$-模。则 $M$ 是有限生成的。
\end{lemma}

\begin{proof}
由定义 $\bigcup_{n\in\mathbb{Z}}F_n M=M$，且对某充分大的 $m_0\in\mathbb{Z}$ 和 $n\in\mathbb{Z}_+$ 有 $F_{n+m_0}M=D_n\cdot F_{m_0}M$。因此 $F_{m_0}M$ 作为 $D$-模生成 $M$。由于 $F_{m_0}M$ 是有限生成 $D_0$-模，结论成立。
\end{proof}

\begin{lemma}
设 $M$ 是有限生成 $D$-模。则 $M$ 有好滤过。
\end{lemma}

\begin{proof}
设 $U$ 是作为 $D$-模生成 $M$ 的有限生成 $D_0$-模。令 $n<0$ 时 $F_n M=0$，$n\ge 0$ 时 $F_n M=D_n\cdot U$。则 $U=\Gr_0 M$，且
$$\Gr_n M=F_n M/F_{n-1}M=(D_n\cdot U)/(D_{n-1}\cdot U)\subset\Gr_n D\cdot\Gr_0 M\subset\Gr_n M,$$
即对所有 $n\in\mathbb{Z}_+$ 有 $\Gr_n M=\Gr_n D\cdot\Gr_0 M$。故 $\Gr M$ 作为 $\Gr D$-模有限生成。结论由 3.1 得出。
\end{proof}

引理 3.1 和 3.3 蕴含：具有好滤过的 $D$-模恰好是有限生成 $D$-模。

\begin{proposition}
环 $D$ 是左诺特且右诺特的。
\end{proposition}

\begin{proof}
设 $L$ 是 $D$ 的左理想。$D$ 的自然滤过在 $L$ 上诱导滤过 $(L_n=L\cap D_n;\ n\in\mathbb{Z})$。这显然是 $D$-模滤过。分次模 $\Gr L$ 自然是 $\Gr D$ 的理想，且由于 $\Gr D$ 是诺特环，它作为 $\Gr D$-模有限生成。因此由 3.1 滤过 $(L_n;\ n\in\mathbb{Z})$ 是好滤过，由 3.3 $L$ 是有限生成的。这证明 $D$ 是左诺特的。

要得到右诺特性质，只需把 $D$ 换成其反环 $D^\circ$。
\end{proof}

若 $D$-模 $M$ 有两个滤过 $F M$ 和 $F'M$，若存在 $k\in\mathbb{Z}_+$ 使得对所有 $n\in\mathbb{Z}$ 有 $F_n M\subset F'_{n+k}M$，则称 $F M$ 比 $F'M$\emph{更细}。若 $F M$ 比 $F'M$ 更细且 $F'M$ 比 $F M$ 更细，则称它们\emph{等价}。

\begin{lemma}
设 $F M$ 是有限生成 $D$-模 $M$ 上的好滤过。则 $F M$ 比 $M$ 上任何其他穷尽 $D$-模滤过更细。
\end{lemma}

\begin{proof}
固定 $m_0\in\mathbb{Z}_+$ 使得对所有 $n\in\mathbb{Z}_+$ 有 $D_n\cdot F_{m_0}M=F_{n+m_0}M$。设 $F'M$ 是 $M$ 上另一个穷尽 $D$-模滤过。则 $F_{m_0}M$ 作为 $D_0$-模有限生成。由于 $F'M$ 是穷尽的，存在 $p\in\mathbb{Z}$ 使得 $F_{m_0}M\subset F'_p M$。由于 $F M$ 是好滤过，存在 $n_0$ 使得 $F_{n_0}M=\{0\}$。令 $k=p+|n_0|$。显然对 $m\le n_0$ 有 $F_m M=0\subset F'_{m+k}M$。对 $n_0<m\le m_0$，有 $-|n_0|\le n_0<m$ 且 $p=-|n_0|+k<m+k$。这给出
$$F_m M\subset F_{m_0}M\subset F'_p M\subset F'_{m+k}M.$$
最后对 $m>m_0$，由于 $m_0$ 为正有 $m-m_0\le m$，且 $p\le k$。于是
$$F_m M=D_{m-m_0}\cdot F_{m_0}M\subset D_m\cdot F'_p M\subset F'_{m+p}M\subset F'_{m+k}M.$$
\end{proof}

\begin{corollary}
有限生成 $D$-模上任意两个好滤过等价。
\end{corollary}

设 $M$ 是有限生成 $D$-模，$F M$ 是其上好滤过。则 $\Gr M$ 是有限生成 $\Gr D$-模，故可应用 §1 关于 Hilbert 多项式的结果。设 $\lambda$ 是有限生成 $D_0$-模上的可加函数。再假设 $\lambda$ 在有限生成 $D_0$-模范畴 $\mathcal{M}_{fg}(D_0)$ 的对象上仅取非负值。则由 1.5，
$$\lambda(F_n M)-\lambda(F_{n-1}M)=\lambda(\Gr_n M)$$
对充分大的 $n\in\mathbb{Z}_+$ 等于 $n$ 的多项式。由 1.7 这意味着 $\lambda(F_n M)$ 对充分大的 $n\in\mathbb{Z}_+$ 等于 $n$ 的多项式。若 $F'M$ 是 $M$ 上另一个好滤过，由 3.7 知 $F M$ 与 $F'M$ 等价，即存在 $k\in\mathbb{Z}_+$ 使得
$$F_n M\subset F'_{n+k}M\subset F_{n+2k}M$$
对所有 $n\in\mathbb{Z}$ 成立。由于 $\lambda$ 可加且仅取非负值，我们断定
$$\lambda(F_n M)\le\lambda(F'_{n+k}M)\le\lambda(F_{n+2k}M)$$
对所有 $n\in\mathbb{Z}$ 成立。这意味着对充分大的 $n$，表示 $\lambda(F_n M)$ 与 $\lambda(F'_n M)$ 的多项式有相等的首项。我们记这些多项式的公共次数为 $d_\lambda(M)$，并称之为 $D$-模 $M$ 关于 $\lambda$ 的\emph{维数}。由 1.6，这些多项式的首项系数形如 $e_\lambda(M)/d_\lambda(M)!$，其中 $e_\lambda(M)\in\mathbb{N}$。我们称 $e_\lambda(M)$ 为 $D$-模 $M$ 关于 $\lambda$ 的\emph{重数}。

设
$$0\longrightarrow M'\xrightarrow{\ f\ }M\xrightarrow{\ g\ }M''\longrightarrow 0$$
是 $D$-模的正合列。若给 $M$ 配置 $D$-模滤过 $F M$，它在 $M'$ 上诱导滤过 $F M'=(f^{-1}(f(M')\cap F_n M);\ n\in\mathbb{Z})$，在 $M''$ 上诱导滤过 $F M''=(g(F_n M);\ n\in\mathbb{Z})$。显然这些滤过是 $D$-模滤过。此外，序列
$$0\longrightarrow \Gr M'\xrightarrow{\ \Gr f\ }\Gr M\xrightarrow{\ \Gr g\ }\Gr M''\longrightarrow 0$$
是正合的。若滤过 $F M$ 是好的，则 $\Gr M$ 是有限生成 $\Gr D$-模，故 $\Gr M'$ 和 $\Gr M''$ 都是有限生成 $\Gr D$-模。由 3.1，$F M'$ 和 $F M''$ 是好滤过。因此我们证明了如下结果。

\begin{lemma}
设
$$0\longrightarrow M'\longrightarrow M\longrightarrow M''\longrightarrow 0$$
是 $D$-模的正合列。若 $F M$ 是 $M$ 上的好滤过，则诱导滤过 $F M'$ 和 $F M''$ 是好的。
\end{lemma}

由上述讨论
$$\lambda(\Gr_n M)=\lambda(\Gr_n M')+\lambda(\Gr_n M'')$$
对所有 $n\in\mathbb{Z}$ 成立。这由 $n$ 上的归纳蕴含
$$\lambda(F_n M)=\lambda(F_n M')+\lambda(F_n M'')$$
对所有 $n\in\mathbb{Z}$ 成立。这导出如下结果。

\begin{proposition}
设
$$0\longrightarrow M'\longrightarrow M\longrightarrow M''\longrightarrow 0$$
是有限生成 $D$-模的正合列。则
\begin{enumerate}[label=(\roman*)]
  \item $d_\lambda(M)=\max(d_\lambda(M'),d_\lambda(M''))$；
  \item 若 $d_\lambda(M)=d_\lambda(M')=d_\lambda(M'')$，则 $e_\lambda(M)=e_\lambda(M')+e_\lambda(M'')$。
\end{enumerate}
\end{proposition}

最后，设 $\varphi$ 是环 $D$ 的自同构且 $\varphi(D_0)=D_0$。可定义从 $D$-模范畴 $\mathcal{M}(D)$ 到自身的函子 $\tilde\varphi$，它把 $D$-模 $M$ 映为 $D$-模 $\tilde\varphi(M)$，其加法群结构不变，而 $D$ 的作用由 $(T,m)\longmapsto\varphi(T)m$（$T\in D$，$m\in M$）给出。显然 $\tilde\varphi$ 是范畴 $\mathcal{M}(D)$ 的自同构，且保持有限生成 $D$-模。

\begin{proposition}
设 $M$ 是有限生成 $D$-模。则
$$d_\lambda(\tilde\varphi(M))=d_\lambda(M).$$
\end{proposition}

\begin{proof}
设 $T_1,T_2,\dots,T_s$ 是 $D_1$ 中代表生成 $\Gr D$ 作为 $D_0$-代数的 $\Gr_1 D$ 中那些类的代表元。则存在 $d\in\mathbb{N}$ 使得对 $1\le i\le s$ 有 $\varphi(T_i)\in D_d$。由于 $T_1,T_2,\dots,T_s$ 和 $1$ 作为 $D_0$-模生成 $D_1$，我们断定 $\varphi(D_1)\subset D_d$。

设 $F M$ 是 $M$ 的好滤过。定义滤过 $F\tilde\varphi(M)$ 为
$$F_p\tilde\varphi(M)=F_{dp}M\quad (p\in\mathbb{Z}).$$
显然 $F\tilde\varphi(M)$ 是由有限生成 $D_0$-子模构成的 $\tilde\varphi(M)$ 的递增滤过。此外，对 $m\in\mathbb{Z}$，
$$D_1\cdot F_m\tilde\varphi(M)=\varphi(D_1)F_{dm}M\subset D_dF_{dm}M\subset F_{d(m+1)}M=F_{m+1}\tilde\varphi(M).$$
故由归纳，对所有 $n,m\in\mathbb{Z}$，
$$D_n\cdot F_m\tilde\varphi(M)=D_1\cdot D^{n-1}\cdot F_m\tilde\varphi(M)\subset D_1F_{m+n-1}\tilde\varphi(M)\subset F_{m+n}\tilde\varphi(M),$$
即 $F\tilde\varphi(M)$ 是 $D$-模滤过。由 3.6，存在比此滤过更细的好滤过 $F'\tilde\varphi(M)$，即存在 $k\in\mathbb{Z}_+$ 使得
$$F'_n\tilde\varphi(M)\subset F_{n+k}\tilde\varphi(M)=F_{d(n+k)}M$$
对所有 $n\in\mathbb{Z}$ 成立。因此
$$\lambda(F'_n\tilde\varphi(M))\le\lambda(F_{d(n+k)}M)$$
对 $n\in\mathbb{Z}$ 成立。对充分大的 $n\in\mathbb{Z}$，$\lambda(F_{d(n+k)}M)$ 等于 $n$ 的多项式，其首项为
$$\frac{e_\lambda(M)d^{d_\lambda(M)}}{d_\lambda(M)!}\,n^{d_\lambda(M)}.$$
由于对充分大的 $n\in\mathbb{Z}$，$\lambda(F'_n\tilde\varphi(M))$ 也由次数为 $d_\lambda(\tilde\varphi(M))$ 的多项式给出，我们断定 $d_\lambda(\tilde\varphi(M))\le d_\lambda(M)$。对 $\varphi^{-1}$ 应用同样推理也得到
$$d_\lambda(M)=d_\lambda(\tilde\varphi^{-1}(\tilde\varphi(M)))\le d_\lambda(\tilde\varphi(M)).$$
\end{proof}
\begin{annote}
“好滤过”者，指滤过模 $M$ 的分次 $\mathrm{gr}M$ 在 $\mathrm{gr}A$ 上有限生成；等价说法是 $M$ 为有限个 $A_{\le d}$ 元素的 $A$-线性组合。好滤过的选择不影响 $\mathrm{gr}M$ 的 Hilbert 函数的主项，这正是文中可放心用 $\mathrm{gr}M$ 替代 $M$ 计算维数的原因：两条好滤过生成的 $\mathrm{gr}$ 只差有限维向量空间。[来源：原书 §3, p15–22；译者按]
\end{annote}



\section{多项式环上模的维数}

设 $A=k[X_1,\dots,X_n]$，其中 $k$ 是代数闭域。可按多项式的次数滤过 $A$，即令
$$A_m=\Bigl\{\sum_I c_I x^I\ \Big|\ c_I\in k,\ |I|\le m\Bigr\}.$$
则 $\Gr A=k[X_1,\dots,X_n]$，故 $A$ 满足前一节的性质 (i)--(vii)。

由于 $A^0=k$，可取函数 $\dim_k$ 作为可加函数 $\lambda$。这导出有限生成 $A$-模 $M$ 的维数 $d(M)$ 和重数 $e(M)$ 的概念。已知对任意 $p\in\mathbb{Z}_+$，
$$\dim_k A_p=\binom{n+p}{n}=\frac{p^n}{n!}+\text{（关于 $p$ 的低阶项）},$$
即 $d(A)=n$，$e(A)=1$。此外，对任意有限生成 $A$-模 $M$ 有正合列
$$0\longrightarrow K\longrightarrow A^p\longrightarrow M\longrightarrow 0,$$
故由 3.9 有 $d(M)\le n$。以后将给出 $d(M)$ 的几何解释。

设 $x\in k^n$，记 $\mathfrak{m}_x$ 为 $k[X_1,\dots,X_n]$ 中所有在 $x$ 处消失的多项式构成的极大理想。记 $A_x$ 为 $A$ 在 $x$ 处的局部化，即 $k^n$ 上在 $x$ 处正则的所有有理 $k$-值函数环。如 2.11 中所见，$A_x$ 是以 $\mathfrak{n}_x=(\mathfrak{m}_x)_x$ 为极大理想的 $n$ 维正则局部环，$\mathfrak{n}_x$ 由 $k^n$ 上在 $x$ 处消失的所有有理 $k$-值函数组成。设 $M$ 是 $A$-模。它在 $x$ 处的局部化 $M_x$ 是 $A_x$-模。定义 $M$ 的\emph{支集}为
$$\operatorname{supp}(M)=\{x\in k^n\mid M_x\neq 0\}.$$

\begin{lemma}
设
$$0\longrightarrow M'\longrightarrow M\longrightarrow M''\longrightarrow 0$$
是 $A$-模的正合列。则
$$\operatorname{supp}(M)=\operatorname{supp}(M')\cup\operatorname{supp}(M'').$$
\end{lemma}

\begin{proof}
由局部化的正合性可知
$$0\longrightarrow M'_x\longrightarrow M_x\longrightarrow M''_x\longrightarrow 0$$
是 $A_x$-模的正合列。这立即蕴含我们的结论。
\end{proof}

对理想 $I\subset k[X_1,\dots,X_n]$，记 $V(I)=\{x\in k^n\mid f(x)=0\ \text{对 }f\in I\}$。

\begin{proposition}
设 $M$ 是有限生成 $A$-模，$I$ 是它在 $A$ 中的零化子。则 $\operatorname{supp}(M)=V(I)$。
\end{proposition}

\begin{proof}
对 $M$ 的生成元个数归纳证明。

先设 $M$ 有一个生成元，即 $M=A/I$。则 $M_x=(A/I)_x=A_x/I_x$。设 $x\in V(I)$。则 $I\subset\mathfrak{m}_x$，$I_x\subset\mathfrak{n}_x$，故 $I_x\neq A_x$。由此可知 $(A/I)_x\neq 0$，即 $x\in\operatorname{supp}(M)$。反之，若 $x\notin V(I)$，存在 $f\in I$ 使得 $f(x)\neq 0$，即 $f\notin\mathfrak{m}_x$。因此 $f$ 在局部环 $A_x$ 中可逆，而 $f\in I_x$ 蕴含 $I_x=A_x$。故 $(A/I)_x=0$，$x\notin\operatorname{supp}(A/I)$。因此 $\operatorname{supp}(A/I)=V(I)$。

现考虑一般情形。设 $m_1,\dots,m_p$ 是 $M$ 的一组生成元。记 $M'$ 为由 $m_1,\dots,m_{p-1}$ 生成的子模。则有正合列
$$0\longrightarrow M'\longrightarrow M\longrightarrow M''\longrightarrow 0$$
且 $M''$ 是循环模。此外由 4.1，$\operatorname{supp}(M)=\operatorname{supp}(M')\cup\operatorname{supp}(M'')$。故由归纳假设，$\operatorname{supp}(M)=V(I')\cup V(I'')$，其中 $I'$ 和 $I''$ 分别是 $M'$ 和 $M''$ 的零化子。

显然 $I'\cdot I''$ 含于 $M$ 的零化子 $I$ 中。另一方面，$I$ 零化 $M'$ 和 $M''$，故 $I\subset I'\cap I''$。于是
$$I'\cdot I''\subset I\subset I'\cap I''.$$
这意味着
$$V(I')\cup V(I'')\subset V(I'\cap I'')\subset V(I)\subset V(I'\cdot I'').$$
设 $x\notin V(I')\cup V(I'')$。则存在 $f\in I'$ 和 $g\in I''$ 使得 $f(x)\neq 0$ 且 $g(x)\neq 0$。由此可知 $(fg)(x)=f(x)g(x)\neq 0$，$x\notin V(I'\cdot I'')$。故 $V(I'\cdot I'')\subset V(I')\cup V(I'')$，上述所有包含均为等式。因此 $V(I)=V(I')\cup V(I'')$ 且 $\operatorname{supp}(M)=V(I)$。
\end{proof}

这立即蕴含如下推论。

\begin{corollary}
设 $M$ 是有限生成 $A$-模。则支集 $\operatorname{supp}(M)$ 是 $k^n$ 中的 Zariski 闭子集。
\end{corollary}

下一个引理在某些归约论证中很有用。

\begin{lemma}
设 $B$ 是诺特交换环，$M\neq 0$ 是有限生成 $B$-模。则存在 $M$ 由 $B$-子模构成的滤过
$$0=M_0\subset M_1\subset\cdots\subset M_{n-1}\subset M_n=M$$
以及 $B$ 的素理想 $\mathfrak{J}_i$，使得对 $1\le i\le n$ 有 $M_i/M_{i-1}\cong B/\mathfrak{J}_i$。
\end{lemma}

\begin{proof}
对 $x\in M$ 令 $\operatorname{Ann}(x)=\{a\in B\mid ax=0\}$。设 $\mathcal{A}$ 是所有这种理想 $\operatorname{Ann}(x)$（$x\in M$，$x\neq 0$）的族。由于 $B$ 是诺特环，$\mathcal{A}$ 有极大元。设 $I$ 是 $\mathcal{A}$ 中的极大元。我们断言 $I$ 是素的。设 $x\in M$ 使得 $I=\operatorname{Ann}(x)$。则 $ab\in I$ 蕴含 $abx=0$。假设 $b\notin I$，即 $bx\neq 0$。则 $I\subset\operatorname{Ann}(bx)$ 且 $a\in\operatorname{Ann}(bx)$。由 $I$ 的极大性，$a\in\operatorname{Ann}(bx)=I$，故 $I$ 是素的。因此存在 $x\in M$ 使得 $\mathfrak{J}_1=\operatorname{Ann}(x)$ 是素的。若令 $M_1=Bx$，则 $M_1\cong B/\mathfrak{J}_1$。

现记 $\mathcal{F}$ 为 $M$ 的所有具有滤过 $0=N_0\subset N_1\subset\cdots\subset N_k=N$ 且 $N_i/N_{i-1}\cong B/\mathfrak{J}_i$（对某素理想 $\mathfrak{J}_i$）的 $B$-子模 $N$ 的族。由于 $M$ 是诺特模，$\mathcal{F}$ 含有一个极大元 $L$。假设 $L\neq M$。则会有正合列
$$0\longrightarrow L\longrightarrow M\longrightarrow L'\longrightarrow 0,$$
由证明的第一部分，$L'$ 会有形如 $B/\mathfrak{J}'$（对某素理想 $\mathfrak{J}'$）的子模 $N'$，与 $L$ 的极大性矛盾。故 $L=M$。这证明了具有所需性质的滤过的存在性。
\end{proof}

\begin{theorem}
设 $M$ 是有限生成 $A$-模，$\operatorname{supp}(M)$ 是其支集。则
$$d(M)=\dim\operatorname{supp}(M).$$
\end{theorem}

这一结果有如下局部版本。$A$ 在 $x\in k^n$ 处的局部化 $A_x$ 是诺特局部环。此外其极大理想 $\mathfrak{n}_x$ 由多项式 $X_i-x_i$（$1\le i\le n$）生成，且它们在 $\mathfrak{n}_x/\mathfrak{n}_x^2$ 中的像作为 $k$-向量空间张成它。因此 $X_i-x_i$（$1\le i\le n$）构成 $A_x$ 中的坐标系。对任意有限生成 $A$-模 $M$，其在 $x$ 处的局部化 $M_x$ 是有限生成 $A_x$-模，故可考虑其维数 $d(M_x)$。

对 $k$ 上任意代数簇 $V$ 和 $x\in V$，记 $\dim_x V$ 为 $V$ 在 $x$ 处的局部维数。

\begin{theorem}
设 $M$ 是有限生成 $A$-模，$x\in\operatorname{supp}(M)$。则
$$d(M_x)=\dim_x(\operatorname{supp}(M)).$$
\end{theorem}

同时证明 4.5 和 4.6。首先注意，若有 $A$-模的正合列
$$0\longrightarrow M'\longrightarrow M\longrightarrow M''\longrightarrow 0$$
且 4.5、4.6 对 $M'$ 和 $M''$ 成立，则由 3.9 和 4.1，
\begin{align*}
d(M)&=\max(d(M'),d(M''))=\max(\dim\operatorname{supp}(M'),\dim\operatorname{supp}(M''))\\
&=\dim(\operatorname{supp}(M')\cup\operatorname{supp}(M''))=\dim\operatorname{supp}(M).
\end{align*}
此外，对任意 $x\in\operatorname{supp}(M)$，由局部化的正合性有正合列
$$0\longrightarrow M'_x\longrightarrow M_x\longrightarrow M''_x\longrightarrow 0;$$
故由 2.7 和 4.1，
\begin{align*}
d(M_x)&=\max(d(M'_x),d(M''_x))=\max(\dim_x\operatorname{supp}(M'),\dim_x\operatorname{supp}(M''))\\
&=\dim_x(\operatorname{supp}(M')\cup\operatorname{supp}(M''))=\dim_x\operatorname{supp}(M).
\end{align*}

假设 4.5 和 4.6 对所有 $M=A/J$（$J$ 为素理想）成立。则前面的注记、4.4 以及滤过长度上的归纳将证明一般情形的陈述。

故可设 $M=A/J$，$J$ 为素理想。先设 $J$ 使得 $A/J$ 是 $k$ 上的有限维向量空间。则 $A/J$ 是整环且在 $k$ 上是整的，故它是 $k$ 的代数扩张的域。由于 $k$ 代数闭，$A/J=k$，$J$ 是极大理想。这时由 Hilbert 零点定理，$\operatorname{supp}(M)=V(J)$ 是 $k^n$ 中一点 $x$，即 $\dim\operatorname{supp}(M)=0$。另一方面，由于 $M_x$ 是一维线性空间，$d(M_x)=0$，结论显然。于是可设 $J$ 在 $A$ 中不是有限余维的，特别地不是极大理想。设 $J_1\supsetneq J$ 是不同于 $J$ 的素理想。则存在 $f\in J_1$ 使得 $f\notin J$。由此可知 $J\subset (f)+J\subset J_1$ 且 $J\neq(f)+J$。因此 $A/J_1$ 是 $A/((f)+J)$ 的商，而 $A/((f)+J)$ 是 $A/J$ 的商。此外 $A/((f)+J)=M/fM$。考虑 $M$ 上由乘以 $f$ 给出的自同态。若 $g+J$ 在此映射的核中，则 $0=f(g+J)=fg+J$，$fg\in J$。由于 $J$ 是素理想且 $f\notin J$，得 $g\in J$，$g+J=0$，此映射是单射。因此有 $A$-模的正合列
$$0\longrightarrow M\xrightarrow{\ f\ }M\longrightarrow M/fM\longrightarrow 0.$$
由 3.9 这意味着 $d(M/fM)\le d(M)$。若 $d(M/fM)=d(M)$，还会有 $e(M)=e(M)+e(M/fM)$，故 $e(M/fM)=0$。这只有在 $d(M/fM)=0$ 时才可能，而这时还会蕴含 $d(M)=0$、$M$ 有限维，与假设矛盾。因此 $d(M/fM)<d(M)$。由于 $A/J_1$ 是 $M/fM$ 的商，这意味着 $d(A/J_1)<d(A/J)$。

设 $x\in V(J_1)$。则由局部化得到 $A_x$-模的正合列
$$0\longrightarrow M_x\xrightarrow{\ f\ }M_x\longrightarrow M_x/fM_x\longrightarrow 0.$$
由 2.7 这意味着 $d(M_x/fM_x)\le d(M_x)$。若 $d(M_x/fM_x)=d(M_x)$，还会有 $e(M_x)=e(M_x)+e(M_x/fM_x)$，故 $e(M_x/fM_x)=0$。这只有在 $d(M_x/fM_x)=0$ 时才可能，而这时会蕴含 $\mathfrak{m}_x(M_x/fM_x)=M_x/fM_x$，由 Nakayama 引理 $M_x/fM_x=0$。这将推出乘以 $f$ 在 $M_x$ 上是满射，且由于 $f\in\mathfrak{m}_x$，由 Nakayama 引理这会蕴含 $M_x=0$，与假设矛盾。因此 $d(M_x/fM_x)<d(M_x)$。由于 $A/J_1$ 是 $M/fM$ 的商，这意味着 $d((A/J_1)_x)<d((A/J)_x)$。

设
$$Z_0=\{x\}\subset Z_1\subset\cdots\subset Z_{n-1}\subset Z_n=k^n$$
是 $k^n$ 中非空不可约闭子集的极大链。则
$$I(Z_0)=\mathfrak{m}_x\supset I(Z_1)\supset\cdots\supset I(Z_{n-1})\supset I(Z_n)=\{0\}$$
是 $A$ 中素理想的极大链。由上述论证有以下严格不等式链
$$0\le d(A/I(Z_0))<d(A/I(Z_1))<\cdots<d(A/I(Z_n))=d(A)=n,$$
以及
$$0\le d((A/I(Z_0))_x)<d((A/I(Z_1))_x)<\cdots<d((A/I(Z_n))_x)=d(A_x)=n,$$
（由 2.11）。由此可知
$$d((A/I(Z_j))_x)=d(A/I(Z_j))=j=\dim Z_j\quad\text{对 }0\le j\le n.$$
由于每个闭不可约子集 $Z$ 都可置于极大链中，可知对任意闭不可约子集 $Z\subset k^n$ 和任意 $x\in Z$ 有 $d((A/I(Z))_x)=d(A/I(Z))=\dim Z$。另一方面，这意味着对 $A$ 中任意素理想 $J$ 和 $x\in V(J)$ 有 $d((A/J)_x)=d(A/J)=\dim V(J)$。由 4.2，这完成了 4.5 和 4.6 的证明。

下一个结果由 4.5 和 4.6 立即得出。

\begin{corollary}
设 $M$ 是有限生成 $A$-模。则
$$d(M)=\sup_{x\in\operatorname{supp}(M)} d(M_x).$$
\end{corollary}

最后证明一个以后需要的结果。

\begin{lemma}
设 $I$ 是 $A$ 中的理想。则 $\dim V(I)=\dim V(\Gr I)$。
\end{lemma}

\begin{proof}
$A$-模的短正合列
$$0\longrightarrow I\longrightarrow A\longrightarrow A/I\longrightarrow 0$$
（其中各模配置由 $A$ 的自然滤过诱导的滤过）导出分次 $A$-模的短正合列
$$0\longrightarrow\Gr I\longrightarrow A\longrightarrow\Gr(A/I)\longrightarrow 0.$$
因此
\begin{align*}
\dim_k F_p(A/I)&=\sum_{q=0}^{p}\bigl(\dim_k F_q(A/I)-\dim_k F_{q-1}(A/I)\bigr)\\
&=\sum_{q=0}^{p}\dim_k\Gr_q(A/I)=\sum_{q=0}^{p}\bigl(\dim_k\Gr_q A-\dim\Gr_q I\bigr)\\
&=\dim_k F_p A-\dim_k F_p\Gr I=\dim_k F_p(A/\Gr I).
\end{align*}
因此 $d(A/I)=d(A/\Gr I)$。结论由 4.4 和 4.5 得出。
\end{proof}

\section{多项式系数微分算子环}

设 $k$ 是特征 $0$ 的域，$A$ 是 $k$ 上的交换代数。设 $\End_k(A)$ 是 $A$ 的所有 $k$-线性自同态构成的代数。以换位子 $[S,T]=ST-TS$（对 $S,T\in\End_k(A)$）它是一个 Lie 代数。显然 $\End_k(A)$ 包含 $A$-线性自同态的集合 $\End_A(A)$ 作为子代数。对任意 $a\in A$，可附着由 $b\longmapsto ab$（$b\in A$）定义的 $A$-线性自同态。由于此自同态在 $1$ 处取值 $a$，该映射显然是代数的单射。

另一方面，若 $T\in\End_A(A)$，则对任意 $b\in A$ 有
$$T(b)=bT(1)=T(1)b,$$
即 $T$ 由乘以 $T(1)$ 给出。这导出如下结果。

\begin{lemma}
把 $a\in A$ 映为 $A$-线性自同态 $b\longmapsto ab$（$b\in A$）的代数同态是 $A$ 到 $\End_A(A)$ 上的同构。
\end{lemma}

以下把 $A$ 与 $\End_k(A)$ 的子代数 $\End_A(A)$ 等同。

$A$ 的一个 $k$-导子是满足
$$T(ab)=T(a)b+aT(b)\quad\text{对任意 }a,b\in A$$
的 $T\in\End_k(A)$。特别地 $[T,a](b)=T(ab)-aT(b)=T(a)b$，即对任意 $a\in A$ 有 $[T,a]=T(a)\in A$。这意味着对任意 $a_0,a_1\in A$ 有 $[[T,a_0],a_1]=0$。

这引出如下定义。设 $n\in\mathbb{Z}_+$。若对任意 $a_0,a_1,\dots,a_n\in A$ 有
$$[\cdots[[T,a_0],a_1],\dots,a_n]=0,$$
则称 $T\in\End_k(A)$ 是 $A$ 上阶数 $\le n$ 的（$k$-线性）微分算子。记 $\Diff_k(A)$ 为 $A$ 上所有微分算子的空间。

\begin{lemma}
设 $T,S$ 分别是阶数 $\le n$、$\le m$ 的两个微分算子。则 $T\circ S$ 是阶数 $\le n+m$ 的微分算子。
\end{lemma}

\begin{proof}
对 $n+m$ 归纳证明。若 $n=m=0$，则 $T,S\in\End_A(A)$，故 $T\circ S\in\End_A(A)$，是阶数 $0$ 的微分算子。

现设 $n+m\le p$。则
$$[T\circ S,a]=TSa-aTS=T[S,a]+[T,a]S,$$
且 $[T,a]$ 和 $[S,a]$ 分别是阶数 $\le n-1$ 和 $\le m-1$ 的微分算子。由归纳假设，此微分算子阶数 $\le n+m-1$。因此 $T\circ S$ 阶数 $\le n+m$。
\end{proof}

因此 $\Diff_k(A)$ 是 $\End_k(A)$ 的子代数。我们称它为 $A$ 上所有 $k$-线性微分算子的代数。此外，令 $n<0$ 时 $F_n\Diff_k(A)=\{0\}$，
$$F_n\Diff_k(A)=\{T\in\Diff_k(A)\mid \operatorname{order}(T)\le n\}\quad(n\ge 0).$$
显然这是 $\Diff_k(A)$ 关于 $k$-向量子空间的递增穷尽滤过。此滤过与 $\Diff_k(A)$ 的环结构相容，即对任意 $n,m\in\mathbb{Z}$，
$$F_n\Diff_k(A)\circ F_m\Diff_k(A)\subset F_{n+m}\Diff_k(A).$$

\begin{lemma}\ 
\begin{enumerate}[label=(\roman*)]
  \item $F_0\Diff_k(A)=A$；
  \item $F_1\Diff_k(A)=\Der_k(A)\oplus A$；
  \item 对任意 $n,m\in\mathbb{Z}_+$，$[F_n\Diff_k(A),F_m\Diff_k(A)]\subset F_{n+m-1}\Diff_k(A)$。
\end{enumerate}
\end{lemma}

\begin{proof}
(i) 显然。

(ii) 如前所述，$\Der_k(A)\subset F_1\Diff_k(A)$。此外，对任意 $T\in\Der_k(A)$ 有 $T(1)=T(1\cdot 1)=2T(1)$，故 $T(1)=0$。这意味着 $\Der_k(A)\cap A=0$。

设 $S\in F_1\Diff_k(A)$，$T=S-S(1)$。则 $T(1)=0$，故 $T(a)=[T,a](1)$，且
$$T(ab)=[T,ab](1)=([T,a]b)(1)+(a[T,b])(1)=(b[T,a])(1)+(a[T,b])(1)=T(a)b+aT(b),$$
即 $T\in\Der_k(A)$。

(iii) 设 $T,S$ 分别是阶数 $\le n$、$\le m$ 的算子。我们断言 $[T,S]$ 阶数 $\le n+m-1$。对 $n+m$ 归纳证明。若 $n=m=0$，无需证明。一般地，由 Jacobi 恒等式，
$$[[T,S],a]=[[T,a],S]+[T,[S,a]],$$
其中 $[T,a]$ 和 $[S,a]$ 分别是阶数 $\le n-1$ 和 $\le m-1$ 的算子。故由归纳假设，$[[T,S],a]$ 阶数 $\le n+m-2$，$[T,S]$ 阶数 $\le n+m-1$。
\end{proof}

这意味着分次环 $\Gr\Diff_k(A)$ 是交换 $A$-代数。此外，$\Diff_k(A)$ 满足 §3 的性质 (i)--(v)。

设 $n\ge 1$。设 $T$ 是 $A$ 上阶数 $\le n$ 的微分算子。则可定义从 $A^n$ 到 $\Diff_k(A)$ 的映射
$$\sigma_n(T)(a_1,a_2,\dots,a_{n-1},a_n)=[\cdots[[T,a_1],a_2],\dots,a_{n-1}],a_n].$$
由于 $\sigma_n(T)(a_1,\dots,a_n)$ 阶数 $\le 0$，可把此映射看作从 $A^n$ 到 $A$ 的映射。

\begin{lemma}
设 $T$ 是 $A$ 上阶数 $\le n$ 的微分算子。则：
\begin{enumerate}[label=(\roman*)]
  \item 映射 $\sigma_n(T):A^n\to A$ 是对称 $k$-多重线性映射；
  \item 算子 $T$ 阶数 $\le n-1$ 当且仅当 $\sigma_n(T)=0$。
\end{enumerate}
\end{lemma}

\begin{proof}
(i) 只需检验对称性。为此注意，由 Jacobi 恒等式，对任意 $S\in\Diff_k(A)$ 和 $a,b\in A$ 有
$$[[S,a],b]=[[S,b],a].$$
这意味着
\begin{align*}
&\sigma_n(T)(a_1,\dots,a_i,a_{i+1},\dots,a_n)\\
&=[\cdots[[\cdots[[T,a_1],a_2],\dots,a_i],a_{i+1}]\dots,a_{n-1}],a_n]\\
&=[\cdots[[\cdots[[T,a_1],a_2],\dots,a_{i+1}],a_i]\dots,a_{n-1}],a_n]\\
&=\sigma_n(T)(a_1,\dots,a_{i+1},a_i,\dots,a_n),
\end{align*}
故 $\sigma(T)$ 是对称的。

(ii) 显然。
\end{proof}

现在讨论一个特例。设 $A=k[X_1,X_2,\dots,X_n]$。则记 $D(n)=\Diff_k(A)$，称 $D(n)$ 为 $k^n$ 上所有微分算子的代数。设 $\partial_1,\partial_2,\dots,\partial_n$ 是 $k[X_1,X_2,\dots,X_n]$ 的标准导子。对 $I,J\in\mathbb{Z}_+^n$ 令
$$X^I=X_1^{i_1}X_2^{i_2}\cdots X_n^{i_n},\qquad \partial^J=\partial_1^{j_1}\partial_2^{j_2}\cdots\partial_n^{j_n}.$$
则 $X^I\partial^J\in D(n)$，且它是阶数 $\le |J|=j_1+j_2+\cdots+j_n$ 的微分算子。此外，若 $T$ 是由
$$T=\sum_{|I|\le p}P_I(X_1,X_2,\dots,X_n)\partial^I$$
给出的微分算子，其中 $P_I\in k[X_1,X_2,\dots,X_n]$，则 $T$ 阶数 $\le p$。

\begin{lemma}
导子 $(\partial_i;\ 1\le i\le n)$ 构成自由 $k[X_1,\dots,X_n]$-模 $\Der_k(k[X_1,X_2,\dots,X_n])$ 的一组基。
\end{lemma}

\begin{proof}
设 $T\in\Der_k(k[X_1,X_2,\dots,X_n])$。令 $P_i=T(X_i)$（$1\le i\le n$），定义 $S=\sum_{i=1}^n P_i\partial_i$。显然对所有 $1\le i\le n$，
$$S(X_i)=\sum_{j=1}^n P_j\partial_j(X_i)=P_i=T(X_i).$$
由于 $X_1,X_2,\dots,X_n$ 作为 $k$-代数生成 $k[X_1,X_2,\dots,X_n]$，可得 $T=S$。因此 $(\partial_i;\ 1\le i\le n)$ 生成 $k[X_1,X_2,\dots,X_n]$-模 $\Der_k(k[X_1,X_2,\dots,X_n])$。若 $\sum_{i=1}^n Q_i\partial_i=0$，其中 $Q_i\in k[X_1,X_2,\dots,X_n]$，则对所有 $1\le i\le n$ 有 $0=(\sum_j Q_j\partial_j)(X_i)=Q_i$。这说明 $\partial_i$（$1\le i\le n$）是自由生成元。
\end{proof}

设 $T$ 是 $k[X_1,X_2,\dots,X_n]$ 上阶数 $\le p$ 的微分算子。若 $p<0$，$T=0$，令 $\Symb_p(T)=0$。若 $p=0$，$T\in A$，令 $\Symb_0(T)=T$。对 $p\ge 1$，按如下方式定义 $k[X_1,\dots,X_n,\xi_1,\dots,\xi_n]$ 中的多项式 $\Symb_p(T)$。设 $(\xi_1,\dots,\xi_n)\in k^n$。定义线性多项式 $\ell_\xi=\sum_{i=1}^n\xi_i X_i\in k[X_1,X_2,\dots,X_n]$，以及函数
$$(\xi_1,\xi_2,\dots,\xi_n)\longmapsto \frac{1}{p!}\sigma_p(T)(\ell_\xi,\ell_\xi,\dots,\ell_\xi)$$
在 $k^n$ 上取值于 $k[X_1,X_2,\dots,X_n]$。显然可把此函数看作关于 $X_1,\dots,X_n,\xi_1,\dots,\xi_n$ 的多项式，关于 $\xi_1,\dots,\xi_n$ 齐次 $p$ 次，记为 $\Symb_p(T)$。多项式 $\Symb_p(T)$ 称为微分算子 $T$ 的 \emph{$p$-主符号}。由定义，当 $T$ 阶数 $<p$ 时 $\Symb_p(T)$ 为零。因此对 $p\ge 0$，它诱导 $\Gr_p D(n)$ 到 $k[X_1,\dots,X_n,\xi_1,\dots,\xi_n]$ 的 $k$-线性映射。记相应的 $\Gr D(n)$ 到 $k[X_1,\dots,X_n,\xi_1,\dots,\xi_n]$ 的 $k$-线性映射为 $\Symb$。

\begin{theorem}
映射 $\Symb:\Gr D(n)\longrightarrow k[X_1,X_2,\dots,X_n,\xi_1,\xi_2,\dots,\xi_n]$ 是 $k$-代数同构。
\end{theorem}

此结果的证明分几步。先证明符号映射是代数同态。

\begin{lemma}
设 $T,S\in D(n)$ 分别是阶数 $\le p$ 和 $\le q$ 的算子。则
$$\Symb_{p+q}(TS)=\Symb_p(T)\,\Symb_q(S).$$
\end{lemma}

\begin{proof}
设 $\xi\in k^n$，定义映射 $\tau_\xi:D(n)\to D(n)$ 为 $\tau_\xi(T)=[T,\ell_\xi]$。则
$$\tau_\xi(TS)=[TS,\ell_\xi]=TS\ell_\xi-\ell_\xi TS=[T,\ell_\xi]S+T[S,\ell_\xi]=\tau_\xi(T)S+T\tau_\xi(S).$$
因此对任意 $k\in\mathbb{Z}_+$，
$$\tau_\xi^k(TS)=\sum_{i=0}^k\binom{k}{i}\tau_\xi^{k-i}(T)\,\tau_\xi^i(S).$$
这意味着
\begin{align*}
\Symb_{p+q}(TS)&=\frac{1}{(p+q)!}\sigma_{p+q}(T)(\ell_\xi,\ell_\xi,\dots,\ell_\xi)=\frac{1}{(p+q)!}\tau_\xi^{p+q}(TS)\\
&=\frac{1}{p!\,q!}\tau_\xi^p(T)\,\tau_\xi^q(S)=\Symb_p(T)\,\Symb_q(S).
\end{align*}
\end{proof}

由于 $\Symb_0(X_i)=X_i$，$\Symb_1(\partial_i)=\xi_i$（$1\le i\le n$），可见对 $X^I\partial^J$（其中 $p=|J|$）有
$$\Symb_p(X^I\partial^J)=X^I\xi^J.$$
特别地，对
$$T=\sum_{|J|\le p}P_I(X_1,X_2,\dots,X_n)\partial^I$$
（其中 $P_I\in k[X_1,X_2,\dots,X_n]$），有
$$\Symb_p(T)=\sum_{|I|=p}P_I(X_1,X_2,\dots,X_n)\xi^I.$$
因此符号同态是满射。还需证明符号映射是单射。

\begin{lemma}
设 $T\in F_p D(n)$。则 $\Symb_p(T)=0$ 当且仅当 $T$ 阶数 $\le p-1$。
\end{lemma}

\begin{proof}
对 $p$ 归纳证明。$p=0$ 时显然。故设 $p>0$。设 $\xi\in k^n$，定义映射 $\tau_\xi:D(n)\to D(n)$ 为 $\tau_\xi(T)=[T,\ell_\xi]$。则对任意 $\lambda\in k$ 和 $\eta\in k^n$，
$$\tau_{\xi+\lambda\eta}(T)=[T,\ell_{\xi+\lambda\eta}]=[T,\ell_\xi]+\lambda[T,\ell_\eta]=\tau_\xi(T)+\lambda\tau_\eta(T).$$
由于 $\tau_\xi$ 和 $\tau_\eta$ 可交换，可见对任意 $k\in\mathbb{Z}_+$，
$$\tau_{\xi+\lambda\eta}^k(T)=\sum_{i=0}^k\binom{k}{i}\lambda^i\tau_\xi^{k-i}(\tau_\eta^i(T)).$$
由假设，对任意 $\lambda\in k$ 有 $\tau_{\xi+\lambda\eta}^p(T)=0$。因此由于域 $k$ 是无限的，对 $0\le i\le p$ 有 $\tau_\xi^{p-i}(\tau_\eta^i(T))=0$。特别地，对任意 $\xi,\eta\in k^n$ 有 $\tau_\xi^{p-1}(\tau_\eta(T))=0$。这意味着对任意 $\eta\in k^n$ 有 $\Symb_{p-1}([T,\ell_\eta])=0$，特别地
$$\Symb_{p-1}([T,X_i])=0$$
对 $1\le i\le n$ 成立。由归纳假设，$[T,X_i]$（$1\le i\le n$）阶数 $\le p-2$。设 $P,Q\in k[X_1,X_2,\dots,X_n]$。则
$$[T,PQ]=TPQ-PQT=[T,P]Q+P[T,Q],$$
故 $[T,PQ]$ 的阶数不超过 $[T,P]$ 和 $[T,Q]$ 阶数的最大值。由于 $X_i$（$1\le i\le n$）生成 $k[X_1,X_2,\dots,X_n]$，我们断定对任意多项式 $P$，$[T,P]$ 的阶数 $\le p-2$。这意味着 $T$ 的阶数 $\le p-1$。
\end{proof}

这也结束了 5.6 的证明。特别地，我们看到 $D(n)$ 满足 §3 的性质 (i)--(vii)。由 3.5 立即推出如下结果。

\begin{theorem}
环 $D(n)$ 是右诺特且左诺特的。
\end{theorem}

\begin{corollary}
$(X^I\partial^J;\ I,J\in\mathbb{Z}_+^n)$ 构成 $D(n)$ 作为 $k$-向量空间的一组基。
\end{corollary}

\begin{proof}
若 $|J|=p$，$X^I\partial^J$ 的 $p$-主符号等于 $X^I\xi^J$，而 $(X^I\xi^J;\ I,J\in\mathbb{Z}_+^n)$ 构成 $k[X_1,\dots,X_n,\xi_1,\dots,\xi_n]$ 作为 $k$-向量空间的一组基。
\end{proof}

$D(n)$ 的如下刻画经常有用。

\begin{theorem}
$k$-代数 $D(n)$ 是由 $X_1,X_2,\dots,X_n$ 和 $\partial_1,\partial_2,\dots,\partial_n$ 生成、满足定义关系 $[X_i,X_j]=0$、$[\partial_i,\partial_j]=0$ 以及 $[\partial_i,X_j]=\delta_{ij}$（对所有 $1\le i,j\le n$）的 $k$-代数。
\end{theorem}

\begin{proof}
设 $B$ 是由 $X_1,X_2,\dots,X_n$ 和 $\partial_1,\partial_2,\dots,\partial_n$ 生成、满足定义关系 $[X_i,X_j]=0$、$[\partial_i,\partial_j]=0$ 以及 $[\partial_i,X_j]=\delta_{ij}$（对所有 $1\le i,j\le n$）的 $k$-代数。由于这些关系在 $D(n)$ 中成立，且 $D(n)$ 由 $X_1,X_2,\dots,X_n$ 和 $\partial_1,\partial_2,\dots,\partial_n$ 生成，我们断定存在唯一的从 $B$ 到 $D(n)$ 上的满射，把生成元映为相应的生成元。显然 $B$ 由 $(X^I\partial^J;\ I,J\in\mathbb{Z}_+^n)$ 张成。因此由 5.10，此同态也是单射。
\end{proof}

\begin{proposition}
$D(n)$ 的中心等于 $k\cdot 1$。
\end{proposition}

\begin{proof}
设 $T$ 是 $D(n)$ 的中心元。则对任意多项式 $P$ 有 $[T,P]=0$，且 $T$ 阶数 $\le 0$。因此由 5.3，$T\in k[X_1,X_2,\dots,X_n]$。另一方面，$0=[\partial_i,T]=\partial_i(T)$（$1\le i\le n$）。这意味着 $T$ 是常数多项式。
\end{proof}

设 $D(n)^\circ$ 是 $D(n)$ 的反代数。则由 5.11，存在唯一的同构 $\varphi:D(n)^\circ\to D(n)$，定义为 $\varphi(X_i)=X_i$，$\varphi(\partial_i)=-\partial_i$（$1\le i\le n$）。此同态 $\varphi$ 称为 $D(n)$ 的\emph{主反自同构}。这证明了如下结果。

\begin{proposition}
代数 $D(n)^\circ$ 同构于 $D(n)$。
\end{proposition}

此外，由 5.11 可定义 $D(n)$ 的自同构 $F$：$F(X_i)=\partial_i$，$F(\partial_i)=-X_i$（$1\le i\le n$）。此自同构称为 $D(n)$ 的\emph{Fourier 自同构}。$F$ 的平方 $F^2$ 是 $D(n)$ 的自同构 $\iota$，作用为 $\iota(X_i)=-X_i$，$\iota(\partial_i)=-\partial_i$（$1\le i\le n$）。显然 $\iota^2=1$。

与按微分算子阶数的滤过不同，$D(n)$ 还有另一个与其环结构相容的滤过，它不能在更一般的微分算子环上定义。对 $p\in\mathbb{Z}$ 令
$$D_p(n)=\Bigl\{\sum_{|I|+|J|\le p}a_{IJ}X^I\partial^J\Bigr\}.$$
显然 $(D_p(n)\mid p\in\mathbb{Z})$ 是 $D(n)$ 关于 $k$ 上有限维向量空间的递增穷尽滤过。

\begin{lemma}
对任意 $p,q\in\mathbb{Z}$ 有
\begin{enumerate}[label=(\roman*)]
  \item $D_p(n)\circ D_q(n)\subset D_{p+q}(n)$；
  \item $[D_p(n),D_q(n)]\subset D_{p+q-2}(n)$。
\end{enumerate}
\end{lemma}

\begin{proof}
由 5.10 和滤过 $(D_p(n);\ p\in\mathbb{Z})$ 的定义，只需检验
$$[\partial^I,X^J]\in D_{|I|+|J|-2}(n).$$
对 $|I|$ 归纳证明。若 $|I|=1$，则对某 $1\le i\le n$ 有 $\partial^I=\partial_i$，且 $[\partial_i,X^J]=\partial_i(X^J)\in D_{|J|-1}(n)$。若 $|I|>1$，可写 $\partial^I=\partial^{I'}\partial_i$，其中 $I'\in\mathbb{Z}_+^n$，$1\le i\le n$。这导出
\begin{align*}
[\partial^I,X^J]&=[\partial^{I'}\partial_i,X^J]=\partial^{I'}\partial_i X^J-X^J\partial^{I'}\partial_i\\
&=\partial^{I'}[\partial_i,X^J]+[\partial^{I'},X^J]\partial_i\\
&=[\partial^{I'},[\partial_i,X^J]]+[\partial_i,X^J]\partial^{I'}+[\partial^{I'},X^J]\partial_i,
\end{align*}
故由归纳假设 $[\partial^I,X^J]\in D_{|I|+|J|-2}(n)$。
\end{proof}

这意味着 $(D_p(n);\ p\in\mathbb{Z})$ 是与 $D(n)$ 上环结构相容的滤过。此外，分次环 $\Gr D(n)$ 是交换 $k$-代数。若定义线性映射 $\Psi_p$（从 $D_p(n)$ 到 $k[X_1,\dots,X_n,\xi_1,\dots,\xi_n]$）为
$$\Psi_p\Bigl(\sum_{|I|+|J|\le p}a_{IJ}X^I\partial^J\Bigr)=\sum_{|I|+|J|=p}a_{IJ}X^I\xi^J,$$
则它是 $\Gr_p D(n)$ 到 $p$ 次齐次多项式的线性同构。因此它扩充为线性同构
$$\Psi:\Gr D(n)\longrightarrow k[X_1,X_2,\dots,X_n,\xi_1,\xi_2,\dots,\xi_n].$$
由 5.14 可见此映射是 $k$-代数同构。因此，配置滤过 $(D_p(n);\ p\in\mathbb{Z})$ 的环 $D(n)$ 满足 §3 的性质 (i)--(vii)。滤过 $(D_p(n);\ p\in\mathbb{Z})$ 称为 $D(n)$ 的\emph{Bernstein 滤过}。

显然，$D(n)$ 的主反自同构和 Fourier 自同构保持 Bernstein 滤过。
\begin{annote}
文中的 $\mathcal D(n)=k[x_1,\dots,x_n]\langle\partial_1,\dots,\partial_n\rangle$ 即 $n$ 阶 Weyl 代数 $A_n(k)$，也是 $A^n$ 上的微分算子环。一般簇 $X$ 上 $\mathcal D_X$ 由 $\mathcal O_X$ 与 $\mathcal T_X$ 生成、满足 Leibniz 关系，其滤过由“微分阶”给出；光滑点处 $\mathrm{gr}\mathcal D_X\cong\mathrm{Sym}\,\mathcal T_X$，所以 §7 的特征簇天然落在余切丛 $T^*X$ 里。正特征下 $\mathcal D_X$ 不再诺特，$\mathcal O_X$-凝聚也不等于 $\mathcal D_X$-凝聚，故全书默认 $k$ 特征零。[来源：nLab 条目 D-module（Last revised 2024-07-27）；原书 §5, p23–32]
\end{annote}



\section{多项式系数微分算子环上的模}

本节研究多项式系数微分算子环 $D(n)$ 上的模范畴。记 $\mathcal{M}_L(D(n))$ 和 $\mathcal{M}_R(D(n))$ 分别为左、右 $D(n)$-模的范畴。它们是 Abel 范畴。$D(n)$ 的主反自同构 $\varphi$ 定义一个从右模范畴 $\mathcal{M}_R(D(n))$ 到左模范畴 $\mathcal{M}_L(D(n))$ 的正合函子，它把模 $M$ 映为其\emph{转置} $M^t$——作为加法群等于 $M$，而 $D(n)$ 的作用由 $(T,m)\longmapsto\varphi(T)m$（$T\in D(n)$，$m\in M$）给出。类似地定义从 $\mathcal{M}_L(D(n))$ 到 $\mathcal{M}_R(D(n))$ 的函子。显然这些函子互为范畴的同构逆。若记 $\mathcal{M}_{fg}^L(D(n))$ 和 $\mathcal{M}_{fg}^R(D(n))$ 为相应的有限生成模全子范畴，则这些函子也诱导它们的等价。因此在以下讨论中可仅限于左模，并在记号中省略上标 $L$（除非想强调处理的是右模）。由于 $D(n)$ 是诺特环，全子范畴 $\mathcal{M}_{fg}(D(n))$ 在取子模、商模和扩张下封闭。

首先把 $D(n)$ 看作配置 Bernstein 滤过的环。由于此时 $D_0(n)=k$，可以用 $k$ 上有限维向量空间范畴上的可加函数 $\dim_k$ 定义 $\mathcal{M}_{fg}^L(D(n))$ 和 $\mathcal{M}_{fg}^R(D(n))$ 中模的维数。模 $M$ 的这个维数 $d(M)$ 及相应重数 $e(M)$ 分别称为 \emph{Bernstein 维数}和\emph{Bernstein 重数}。由于主反自同构保持 Bernstein 滤过，对任意有限生成 $D(n)$-模 $M$ 有 $d(M)=d(M^t)$。

对任意有限生成 $D(n)$-模 $M$ 有正合列 $D(n)^p\longrightarrow M\longrightarrow 0$，故 $d(M)\le d(D(n))$。此外，由 5.6 得出如下结果。

\begin{lemma}
对任意有限生成 $D(n)$-模 $M$ 有 $d(M)\le 2n$。
\end{lemma}

\begin{example}
考虑单变量的多项式微分算子代数 $D(1)$。设 $M$ 是非零有限生成 $D(1)$-模。则其 Bernstein 维数 $d(M)$ 可为 $0$、$1$ 或 $2$。显然 $d(M)=0$ 会蕴含：对 $M$ 的任意好滤过 $F M$，函数 $p\longmapsto\dim F_p M$ 对充分大的 $p\in\mathbb{Z}$ 为常数。由于 $F M$ 是穷尽的，这意味着 $M$ 有限维。记 $\pi(x)$ 和 $\pi(\partial)$ 为由 $x$ 和 $\partial$ 的作用诱导的 $M$ 上的线性变换。则 $[\pi(x),\pi(\partial)]=1_M$。对等式两边取迹可得 $\dim_k M=0$，与 $M\neq 0$ 的假设矛盾。故 $d(M)$ 只能是 $1$ 或 $2$。
\end{example}

$D(n)$ 维数理论的主要结果是把上述例子推广的下述陈述。

\begin{theorem}[Bernstein]
设 $M$ 是有限生成 $D(n)$-模且 $M\neq 0$。则 $d(M)\ge n$。
\end{theorem}

\begin{proof}
由于 $M$ 是有限生成 $D(n)$-模，由 3.4 可给它配置好滤过。此外，通过平移指标，显然可设 $n<0$ 时 $F_n M=0$ 且 $F_0 M\neq 0$。

对任意 $p\in\mathbb{Z}_+$，考虑线性映射 $D_p(n)\longrightarrow\Hom_k(F_p M,F_{2p}M)$，它把 $T\in D_p(n)$ 映为线性映射 $m\longmapsto Tm$。我们断言此映射是单射。$p\le 0$ 时显然。设它对 $p-1$ 成立，且 $T\in D_p(n)$ 满足对所有 $m\in F_p M$ 有 $Tm=0$。则对任意 $v\in F_{p-1}M$ 和 $1\le i\le n$，有 $X_i v\in F_p M$ 和 $\partial_i v\in F_p M$，故
$$[X_i,T]v=X_iTv-TX_iv=0$$
以及
$$[\partial_i,T]v=\partial_i T v-T\partial_i v=0,$$
且由 5.14 有 $[X_i,T],[\partial_i,T]\in D_{p-1}(n)$。由归纳假设这意味着对 $1\le i\le n$ 有 $[X_i,T]=0$、$[\partial_i,T]=0$，即 $T$ 在 $D(n)$ 的中心中。由于 $D(n)$ 的中心由 5.12 等于 $k$，我们断定 $T=0$。因此对任意 $p\in\mathbb{Z}$，
$$\dim_k\bigl(D_p(n)\bigr)\le\dim_k\bigl(\Hom_k(F_p M,F_{2p}M)\bigr)=\dim_k(F_p M)\cdot\dim_k(F_{2p}M).$$
另一方面，对充分大的 $p\in\mathbb{Z}_+$，左边等于 $p$ 的 $2n$ 次多项式（首项系数为正），右边等于 $p$ 的 $2d(M)$ 次多项式（首项系数为正）。这只有在 $d(M)\ge n$ 时才可能。
\end{proof}

下一节将给出 Bernstein 维数的几何解释。

最后，若 $M$ 是 $D(n)$-模，可定义其 Fourier 变换 $F(M)$ 为这样的模：作为加法群等于 $M$，$D(n)$ 的作用由 $(T,m)\longmapsto F(T)m$（$T\in D(n)$，$m\in M$）给出。显然 Fourier 变换是范畴 $\mathcal{M}(D(n))$ 的自同构。它也诱导 $\mathcal{M}_{fg}(D(n))$ 的自同构。由 Fourier 自同构 $F$ 保持 Bernstein 滤过（或 3.9）这一事实得出如下结果。

\begin{lemma}
设 $M$ 是有限生成 $D(n)$-模。则 $d(F(M))=d(M)$。
\end{lemma}

\section{特征簇}

现在研究有限生成 $D(n)$-模的一个更具几何味道的不变量。特别地，它将用 $D(n)$ 按微分算子次数的滤过 $F D(n)$ 来构造，而非 Bernstein 滤过。与 Bernstein 滤过不同，次数滤过对任意光滑仿射代数簇上的微分算子环都有意义。

首先，由于任何 $D(n)$-模 $M$ 都可看作 $k[X_1,X_2,\dots,X_n]$-模，可考虑其支集 $\operatorname{supp}(M)\subset k^n$。

\begin{proposition}
设 $M$ 是有限生成 $D(n)$-模。则 $\operatorname{supp}(M)$ 是 $k^n$ 的闭子簇。
\end{proposition}

\begin{proof}
固定 $M$ 上的好滤过 $F M$。则对 $x\in k^n$，$M_x=0$ 等价于对所有 $p\in\mathbb{Z}$ 有 $(F_p M)_x=0$。因此由局部化的正合性，等价于 $(\Gr M)_x=0$。设 $I_p$ 是 $k[X_1,X_2,\dots,X_n]$-模 $\Gr_p M$（$p\in\mathbb{Z}$）的零化子。由于 $\Gr_p M$ 是有限生成 $k[X_1,X_2,\dots,X_n]$-模，由 4.2 其支集 $\operatorname{supp}(\Gr_p M)=V(I_p)$。这意味着 $\operatorname{supp}(M)=\bigcup_{p\in\mathbb{Z}}V(I_p)$。设 $m_1,m_2,\dots,m_s$ 是 $\Gr D(n)$-模 $\Gr M$ 的一组齐次生成元。则 $m_1,m_2,\dots,m_s$ 在 $k[X_1,X_2,\dots,X_n]$ 中的零化子 $I$ 零化整个 $\Gr M$。因此存在 $\mathbb{Z}$ 的有限子集 $S$ 使得 $\bigcap_{p\in S}I_p=I\subset I_q$ 对所有 $q\in\mathbb{Z}$ 成立。这意味着 $\bigcup_{p\in S}V(I_p)=V(I)\supset V(I_q)$ 对所有 $q\in\mathbb{Z}$ 成立，且 $\operatorname{supp}(M)=V(I)$。
\end{proof}

设 $D$ 是配置满足 §3 开头性质 (i)--(vii) 的滤过 $F D$ 的滤过环。设 $M$ 是有限生成 $D$-模，$F M$ 是 $M$ 的好滤过。则 $\Gr M$ 是分次 $\Gr D$-模。设 $I$ 是 $\Gr M$ 在 $\Gr D$ 中的零化子。这显然是 $\Gr D$ 中的分次理想，故其根 $r(I)$ 也是分次理想。一般地 $I$ 依赖于 $M$ 上好滤过的选择，但也有如下结果。

\begin{lemma}
设 $M$ 是有限生成 $D$-模，$F M$ 和 $F'M$ 是 $M$ 上两个好滤过。设 $I$、$I'$ 分别是相应分次 $\Gr D$-模 $\Gr M$、$\Gr' M$ 的零化子。则 $r(I)=r(I')$。
\end{lemma}

\begin{proof}
设 $T\in r(I)\cap\Gr_p D$。则存在 $s\in\mathbb{Z}_+$ 使得 $T^s\in I$。若取 $Y\in F_p D$ 使得 $Y+F_{p-1}D=T$，则对所有 $q\in\mathbb{Z}$ 有 $Y^s F_q M\subset F_{q+sp-1}M$。故由归纳得
$$Y^{ms}F_q M\subset F_{q+msp-m}M$$
对所有 $m\in\mathbb{N}$ 和 $q\in\mathbb{Z}$ 成立。另一方面，由 3.7 知 $F M$ 与 $F'M$ 等价。故存在 $l\in\mathbb{Z}_+$ 使得对所有 $q\in\mathbb{Z}$ 有 $F_q M\subset F'_{q+l}M\subset F_{q+2l}M$。这导出
$$Y^{ms}F'_q M\subset Y^{ms}F_{q+l}M\subset F_{q+l+msp-m}M\subset F'_{q+2l+msp-m}M$$
对所有 $q\in\mathbb{Z}$ 和 $m\in\mathbb{N}$ 成立。若取 $m>2l$，则对所有 $q\in\mathbb{Z}$ 有 $Y^{ms}F'_q M\subset F'_{q+msp-1}M$，即 $T^{ms}\in I'$。因此 $T\in r(I')$，有 $r(I)\subset r(I')$。由于 $I$ 和 $I'$ 的角色对称，断定 $r(I)=r(I')$。
\end{proof}

因此 $\Gr M$ 零化子的根与 $\Gr M$ 上好滤过的选择无关。我们称之为 $M$ 的\emph{特征理想}，记为 $J(M)$。

现在可把此构造应用于 $D(n)$。由于由 5.6 有 $\Gr D(n)=k[X_1,\dots,X_n,\xi_1,\dots,\xi_n]$，可定义闭代数集
$$\operatorname{Ch}(M)=V(J(M))\subset k^{2n},$$
称之为 $M$ 的\emph{特征簇}。

由于 $J(M)$ 关于后 $n$ 个变量是齐次理想，立即得到如下结果。

\begin{lemma}
有限生成 $D(n)$-模 $M$ 的特征簇 $\operatorname{Ch}(M)$ 具有如下性质：若 $(x,\xi)\in\operatorname{Ch}(M)$，则对任意 $\lambda\in k$ 有 $(x,\lambda\xi)\in\operatorname{Ch}(M)$。
\end{lemma}

我们称 $\operatorname{Ch}(M)$ 是\emph{锥簇}。

\begin{proposition}
设
$$0\longrightarrow M'\longrightarrow M\longrightarrow M''\longrightarrow 0$$
是有限生成 $D(n)$-模的正合列。则
$$\operatorname{Ch}(M)=\operatorname{Ch}(M')\cup\operatorname{Ch}(M'').$$
\end{proposition}

\begin{proof}
设 $F M$ 是 $M$ 上的好滤过。则它在 $M'$ 上诱导滤过 $F M'$，在 $M''$ 上诱导滤过 $F M''$。由 3.8 知这些滤过也是好的。此外有有限生成 $k[X_1,\dots,X_n,\xi_1,\dots,\xi_n]$-模的正合列
$$0\longrightarrow\Gr M'\longrightarrow\Gr M\longrightarrow\Gr M''\longrightarrow 0,$$
由 4.2 它们的支集分别是 $D(n)$-模 $M$、$M'$、$M''$ 的特征簇。因此结论由 4.1 得出。
\end{proof}

接下来两个结果阐明特征簇与有限生成 $D(n)$-模支集之间的关系。设 $\pi:k^{2n}\longrightarrow k^n$ 是由 $\pi(x,\xi)=x$（对 $x,\xi\in k^n$）定义的映射。

\begin{proposition}
设 $M$ 是有限生成 $D(n)$-模。则 $\operatorname{supp}(M)=\pi(\operatorname{Ch}(M))$。
\end{proposition}

\begin{proof}
记 $m_1,m_2,\dots,m_s$ 为 $\Gr M$ 的一组齐次生成元。则如 7.1 的证明，$m_1,m_2,\dots,m_s$ 在 $k[X_1,X_2,\dots,X_n]$ 中的零化子 $I$ 满足 $\operatorname{supp}(M)=V(I)$。另一方面，若 $J$ 是 $m_1,m_2,\dots,m_s$ 在 $k[X_1,X_2,\dots,X_n,\xi_1,\xi_2,\dots,\xi_n]$ 中的零化子，它是关于 $\xi_1,\xi_2,\dots,\xi_n$ 的齐次理想，满足 $I=k[X_1,X_2,\dots,X_n]\cap J$，且 $\operatorname{Ch}(M)=V(J)$。这意味着 $x\in V(I)=\operatorname{supp}(M)$ 等价于 $(x,0)\in V(J)=\operatorname{Ch}(M)$。由于 $\operatorname{Ch}(M)$ 是锥的，这蕴含结论。
\end{proof}

设 $M$ 是有限生成 $D(n)$-模。定义 $M$ 的\emph{奇异支集}为
$$\operatorname{sing\,supp}(M)=\{x\in k^n\mid (x,\xi)\in\operatorname{Ch}(M)\ \text{对某个 }\xi\neq 0\}.$$
显然有 $\operatorname{sing\,supp}(M)\subset\operatorname{supp}(M)$。

\begin{lemma}
设 $M$ 是有限生成 $D(n)$-模。则 $\operatorname{sing\,supp}(M)$ 是 $\operatorname{supp}(M)$ 的闭子簇。
\end{lemma}

\begin{proof}
设 $p:k^n-\{0\}\longrightarrow P^{n-1}(k)$ 为自然投射。则
$$1\times p:k^n\times(k^n-\{0\})\longrightarrow k^n\times P^{n-1}(k)$$
把 $\operatorname{Ch}(M)-(k^n\times\{0\})$ 投射到 $k^n\times P^{n-1}(k)$ 中对应于理想 $J(M)$（它关于 $\xi_1,\xi_2,\dots,\xi_n$ 齐次）的闭子簇上。最后，到第一个因子的投射 $k^n\times P^{n-1}(k)\longrightarrow k^n$ 把它映到 $\operatorname{sing\,supp}(M)$ 上。由于 $P^{n-1}(k)$ 是完全簇，投射 $k^n\times P^{n-1}(k)\longrightarrow k^n$ 是闭映射。因此 $\operatorname{sing\,supp}(M)$ 是闭的。
\end{proof}

关于特征簇的基本结果如下定理，它也给出 Bernstein 维数的几何描述。

\begin{theorem}
设 $M$ 是有限生成 $D(n)$-模。则
$$\dim\operatorname{Ch}(M)=d(M).$$
\end{theorem}

证明此定理需要一些准备。先考虑 $D(n)$-模 $M=D(n)/L$，其中 $L$ 是 $D(n)$ 中的左理想。$D(n)$-模的正合列
$$0\longrightarrow L\longrightarrow D(n)\longrightarrow M\longrightarrow 0$$
（对 $D(n)$ 配置由微分算子次数滤过诱导的滤过）导出 $A$-模的正合列
$$0\longrightarrow\Gr L\longrightarrow\Gr D(n)\longrightarrow\Gr M\longrightarrow 0,$$
其中 $A=\Gr D(n)=k[X_1,\dots,X_n,\xi_1,\dots,\xi_n]$。因此 $\Gr M$ 是商 $A/\Gr L$，$\Gr M$ 的零化子等于 $\Gr L$。故由定义，$V(\Gr L)$ 是 $M$ 的特征簇。因此在这种情况下须建立 $\dim V(\Gr L)=d(D(n)/L)$。

为证明这一点，先作一些一般性注记。设 $A=k[X_1,X_2,\dots,X_n]$。设 $\mathbf{t}=(t_1,\dots,t_n)\in\mathbb{N}^n$。通过令 $\Gr^{\mathbf{t}}_m A$ 为满足 $\sum_{j=1}^n t_j i_j=m$ 的 $X^I$ 的线性张成，定义 $\Gr^{\mathbf{t}} A$ 的分次。显然这样 $A$ 成为分次环。此外可定义相应滤过 $F^{\mathbf{t}}A$ 为 $F^{\mathbf{t}}_p A=\sum_{m\le p}\Gr^{\mathbf{t}}_m A$。显然，若以 $F A$ 记 $A$ 按多项式次数的自然滤过，并令 $t=\max_{1\le i\le n}t_i$，则对任意 $p\in\mathbb{Z}$ 有
$$F^{\mathbf{t}}_p A\subset F_p A,\qquad F_p A\subset F^{\mathbf{t}}_{tp}A.$$

设 $I$ 是 $A$ 中的理想。则可考虑配置由 $A$ 上滤过诱导的滤过的 $A$-模正合列
$$0\longrightarrow I\longrightarrow A\longrightarrow A/I\longrightarrow 0.$$
则对任意 $p\in\mathbb{Z}$ 有
$$F^{\mathbf{t}}_p(A/I)\subset F_p(A/I),\qquad F_p(A/I)\subset F^{\mathbf{t}}_{tp}(A/I).$$
这进而蕴含如下引理。

\begin{lemma}
对任意 $p\in\mathbb{Z}$，有
$$\dim_k F^{\mathbf{t}}_p(A/I)\le\dim_k F_p(A/I),\qquad \dim_k F_p(A/I)\le\dim_k F^{\mathbf{t}}_{tp}(A/I).$$
\end{lemma}

设 $s\in\mathbb{N}$。则定义代数 $D(n)$ 的滤过 $F^{(s)}D(n)$ 为
$$F^{(s)}_m D(n)=\Bigl\{T\in D(n)\ \Big|\ T=\sum_{|I|+s|J|\le m}c_{I,J}X^I\partial^J,\ c_{I,J}\in k\Bigr\}.$$
显然 $F^{(1)}D(n)$ 是 $D(n)$ 的 Bernstein 滤过。滤过 $F^{(s)}D(n)$ 具有 §3 中考虑的环滤过的性质 (i)--(iii)。此外，$T\in F^{(s)}_m D(n)$ 当且仅当 $T\in F^{(1)}_p D(n)$ 且 $T$ 的阶数 $\le q$，其中 $p,q$ 满足 $m=p+(s-1)q$。因此，若 $T\in F^{(s)}_m D(n)$、$S\in F^{(s)}_{m'}D(n)$，则存在 $p,p'$ 和 $q,q'$ 使得 $m=p+(s-1)q$、$m'=p'+(s-1)q'$，$T\in F^{(1)}_p D(n)$，$S\in F^{(1)}_{p'}D(n)$，且 $T$、$S$ 的阶数分别 $\le q$、$\le q'$。这意味着 $TS$ 的阶数 $\le q+q'$ 且 $TS\in F^{(1)}_{p+p'}D(n)$。由此可知 $TS\in F^{(s)}_{m+m'}D(n)$。故滤过 $F^{(s)}D(n)$ 也满足 (iv)，即是环滤过。同样可检验 (v) 成立，即分次环 $\Gr^{(s)}D(n)$ 是交换的。此外，分次环 $\Gr^{(s)}D(n)$ 同构于分次环 $A=k[X_1,\dots,X_n,\xi_1,\dots,\xi_n]$，其分次结构对应于 $\mathbf{t}=(1,\dots,1,s,\dots,s)$。我们记该分次模为 $\Gr^{(s)}A$，其相应滤过为 $F^{(s)}A$。

此外，对任意 $p\in\mathbb{Z}$ 有
$$F^{(s)}_p D(n)\subset F^{(1)}_p D(n),\qquad F^{(1)}_p D(n)\subset F^{(s)}_{sp}D(n).$$

再考虑配置由 $D(n)$ 上滤过诱导的滤过的 $D(n)$-模正合列
$$0\longrightarrow L\longrightarrow D(n)\longrightarrow D(n)/L\longrightarrow 0.$$
则对任意 $p\in\mathbb{Z}$ 有
$$F^{(s)}_p(D(n)/L)\subset F^{(1)}_p(D(n)/L),\qquad F^{(1)}_p(D(n)/L)\subset F^{(s)}_{sp}(D(n)/L).$$
这进而蕴含与 7.8 类似的如下引理。

\begin{lemma}
对任意 $p\in\mathbb{Z}$，有
$$\dim_k F^{(s)}_p(D(n)/L)\le\dim_k F^{(1)}_p(D(n)/L)$$
以及
$$\dim_k F^{(1)}_p(D(n)/L)\le\dim_k F^{(s)}_{sp}(D(n)/L).$$
\end{lemma}

\begin{lemma}
设 $L$ 是 $D(n)$ 中的左理想。则对任意 $s\in\mathbb{N}$，
$$d(D(n)/L)=\dim V(\Gr^{(s)}L).$$
\end{lemma}

\begin{proof}
正合列
$$0\longrightarrow L\longrightarrow D(n)\longrightarrow D(n)/L\longrightarrow 0$$
（其中 $D(n)$ 配置滤过 $F^{(s)}D(n)$，$L$ 和 $D(n)/L$ 分别配置诱导滤过 $F^{(s)}L$ 和 $F^{(s)}(D(n)/L)$）导出分次 $A$-模（$A=\Gr^{(s)}D(n)$）的正合列
$$0\longrightarrow\Gr^{(s)}L\longrightarrow\Gr^{(s)}D(n)\longrightarrow\Gr^{(s)}(D(n)/L)\longrightarrow 0.$$
显然作为分次环 $\Gr^{(s)}D(n)=\Gr^{(s)}A$。这意味着对任意 $p\in\mathbb{Z}$，
\begin{align*}
\dim_k F^{(s)}_p(D(n)/L)&=\sum_{q=0}^{p}\bigl(\dim_k F^{(s)}_q(D(n)/L)-\dim_k F^{(s)}_{q-1}(D(n)/L)\bigr)\\
&=\sum_{q=0}^{p}\dim_k\Gr^{(s)}_q(D(n)/L)=\sum_{q=0}^{p}\bigl(\dim_k\Gr^{(s)}_q D(n)-\dim_k\Gr^{(s)}_q L\bigr)\\
&=\sum_{q=0}^{p}\bigl(\dim_k\Gr^{(s)}_q A-\dim_k\Gr^{(s)}_q L\bigr)=\sum_{q=0}^{p}\dim_k\Gr^{(s)}_q(A/\Gr^{(s)}L)\\
&=\dim_k F^{(s)}_p(A/\Gr^{(s)}L).
\end{align*}
这进而利用 7.8 和 7.9 蕴含
$$\dim_k F^{(1)}_p(D(n)/L)\le\dim_k F^{(s)}_{sp}(D(n)/L)=\dim_k F^{(s)}_{sp}(A/\Gr^{(s)}L)\le\dim_k F_{sp}(A/\Gr^{(s)}L)$$
以及
$$\dim_k F_p(A/\Gr^{(s)}L)\le\dim_k F^{(s)}_{sp}(A/\Gr^{(s)}L)=\dim_k F^{(s)}_{sp}(D(n)/L)\le\dim_k F^{(1)}_{sp}(D(n)/L).$$
由于函数 $p\longmapsto\dim_k F^{(1)}_p(D(n)/L)$ 和 $p\longmapsto\dim_k F_p(A/\Gr^{(s)}L)$ 对充分大的 $p\in\mathbb{Z}$ 由多项式表示，这些多项式必有相等的次数。这进而蕴含 $d(D(n)/L)=d(A/\Gr^{(s)}L)$。由于 $\Gr^{(s)}L$ 是 $A/\Gr^{(s)}L$ 的零化子，结论由 4.5 得出。
\end{proof}

对任意 $s\in\mathbb{N}$，记 $\sigma^{(s)}_p(T)$ 为 $T\in F^{(s)}_p D(n)$ 在 $\Gr^{(s)}_p D(n)=A$ 中的投影。此外，对 $A$ 上由多项式次数给出的自然滤过，记 $\sigma_p$ 为把 $p$ 次多项式映为其 $p$ 次齐次分量的映射。

\begin{example}
设 $D=D(1)$，$T\in D$ 由 $T=x^3\partial+\partial^2$ 给出。则 $T$ 的次数等于 $2$，且 $\Symb_2(T)=\xi^2$。故 $\sigma_2(\Symb_2(T))=\xi^2$。

另一方面，有 $\sigma^{(1)}_4(T)=x^3\xi$，$\sigma^{(2)}_5(T)=x^3\xi$，$\sigma^{(3)}_6(T)=x^3\xi+\xi^2$，且当 $s>3$ 时 $\sigma^{(s)}_{2s}(T)=\xi^2$。
\end{example}

因此，对充分大的 $s$，$\sigma^{(s)}(T)$ 变得等于 $\sigma(\Symb(T))$。一般情形也成立，更精确地说有如下结果。

\begin{lemma}
设 $T$ 是 $D(n)$ 中阶数 $\le m$ 的微分算子，其符号 $\Symb_m(T)$ 是 $p$ 次多项式。则存在 $s_0$ 使得对 $s\ge s_0$，
$$\sigma^{(s)}_{p+(s-1)m}(T)=\sigma_p(\Symb_m(T)).$$
\end{lemma}

\begin{proof}
由假设
$$T=\sum_{|J|\le m}c_{I,J}X^I\partial^J.$$
此外可固定 $q_0$ 使得 $c_{I,J}\neq 0$ 蕴含 $|I|\le q_0$。则
$$\Symb_m(T)=\sum_{|J|=m}c_{I,J}X^I\xi^J$$
是 $p$ 次多项式，其首项为
$$\sigma_p(\Symb_m(T))=\sum_{|I|=p-m,\,|J|=m}c_{I,J}X^I\xi^J.$$
另一方面，项 $X^I\partial^J$ 属于 $F^{(s)}_{|I|+s|J|}D(n)$。设 $c_{I,J}\neq 0$。则有以下可能：

(i) $|J|=m$，$|I|=p-m$：$X^I\partial^J$ 属于 $F^{(s)}_{p+(s-1)m}D(n)$。

(ii) $|J|=m$，$|I|<p-m$：$X^I\partial^J$ 属于 $F^{(s)}_{p+(s-1)m-1}D(n)$。

(iii) $m\ge 1$，$|J|<m$，$|I|\le q_0$：$X^I\partial^J$ 属于 $F^{(s)}_{q_0+s(m-1)}D(n)$。此外
$$q_0+s(m-1)=q_0+sm-s=q_0+m-s+(s-1)m.$$
故若 $s\ge s_0=q_0+m-p+1$，则 $q_0+s(m-1)\le p+(s-1)m-1$。由此可知此时微分算子 $X^I\partial^J$ 也属于 $F^{(s)}_{p+(s-1)m-1}D(n)$。

这意味着对 $s\ge s_0$ 有
\begin{align*}
\sigma^{(s)}_{p+(s-1)m}(T)&=\sigma^{(s)}_{p+(s-1)m}\Bigl(\sum_{|I|=p-m,\,|J|=m}c_{I,J}X^I\partial^J\Bigr)\\
&=\sum_{|I|=p-m,\,|J|=m}c_{I,J}X^I\xi^J=\sigma_p(\Symb_m(T)).
\end{align*}
\end{proof}

特别地，可取有限多个 $D_i\in L$（$1\le i\le k$），使得 $\sigma(\Symb(D_i))$ 生成 $\Gr(\Gr L)$。由 7.12，存在 $s$ 使得对任意 $1\le i\le k$ 有 $\sigma(\Symb(D_i))=\sigma^{(s)}(D_i)$。因此有 $\Gr(\Gr L)\subset\Gr^{(s)}L$ 和 $V(\Gr(\Gr L))\supset V(\Gr^{(s)}L)$。这进而蕴含 $\dim V(\Gr(\Gr L))\ge\dim V(\Gr^{(s)}L)$。由 4.8 和 7.10，最终得 $\dim V(\Gr L)\ge d(D(n)/L)$。

反向不等式简单得多。首先设 $D\in F^{(1)}_p(D(n))$。则其次数 $\le p$，且 $\Symb(D)$ 是次数 $\le p$ 的多项式。因此 $F^{(1)}_p(L)$ 只含次数 $\le p$ 的微分算子，且 $\Gr F^{(1)}_p(L)\subset F_p(\Gr L)$。由于 $F^{(1)}_p(L)$ 有限维，且其按微分算子次数的滤过是 Hausdorff 且穷尽的，有
$$\dim_k F^{(1)}_p(L)=\dim_k\Gr F^{(1)}_p(L)\le\dim_k F_p(\Gr L).$$
因此，对任意 $p\in\mathbb{N}$，
\begin{align*}
\dim_k F^{(1)}_p(D(n)/L)&=\dim_k D_p(n)-\dim_k F^{(1)}_p(L)\\
&\ge\dim_k F_p A-\dim_k F_p(\Gr L)=\dim_k F_p(A/\Gr L).
\end{align*}
这进而蕴含 $d(D(n)/L)\ge d(A/\Gr L)$。由于 $\Gr L$ 是 $A/\Gr L$ 的零化子，由 4.5 得 $d(D(n)/L)\ge\dim V(\Gr L)$。

故证明了
$$d(D(n)/L)=\dim V(\Gr L).$$
为证明一般情形的定理，考虑正合列
$$0\longrightarrow M'\longrightarrow M\longrightarrow M''\longrightarrow 0,$$
其中 $M$ 有 $q$ 个生成元，$M'$ 有 $q-1$ 个生成元，$M''$ 是循环模。因此 $M'$ 同构于某左理想 $L$ 的 $D(n)/L$。

由证明的第一部分有 $d(M'')=\dim\operatorname{Ch}(M'')$。此外由归纳假设有 $d(M')=\dim\operatorname{Ch}(M')$。由 3.9 和 7.4 得
\begin{align*}
d(M)&=\max(d(M'),d(M''))=\max(\dim\operatorname{Ch}(M'),\dim\operatorname{Ch}(M''))\\
&=\dim(\operatorname{Ch}(M')\cup\operatorname{Ch}(M''))=\dim\operatorname{Ch}(M).
\end{align*}
这完成 7.7 的证明。

特别地，结合 6.3 和 7.7 得到如下结果。

\begin{theorem}
设 $M$ 是有限生成 $D(n)$-模，$M\neq 0$，$\operatorname{Ch}(M)$ 是其特征簇。则 $\dim\operatorname{Ch}(M)\ge n$。
\end{theorem}
\begin{annote}
$\mathrm{Ch}(M)$ 定义为 $k[x,\xi]$ 中 $\mathrm{gr}M$ 的零化子所定义的闭集，因此它是 $T^*A^n$ 中的锥：$\xi$ 方向齐次。这一点不依赖好滤过的选取，因为两条好滤过的 $\sqrt{\mathrm{Ann}}$ 相同。nLab 的 characteristic variety 条目只列文献、无陈述，故本条不作外部断言。[来源：nLab 条目 characteristic+variety（Last revised 2017-08-29，仅文献列表）；原书 §7, p32–38]
\end{annote}



\section{完整模}

设 $M$ 是非平凡有限生成 $D(n)$-模。则由 7.13，其特征簇 $\operatorname{Ch}(M)$ 的维数 $\ge n$。

若有限生成 $D(n)$-模的特征簇 $\operatorname{Ch}(M)$ 的维数 $\le n$，则称它是\emph{完整（holonomic）的}。因此 $M$ 是完整的，当且仅当 $M=0$ 或 $\dim\operatorname{Ch}(M)=n$。

粗略地说，完整模是具有最小可能特征簇的模。

以下是关于完整模的基本观察。

\begin{theorem}\ 
\begin{enumerate}[label=(\roman*)]
  \item 完整模具有有限长度；
  \item 完整模的子模、商模和扩张都是完整的。
\end{enumerate}
\end{theorem}

\begin{proof}
(ii) 由 3.9 立即得出。

(i) 设 $M$ 是非零完整 $D(n)$-模。则由定义其特征簇 $\operatorname{Ch}(M)$ 的维数等于 $n$。由 7.7，其 Bernstein 维数 $d(M)$ 也等于 $n$。由于 $M$ 有限生成且 $D(n)$ 是诺特环，存在 $M$ 中不同于 $M$ 的极大 $D(n)$-子模 $M'$。因此有正合列
$$0\longrightarrow M'\longrightarrow M\longrightarrow M/M'\longrightarrow 0.$$
由 (ii)，$M'$ 和 $M/M'$ 是完整的，且 $M/M'$ 是不可约 $D(n)$-模。若 $M'\neq 0$，由 3.9 得 $e(M')<e(M)$。因此由 $e(M)$ 上的归纳，$M$ 具有有限长度。
\end{proof}

因此范畴 $\mathcal{M}_{fg}(D(n))$ 的全子范畴 $\operatorname{Hol}(D(n))$ 在取子模、商模和扩张下封闭。此外，若记 $\mathcal{M}_{fl}(D(n))$ 为 $\mathcal{M}_{fg}(D(n))$ 中由有限长度 $D(n)$-模构成的全子范畴，则 $\operatorname{Hol}(D(n))$ 是 $\mathcal{M}_{fl}(D(n))$ 的子范畴。可以证明当 $n>1$ 时 $\operatorname{Hol}(D(n))$ 严格小于 $\mathcal{M}_{fl}(D(n))$。

此外，转置函子和 Fourier 函子把完整模映为完整模。

现在讨论完整模的一些例子。

\begin{example}
设 $\mathcal{O}_n=k[X_1,X_2,\dots,X_n]$。则
$$\mathcal{O}_n=D(n)/\bigl(D(n)(\partial_1,\partial_2,\dots,\partial_n)\bigr)$$
是有限生成 $D(n)$-模。此外，若令 $p<0$ 时 $F_p\mathcal{O}_n=0$，$p\ge 0$ 时 $F_p\mathcal{O}_n=\mathcal{O}_n$，则滤过 $F\mathcal{O}_n$ 是关于 $D(n)$ 次数滤过的好滤过。相应的分次模 $\Gr\mathcal{O}_n$ 满足：$p\neq 0$ 时 $\Gr_p\mathcal{O}_n=0$，$\Gr_0\mathcal{O}_n=k[X_1,X_2,\dots,X_n]$。由此可知 $\Gr\mathcal{O}_n$ 的零化子等于 $k[X_1,X_2,\dots,X_n,\xi_1,\xi_2,\dots,\xi_n]$ 中由 $\xi_1,\xi_2,\dots,\xi_n$ 生成的理想。这意味着 $\operatorname{Ch}(\mathcal{O}_n)=k^n\times\{0\}\subset k^{2n}$。特别地 $\dim\operatorname{Ch}(\mathcal{O}_n)=n$，$\mathcal{O}_n$ 是完整的。此外 $\operatorname{supp}(\mathcal{O}_n)=k^n$，投射 $\pi:k^{2n}\longrightarrow k^n$ 是 $\operatorname{Ch}(\mathcal{O}_n)$ 到 $\mathcal{O}$ 上的双射。

通过求微分可见 $\mathcal{O}_n$ 的任何子模必含常数，因此 $\mathcal{O}_n$ 是不可约的。
\end{example}

\begin{example}
现考虑 $\Delta_n=F(\mathcal{O}_n)$。则有
$$\Delta_n=D(n)/\bigl(D(n)(X_1,X_2,\dots,X_n)\bigr).$$
显然 $\Delta_n$ 是完整的且不可约的。设 $\delta$ 是对应于 $1\in\mathcal{O}_1$ 的向量。则对任意 $1\le i\le n$ 有 $X_i\delta=0$。显然 $\Delta_n$ 由 $\delta^{(I)}=\partial^I\delta$（$I\in\mathbb{Z}_+^n$）张成。设 $F\Delta_n$ 是 $\Delta_n$ 的滤过：$p<0$ 时 $F_p\Delta_n=\{0\}$，$p\ge 0$ 时 $F_p\Delta_n$ 由 $\delta^{(I)}$（$|I|\le p$）张成。记 $\epsilon_i$ 为第 $i$ 个位置为 $1$ 的多重指标 $(0,\dots,0,1,0,\dots,0)$。则由 Fourier 变换的定义，
$$\partial_j\delta^{(I)}=\delta^{(I+\epsilon_j)},\qquad X_j\delta^{(I)}=-i_j\delta^{(I-\epsilon_j)}$$
对所有 $1\le j\le n$ 和 $I\in\mathbb{Z}_+^n$ 成立。这意味着 $F_p\Delta_n$ 是 $k[X_1,X_2,\dots,X_n]$-子模。此外对所有 $1\le i\le n$ 和 $p\in\mathbb{Z}$ 有 $\partial_i F_p\Delta_n\subset F_{p+1}\Delta_n$。故 $F\Delta_n$ 是关于按微分算子次数滤过的 $D(n)$ 的穷尽 $D(n)$-模滤过。设 $\bar\delta^{(I)}$ 为 $\delta^{(I)}$ 在 $\Gr_{|I|}\Delta_n$ 中代表的陪集。则 $\Gr_p\Delta_n$ 由 $\bar\delta^{(I)}$（$I\in\mathbb{Z}_+$，$|I|=p$）张成。显然 $X_i$ 在 $\Gr\Delta_n$ 上作为 $0$ 作用，而 $\xi_i$ 把 $\bar\delta^{(I)}$ 映为 $\bar\delta^{(I+\epsilon_i)}$。因此 $\bar\delta$ 作为 $k[X_1,X_2,\dots,X_n,\xi_1,\xi_2,\dots,\xi_n]$-模生成 $\Gr\Delta_n$，且 $F\Delta_n$ 是好滤过。此外 $\Gr\Delta_n$ 的零化子是由 $X_i$（$1\le i\le n$）生成的理想。故 $\Delta_n$ 的特征簇为 $\operatorname{Ch}(\Delta_n)=\{0\}\times k^n\subset k^{2n}$，支集 $\operatorname{supp}(\Delta_n)=\{0\}\subset k^n$。
\end{example}

现在构造更多完整模。从一个简单的完整性判别法开始。

\begin{lemma}
设 $D(n)$ 配置 Bernstein 滤过。设 $M$ 是 $D(n)$-模，$F M$ 是 $M$ 上穷尽的 $D(n)$-模滤过。若
$$\dim_k F_p M\le\frac{c^n}{n!}\,p+\text{（关于 $p$ 的低阶项）}$$
对所有 $p\in\mathbb{Z}_+$ 成立，则 $M$ 是完整 $D(n)$-模且其长度 $\le c$。特别地，$M$ 是有限生成 $D(n)$-模。
\end{lemma}

\begin{proof}
设 $N$ 是 $M$ 的有限生成 $D(n)$-子模。则 $F M$ 在 $N$ 上诱导穷尽 $D(n)$-模滤过。由 3.6 存在 $N$ 的好滤过 $F'N$ 和 $s\in\mathbb{Z}_+$ 使得对任意 $p\in\mathbb{Z}$ 有 $F'_p N\subset F_{p+s}N$。由此对 $p\in\mathbb{Z}_+$ 有
$$\dim_k F'_p N\le\dim_k F_{p+s}N\le\dim_k F_{p+s}M\le\frac{c^n}{n!}\,p+\text{（关于 $p$ 的低阶项）}.$$
因此 $d(N)\le n$，$N$ 是完整的。若 $N\neq 0$，则 $e(N)\le c$。显然这意味着 $N$ 的长度 $\le e(N)\le c$。由此可知 $M$ 的有限生成 $D(n)$-子模的任意递增序列稳定，且 $M$ 本身有限生成。
\end{proof}

\begin{example}
设 $n=1$，令 $D=D(1)$。对任意 $\alpha\in k$ 考虑 $D$-模 $M_\alpha=D/D(z\partial-\alpha)$。

令 $E=z\partial$。如 5.10 的证明，算子 $(z^pE^q,\ \partial^pE^q;\ p,q\in\mathbb{Z}_+)$ 构成 $D$ 作为 $k$-线性空间的一组基。此外，理想 $D(z\partial-\alpha)$ 由元素 $(z^pE^q(E-\alpha),\ \partial^pE^q(E-\alpha);\ p,q\in\mathbb{Z}_+)$ 张成。故 $M_\alpha$ 由对应于 $(z^p,\partial^p;\ p\in\mathbb{Z}_+)$ 的陪集张成。

显然
$$[E,z]=z\partial z-z^2\partial=z,\qquad [E,\partial]=z\partial^2-\partial z\partial=-\partial.$$
因此
$$Ez=z(E+1),\qquad E\partial=\partial(E-1).$$
这立即蕴含对任意 $n\in\mathbb{Z}_+$，$z^n$ 的陪集是 $E$ 的属于特征值 $\alpha+n$ 的特征向量。另一方面，对任意 $n\in\mathbb{Z}_+$，$\partial^n$ 的陪集是 $E$ 的属于特征值 $\alpha-n$ 的特征向量。因此 $E$ 在 $M_\alpha$ 上的谱等于 $\{\alpha+n;\ n\in\mathbb{Z}\}$，且每个特征值的重数等于 $1$。

$M_\alpha$ 的 Fourier 变换同构于
$$D/D(-\partial z-\alpha)=D/D(z\partial+\alpha+1)=M_{-\alpha-1}.$$

先设 $\alpha\notin\mathbb{Z}$。则 $E$ 是线性同构，$z$ 必为满射。由于 $z$ 把属于特征值 $\alpha+n$ 的特征空间映到属于特征值 $\alpha+n+1$ 的特征空间，$z$ 也是单射。因此可归纳构造一族向量 $z^{\alpha+n}$（$n\in\mathbb{Z}$），使得 $Ez^{\alpha+n}=(\alpha+n)z^{\alpha+n}$ 且 $zz^{\alpha+n}=z^{\alpha+n+1}$。显然这些向量构成 $M_\alpha$ 的一组基。此外对任意 $n\in\mathbb{Z}$，
$$\partial z^{\alpha+n}=\partial zz^{\alpha+n-1}=[\partial,z]z^{\alpha+n-1}+Ez^{\alpha+n-1}=(\alpha+n)z^{\alpha+n-1}.$$
这立即蕴含对任意整数 $p\in\mathbb{Z}$ 有 $M_\alpha\cong M_{\alpha+p}$。

此外，$M$ 的任何非零 $D$-子模在 $E$ 下不变，故含 $E$ 的一个特征向量。这进而蕴含它含某 $z^{\alpha+p}$（$p\in\mathbb{Z}$）。由此可知它含所有 $z^{\alpha+n}$（$n\in\mathbb{Z}$），即等于 $M_\alpha$。故 $M_\alpha$ 是不可约 $D$-模。

定义 $M_\alpha$ 的滤过 $F M_\alpha$：$p<0$ 时 $F_pM_\alpha=\{0\}$；$p\ge 0$ 时 $F_pM_\alpha$ 为 $\{z^{\alpha+n};\ |n|\le p\}$ 的张成。显然 $F M_\alpha$ 是 $M_\alpha$ 关于线性子空间的递增穷尽滤过。此外由上述注记，对任意 $p\in\mathbb{Z}$ 有 $zF_pM_\alpha\subset F_{p+1}M_\alpha$ 和 $\partial F_pM_\alpha\subset F_{p+1}M_\alpha$。因此 $F M_\alpha$ 是关于配置 Bernstein 滤过的 $D$ 的 $D$-模滤过。由于 $p\ge 0$ 时 $\dim_k F_pM_\alpha=2p+1$，由 8.4 知 $M_\alpha$ 是完整的。

为计算其特征簇，考虑另一个滤过 $F M_\alpha$：$n<0$ 时 $F_nM_\alpha=\{0\}$，$n\ge 0$ 时 $F_nM_\alpha$ 由 $\{z^{\alpha+p};\ p\ge -n\}$ 张成。显然这是 $F M_\alpha$ 关于 $z$ 中多项式环的模的穷尽滤过。此外对任意 $p\in\mathbb{Z}_+$ 有 $\partial F_pM_\alpha=F_{p+1}M_\alpha$，这是关于按微分算子次数滤过的 $D$ 的好 $D$-模滤过。分次模 $\Gr M_\alpha$ 是 $\Gr_pM_\alpha$ 的直和，其中 $n<0$ 时 $\Gr_nM_\alpha=0$；$\Gr_0M_\alpha$ 等于 $z^{\alpha+p}$（$p\ge 0$）的张成；$\Gr_pM_\alpha$ 由 $z^{\alpha-p}$ 模 $\{z^{\alpha+q};\ q>-p\}$ 的陪集张成。因此 $z$ 零化 $\Gr_pM_\alpha$（$p\neq 0$），而 $\partial$ 的符号 $\xi$ 零化 $\Gr_0M_\alpha$ 并把 $\Gr_pM_\alpha$ 映到 $\Gr_{p+1}M_\alpha$（$p>0$）上。由此可知 $\Gr M_\alpha$ 的零化子是 $k[z,\xi]$ 中由 $z\xi$ 生成的理想。故特征簇 $\operatorname{Ch}(M_\alpha)$ 是 $k^2$ 中直线 $\{z=0\}$ 和 $\{\xi=0\}$ 的并。

现设 $\alpha\in\mathbb{Z}$。则 $E$ 的特征值是整数。若 $v$ 是 $E$ 的属于特征值 $m\neq 0$ 的非零特征向量，则 $\partial v$ 是 $E$ 的属于特征值 $m-1$ 的特征向量，且 $z\partial v=mv\neq 0$。因此 $z$ 把 $E$ 的属于特征值 $q\neq -1$ 的所有特征空间映到属于特征值 $q+1$ 的特征空间。

先设 $n=-\alpha>0$。则 $z^{n-1}$ 的陪集是 $E$ 的属于特征值 $-1$ 的特征向量。因此 $z$ 把 $E$ 的属于特征值 $-1$ 的特征空间映到属于特征值 $0$ 的特征空间。故此时可选取特征值 $m\in\mathbb{Z}$ 的特征空间的基向量 $v_m$，使得对 $m\in\mathbb{Z}$ 有 $zv_m=v_{m+1}$。有
$$\partial v_m=\partial z v_{m-1}=[\partial,z]v_{m-1}+Ev_{m-1}=mv_{m-1}$$
对所有 $m\in\mathbb{Z}$ 成立。这意味着所有 $M_{-n}$（$n>0$）互相同构。

此外，通过考察 $z$ 和 $\partial$ 的作用，可见向量 $v_m$（$m\in\mathbb{Z}_+$）张成一个同构于 8.2 中 $\mathcal{O}_1$ 的 $D$-子模 $N_{-n}$。特别地 $M_{-n}$ 是可约的。

显然 $zv_{-1}\in N_{-n}$。由此可知 $v_{-1}$ 在 $M_{-n}/N_{-n}$ 中的陪集 $\delta$ 满足 $z\delta=0$。

$E$ 在 $L_{-n}=M_{-n}/N_{-n}$ 上的谱由所有严格负整数构成。因此 $\partial$ 在 $L_{-n}$ 上是单射。故 $\delta^{(m)}=\partial^m\delta$ 是 $E$ 的属于特征值 $-(m+1)$（$m\in\mathbb{Z}_+$）的非零特征向量，即它们与 $v_{-(m+1)}$ 的陪集成比例。显然对所有 $m>0$ 有
$$z\delta^{(m)}=z\partial\delta^{(m-1)}=-m\delta^{(m-1)}.$$
由此可知 $L_{-n}$ 同构于 8.3 中描述的 $D$-模 $\Delta_1$。故有正合列
$$0\longrightarrow\mathcal{O}_1\longrightarrow M_{-n}\longrightarrow\Delta_1\longrightarrow 0,$$
且此正合列不分裂。特别地，所有这些 $D$-模都同构于 $M_{-1}$。

由 Fourier 变换可见 $D$-模 $M_n$（$n\ge 0$）同构于 $M_0$。此外有正合列
$$0\longrightarrow\Delta_1\longrightarrow M_n\longrightarrow\mathcal{O}_1\longrightarrow 0,$$
它也不分裂。

由于 $\mathcal{O}_1$ 和 $\Delta_1$ 由 8.2 和 8.3 是完整的，由 8.1 知 $M_n$（$n\in\mathbb{Z}$）是完整的。此外由 7.4 得
$$\operatorname{Ch}(M_n)=\operatorname{Ch}(\mathcal{O}_1)\cup\operatorname{Ch}(\Delta_1)$$
对所有 $n\in\mathbb{Z}$ 成立。故由 8.2 和 8.3，它们等于 $k^2$ 中直线 $\{z=0\}$ 和 $\{\xi=0\}$ 的并。

由上述例子可见特征簇不能决定相应的 $D$-模。此外，不可约完整 $D(n)$-模的特征簇可以是可约的。
\end{example}

现在把上述例子中 $M_{-1}$ 的构造推广。

设 $M$ 是 $D(n)$-模，$P\in k[X_1,X_2,\dots,X_n]$。则在 $M$ 的局部化 $M_P$ 上可定义 $k$-线性映射 $\partial_i:M_P\to M_P$ 为
$$\partial_i\Bigl(\frac{m}{P^p}\Bigr)=-\frac{p\,\partial_i(P)m}{P^{p+1}}+\frac{\partial_i m}{P^p}$$
对任意 $m\in M$ 和 $p\in\mathbb{Z}_+$。通过直接计算可检验
$$[\partial_i,\partial_j]\Bigl(\frac{m}{P^p}\Bigr)=0,\qquad [\partial_i,x_j]\Bigl(\frac{m}{P^p}\Bigr)=\delta_{ij}\frac{m}{P^p}$$
对任意 $1\le i,j\le n$ 和 $p\in\mathbb{Z}_+$。由 5.11 这在 $M_P$ 上定义了 $D(n)$-模结构。

\begin{proposition}
设 $M$ 是完整 $D(n)$-模，$P\in k[X_1,X_2,\dots,X_n]$。则 $M_P$ 是完整 $D(n)$-模。
\end{proposition}

\begin{proof}
显然可设 $P\neq 0$。设 $F M$ 是 $M$ 上好滤过，使得 $p\le 0$ 时 $F_pM=0$，且 $m=\deg P$。定义 $p<0$ 时 $F_pM_P=0$，
$$F_pM_P=\Bigl\{\frac{v}{P^p}\ \Big|\ v\in F_{(m+1)p}M\Bigr\}\quad(p\in\mathbb{Z}_+).$$
显然 $F_pM_P$（$p\in\mathbb{Z}$）是 $M_P$ 的向量子空间。

设 $w\in F_pM_P$，$p\ge 0$。则对某 $v\in F_{(m+1)p}M$ 有 $w=\frac{Pv}{P^{p+1}}$。由于 $Pv\in F_{(m+1)p+m}M\subset F_{(m+1)(p+1)}M$，可见 $w\in F_{p+1}M_P$。这证明滤过 $F M_P$ 是递增的。

设 $v\in F_qM$。则对任意 $s\in\mathbb{Z}_+$ 有 $\frac{v}{P^p}=\frac{P^sv}{P^{p+s}}$。此外对任意 $s\in\mathbb{Z}_+$ 有 $P^sv\in F_{q+sm}M$。而且当 $s\ge q-(m+1)p$ 时，$(m+1)(p+s)-(q+sm)=s+(m+1)p-q\ge 0$。故
$$P^sv\in F_{q+sm}M\subset F_{(m+1)(p+s)}M,\quad \frac{v}{P^p}\in F_{p+s}M_P.$$
因此滤过 $F M_P$ 是穷尽的。

还需证明它是 $D(n)$-模滤过。首先，对 $v\in F_{(m+1)p}M$，有 $x_iPv\in F_{(m+1)(p+1)}M$，故 $x_i\frac{v}{P^p}=\frac{Pv}{P^{p+1}}x_i\in F_{p+1}M_P$。此外，
$$\partial_i\Bigl(\frac{v}{P^p}\Bigr)=\frac{-p\,\partial_i(P)v+P\,\partial_i v}{P^{p+1}},$$
且 $-p\,\partial_i(P)v+P\,\partial_i v\in F_{(m+1)(p+1)}M$，故 $\partial_i\frac{v}{P^p}\in F_{p+1}M_P$。

因此在 $M_P$ 上构造了穷尽 $D(n)$-模滤过。由于对 $p\in\mathbb{Z}_+$，
$$\dim_k F_pM_P\le\dim_k F_{(m+1)p}M\le e(M)\frac{((m+1)p)^n}{n!}+\text{（关于 $p$ 的低阶项）},$$
由 8.4 知 $M_P$ 是完整的。
\end{proof}

\begin{corollary}
设 $P\in k[X_1,X_2,\dots,X_n]$。则 $k[X_1,X_2,\dots,X_n]_P$ 是完整 $D(n)$-模。
\end{corollary}
\begin{annote}
“完整”（holonomic）指 $\dim\mathrm{Ch}(M)=n$；因 $\mathrm{Ch}(M)$ 总是维数 $\ge n$ 的锥，这是可能的最小值。等价刻画：$M$ 的每个非零元在 $D(n)$ 上的零化子理想的维数…（Gabber 定理）——nLab 直言完整当且仅当特征簇维数为 $n$，此时特征簇既是锥又是 Lagrangian 子簇。这也解释了为什么完整模“几乎”是有限维的：解空间有限维（Riemann–Hilbert 对应的 D-模侧）。[来源：nLab 条目 holonomic+D-module（Last revised 2014-02-16）；原书 §8, p38–40]
\end{annote}



\section{外张量积}

以下设 $X=k^n$，$Y=k^m$，并以 $D_X$、$D_Y$ 记相应的多项式系数微分算子代数。则可考虑代数 $D_X\boxtimesext D_Y$，它作为 $k$-向量空间等于 $D_X\otimes_k D_Y$，乘法定义为 $(T\otimes S)(T'\otimes S')=TT'\otimes SS'$（$T,T'\in D_X$，$S,S'\in D_Y$）。称 $D_X\boxtimesext D_Y$ 为 $D_X$ 和 $D_Y$ 的\emph{外张量积}。

以下结果显然。

\begin{lemma}
$D_X\boxtimesext D_Y=D_{X\times Y}$。
\end{lemma}

若 $M$、$N$ 分别是 $D_X$-、$D_Y$-模，可定义 $D_{X\times Y}$-模 $M\boxtimesext N$，它作为 $k$-向量空间等于 $M\otimes_k N$，$D_X\boxtimesext D_Y=D_{X\times Y}$ 的作用由 $(T\otimes S)(m\otimes n)=Tm\otimes Sn$ 给出（$T\in D_X$，$S\in D_Y$，$m\in M$，$n\in N$）。

\begin{lemma}
设 $M$ 是有限生成 $D_X$-模，$N$ 是有限生成 $D_Y$-模。则 $M\boxtimesext N$ 是有限生成 $D_{X\times Y}$-模。
\end{lemma}

\begin{proof}
设 $e_1,e_2,\dots,e_p$ 和 $f_1,f_2,\dots,f_q$ 分别是 $M$ 和 $N$ 的生成元。则对任意 $m\in M$ 和 $n\in N$，有 $m=\sum_i T_ie_i$（$T_i\in D_X$），$n=\sum_j S_jf_j$（$S_j\in D_Y$）。这意味着 $m\otimes n=\sum_{i,j}T_ie_i\otimes S_jf_j=\sum_{i,j}(T_i\otimes S_j)(e_i\otimes f_j)$，且 $e_i\otimes f_j$（$1\le i\le p$，$1\le j\le q$）生成 $M\boxtimesext N$。
\end{proof}

本节的主要目标是证明如下结果。

\begin{theorem}
设 $M$ 是有限生成 $D_X$-模，$N$ 是有限生成 $D_Y$-模。则
$$d(M\boxtimesext N)=d(M)+d(N).$$
\end{theorem}

这有如下重要推论。

\begin{corollary}
设 $M$ 是完整 $D_X$-模，$N$ 是完整 $D_Y$-模。则 $M\boxtimesext N$ 是完整 $D_{X\times Y}$-模。
\end{corollary}

设 $D_X$ 和 $D_Y$ 配置 Bernstein 滤过。设 $M$、$N$ 分别是带好滤过 $F M$、$F N$ 的有限生成 $D_X$-、$D_Y$-模。定义 $M\boxtimesext N$ 上的\emph{乘积滤过}为
$$F_j(M\boxtimesext N)=\sum_{p+q=j}F_pM\otimes_k F_qN$$
对任意 $j\in\mathbb{Z}$。显然 $D_X\boxtimesext D_Y=D_{X\times Y}$ 上的乘积滤过与 Bernstein 滤过一致。因此 $F(M\boxtimesext N)$ 是穷尽 Hausdorff 的 $D_{X\times Y}$-模滤过。

为证明此滤过是好的，需要线性代数中的一些准备。从如下引理开始。

\begin{lemma}
设 $M,M',N,N'$ 是 $k$ 上线性空间，$\varphi:M\to M'$ 和 $\psi:N\to N'$ 是线性映射。则它们定义线性映射 $\varphi\otimes\psi:M\otimes_k N\to M'\otimes_k N'$。有
\begin{enumerate}[label=(\roman*)]
  \item $\operatorname{im}(\varphi\otimes\psi)=\operatorname{im}\varphi\otimes\operatorname{im}\psi$；
  \item $\ker(\varphi\otimes\psi)=\ker\varphi\otimes_k N+M\otimes_k\ker\psi$。
\end{enumerate}
\end{lemma}

\begin{proof}
(i) 由定义显然。

(ii) 由 (i)，为证 (ii) 可设 $\varphi$ 和 $\psi$ 是满射。这时有短正合列
$$0\longrightarrow M''\longrightarrow M\xrightarrow{\ \varphi\ }M'\longrightarrow 0,$$
其中 $M''=\ker\varphi$，以及
$$0\longrightarrow N''\longrightarrow N\xrightarrow{\ \psi\ }N'\longrightarrow 0,$$
其中 $N''=\ker\psi$。

显然 $\varphi\otimes\psi=(\varphi\otimes\operatorname{id}_{N'})\circ(\operatorname{id}_M\otimes\psi)$。由于与 $N'$ 做张量是正合的，第一正合列蕴含序列
$$0\longrightarrow M''\otimes_k N'\longrightarrow M\otimes_k N'\xrightarrow{\ \varphi\otimes\operatorname{id}\ }M'\otimes_k N'\longrightarrow 0$$
是正合的。故 $\ker(\varphi\otimes\operatorname{id}_{N'})=M''\otimes_k N'=\ker\varphi\otimes_k N'$。因此 $M\otimes_k N$ 中元素 $z$ 属于 $\varphi\otimes\psi$ 的核，当且仅当 $(\operatorname{id}_M\otimes\psi)(z)$ 属于 $\ker\varphi\otimes_k N'$。

由于
$$0\longrightarrow M\otimes_k N''\longrightarrow M\otimes_k N\xrightarrow{\ \operatorname{id}_M\otimes\psi\ }M\otimes_k N'\longrightarrow 0$$
也正合，$\ker\varphi\otimes_k N$ 满射到 $\ker\varphi\otimes_k N'$ 上，且 $\ker(\operatorname{id}_M\otimes\psi)=M\otimes_k N''=M\otimes_k\ker\psi$。因此 $z$ 属于 $\varphi\otimes\psi$ 的核当且仅当
$$z\in\ker\varphi\otimes_k N+M\otimes_k\ker\psi.$$
\end{proof}

\begin{lemma}
设 $X_1,X_2,\dots,X_n$ 是张成线性空间 $X$ 的线性子空间。若对 $1\le i\le n$ 有
$$X_i\cap\sum_{j\neq i}X_j=\{0\},$$
则线性空间 $X$ 是 $X_1,X_2,\dots,X_n$ 的直和。
\end{lemma}

\begin{proof}
设 $x_i\in X_i$（$1\le i\le n$）满足 $x_1+x_2+\cdots+x_n=0$。则 $x_i=-\sum_{j\neq i}x_j\in X_i\cap\sum_{j\neq i}X_j$，由假设对每个 $1\le i\le n$ 它等于 $0$。
\end{proof}

现在描述 $\Gr(M\boxtimesext N)$。设 $j\in\mathbb{Z}$。若 $p+q=j$，有良定义的 $k$-线性映射 $F_pM\otimes_k F_qN\to F_j(M\boxtimesext N)$，故有良定义的 $k$-线性映射 $F_pM\otimes_k F_qN\to\Gr_j(M\boxtimesext N)$。由 9.5，自然映射 $F_pM\otimes_k F_qN\to\Gr_pM\otimes_k\Gr_qN$ 的核为 $F_{p-1}M\otimes_k F_qN+F_pM\otimes_k F_{q-1}N$，即含于 $F_{j-1}(M\boxtimesext N)$。故线性映射 $F_pM\otimes_k F_qN\to\Gr_j(M\boxtimesext N)$ 通过 $\Gr_pM\otimes_k\Gr_qN$ 分解。这导出线性映射
$$\pi:\bigoplus_{p+q=j}\Gr_pM\otimes_k\Gr_qN\longrightarrow\Gr_j(M\boxtimesext N).$$
显然由其构造此映射是满射。此外它在直和中每个加项 $\Gr_pM\otimes_k\Gr_qN$ 上的限制是单射。设 $X_{p,q}$ 为 $\Gr_pM\otimes_k\Gr_qN$ 在 $\Gr_j(M\boxtimesext N)$ 中的像。由于
$$F_{p-1}M\otimes_k F_qN+F_pM\otimes_k F_{q-1}N=(F_pM\otimes_k F_qN)\cap\Bigl(\sum_{\substack{p'+q'=j\\ p'\neq p,\ q'\neq q}}F_{p'}M\otimes_k F_{q'}N\Bigr),$$
可见
$$X_{p,q}\cap\Bigl(\sum_{\substack{p'+q'=j\\ p'\neq p,\ q'\neq q}}X_{p',q'}\Bigr)=\{0\}.$$
故由 9.6 此映射是同构。这意味着对任意 $j\in\mathbb{Z}$，
$$\Gr_j(M\boxtimesext N)=\bigoplus_{p+q=j}\Gr_pM\otimes_k\Gr_qN.$$

若类似地定义代数 $\Gr D_X\boxtimesext\Gr D_Y$，其分次由总次数给出，则 $\Gr D_X\boxtimesext\Gr D_Y=\Gr D_{X\times Y}$。此外，$\Gr M\boxtimesext\Gr N$ 成为分次 $\Gr D_{X\times Y}$-模，由上述讨论同构于 $\Gr(M\boxtimesext N)$。由于滤过 $F M$ 和 $F N$ 是好的，由 3.1，$\Gr M$ 和 $\Gr N$ 分别是有限生成 $\Gr D_X$-、$\Gr D_Y$-模。由 9.2 的类似结论，$\Gr(M\boxtimesext N)$ 是有限生成 $\Gr D_{X\times Y}$-模。这意味着乘积滤过是 $M\boxtimesext N$ 上的好滤过。

设
$$P(M,t)=\sum_{p\in\mathbb{Z}}\dim_k(\Gr_p M)\,t^p,\qquad P(N,t)=\sum_{q\in\mathbb{Z}}\dim_k(\Gr_q N)\,t^q$$
为 $\Gr M$ 和 $\Gr N$ 的 Poincaré 级数。则
\begin{align*}
P(M,t)P(N,t)&=\sum_{p\in\mathbb{Z}}\sum_{q\in\mathbb{Z}}\dim_k(\Gr_p M)\dim_k(\Gr_q N)\,t^{p+q}\\
&=\sum_{j\in\mathbb{Z}}\Bigl(\sum_{p+q=j}\dim_k(\Gr_p M)\dim_k(\Gr_q N)\Bigr)t^j\\
&=\sum_{j\in\mathbb{Z}}\Bigl(\sum_{p+q=j}\dim_k(\Gr_p M\otimes_k\Gr_q N)\Bigr)t^j\\
&=\sum_{j\in\mathbb{Z}}\dim_k\Gr_j(M\boxtimesext N)\,t^j=P(M\boxtimesext N,t)
\end{align*}
是 $M\boxtimesext N$ 的 Poincaré 级数。因此 $P(M\boxtimesext N,t)$ 在 $1$ 处极点的阶数是 $P(M,t)$ 和 $P(N,t)$ 极点阶数之和。由 1.5 可见这立即蕴含 9.3。

也可通过考虑特征簇推导 9.3。把 $D_X$、$D_Y$、$D_{X\times Y}$ 看作按微分算子次数滤过的环。设 $M$、$N$ 是带好滤过 $F M$、$F N$ 的有限生成 $D_X$-、$D_Y$-模。如上定义 $M\boxtimesext N$ 上的 $D_{X\times Y}$-模滤过 $F(M\boxtimesext N)$。则如同上述论证，$F(M\boxtimesext N)$ 是 $M\boxtimesext N$ 的好滤过。设 $I$ 为 $\Gr M$ 在 $\Gr D_X$ 中的零化子，$J$ 为 $\Gr N$ 在 $\Gr D_Y$ 中的零化子。则由 9.5，$\Gr(M\boxtimesext N)$ 的零化子等于 $\Gr D_{X\times Y}=\Gr D_X\boxtimesext\Gr D_Y$ 中的理想 $I\otimes_k\Gr D_Y+\Gr D_X\otimes_k J$。

可把 $\Gr D_X$ 与多项式环 $k[x_1,\dots,x_n,\xi_1,\dots,\xi_n]$ 等同，把 $\Gr D_Y$ 与
$k[y_1,\dots,y_m,\eta_1,\dots,\eta_m]$ 等同。此外可把 $\Gr D_{X\times Y}$ 与多项式环
$$k[x_1,\dots,x_n,y_1,\dots,y_m,\xi_1,\dots,\xi_n,\eta_1,\dots,\eta_m]$$
等同。则 $\Gr(M\boxtimesext N)$ 的零化子对应于该环中由 $I$ 和 $J$ 的像生成的理想。若定义映射
$q:k^{2n}\times k^{2m}\to k^{2(n+m)}$ 为
\begin{multline*}
q(x_1,\dots,x_n,\xi_1,\dots,\xi_n,y_1,\dots,y_m,\eta_1,\dots,\eta_m)\\
=(x_1,\dots,x_n,y_1,\dots,y_m,\xi_1,\dots,\xi_n,\eta_1,\dots,\eta_m),
\end{multline*}
则有如下结果。

\begin{theorem}
设 $M$ 和 $N$ 是有限生成 $D_X$-、$D_Y$-模。则
$$\operatorname{Ch}(M\boxtimesext N)=q\bigl(\operatorname{Ch}(M)\times\operatorname{Ch}(N)\bigr).$$
\end{theorem}

这进而蕴含 $\dim\operatorname{Ch}(M\boxtimesext N)=\dim\operatorname{Ch}(M)+\dim\operatorname{Ch}(N)$，由 7.7 得到 9.3 的另一个证明。

利用 7.1 的证明并如上论证，或利用 7.5，也可得到如下结果。

\begin{proposition}
设 $M$ 和 $N$ 是有限生成 $D_X$-、$D_Y$-模。则
$$\operatorname{supp}(M\boxtimesext N)=\operatorname{supp}(M)\times\operatorname{supp}(N).$$
\end{proposition}

\section{逆像}

设 $X=k^n$，$Y=k^m$，以 $x_1,x_2,\dots,x_n$ 和 $y_1,y_2,\dots,y_m$ 分别记 $X$ 和 $Y$ 上的典范坐标函数。以 $R(X)=k[x_1,x_2,\dots,x_n]$ 和 $R(Y)=k[y_1,y_2,\dots,y_m]$ 分别记 $X$ 和 $Y$ 上正则函数环。

设 $F:X\longrightarrow Y$ 是多项式映射，即
$$F(x_1,x_2,\dots,x_n)=\bigl(F_1(x_1,\dots,x_n),F_2(x_1,\dots,x_n),\dots,F_m(x_1,\dots,x_n)\bigr),$$
其中 $F_i\in R(X)$。则 $F$ 定义环同态 $\varphi_F:R(Y)\to R(X)$：$\varphi_F(P)=P\circ F$（$P\in R(Y)$）。因此可把 $R(X)$ 看作 $R(Y)$-模。于是可定义函子 $F^*$，从 $R(Y)$-模范畴 $\mathcal{M}(R(Y))$ 到 $R(X)$-模范畴 $\mathcal{M}(R(X))$，由
$$F^*(N)=R(X)\otimes_{R(Y)}N$$
给出（对任意 $R(Y)$-模 $N$）。显然 $F^*:\mathcal{M}(R(Y))\to\mathcal{M}(R(X))$ 是右正合函子。我们称它为从 $\mathcal{M}(R(Y))$ 到 $\mathcal{M}(R(X))$ 的\emph{逆像函子}。

现在想把此函子扩充到 $D$-模。以 $D_X$、$D_Y$ 分别记 $X$、$Y$ 上多项式系数微分算子代数。若 $N$ 是左 $D_Y$-模，想在逆像 $F^*(N)$ 上定义 $D_X$-模结构。（如 §6 开头所述，转置函子是左 $D$-模范畴与右 $D$-模范畴的等价，故可类似处理右模。）先考虑双线性映射
$$(P,v)\longmapsto\frac{\partial P}{\partial x_i}\otimes v+\sum_{j=1}^m P\frac{\partial F_j}{\partial x_i}\otimes\frac{\partial}{\partial y_j}v,$$
从 $R(X)\times N$ 到 $R(X)\otimes_{R(Y)}N$。由于对任意 $Q\in R(Y)$，
\begin{align*}
&\frac{\partial P\,(Q\circ F)}{\partial x_i}\otimes v+\sum_j P\frac{\partial F_j}{\partial x_i}(Q\circ F)\otimes\frac{\partial}{\partial y_j}v\\
&=\frac{\partial P}{\partial x_i}\otimes Qv+\sum_j P\Bigl(\frac{\partial Q}{\partial y_j}\circ F\Bigr)\frac{\partial F_j}{\partial x_i}\otimes v+\sum_j P(Q\circ F)\frac{\partial F_j}{\partial x_i}\otimes\frac{\partial}{\partial y_j}v\\
&=\frac{\partial P}{\partial x_i}\otimes Qv+\sum_j P\frac{\partial F_j}{\partial x_i}\otimes\Bigl(\frac{\partial Q}{\partial y_j}v+Q\frac{\partial}{\partial y_j}v\Bigr)\\
&=\frac{\partial P}{\partial x_i}\otimes Qv+\sum_j P\frac{\partial F_j}{\partial x_i}\otimes\frac{\partial}{\partial y_j}(Qv),
\end{align*}
此映射通过 $F^*(N)$ 的线性自同态分解，记为 $\dfrac{\partial}{\partial x_i}$。直接计算得
$$\Bigl[\frac{\partial}{\partial x_i},\frac{\partial}{\partial x_j}\Bigr](P\otimes v)=0,\qquad \Bigl[\frac{\partial}{\partial x_i},x_j\Bigr](P\otimes v)=\delta_{ij}(P\otimes v),$$
故由 5.11，$F^*(N)$ 具有自然的左 $D_X$-模结构。

其结构可用另一方式描述。令
$$D_{X\to Y}=F^*(D_Y)=R(X)\otimes_{R(Y)}D_Y.$$
则如上所述，$D_{X\to Y}$ 有左 $D_X$-模结构。但它也有由 $D_Y$ 中右乘给出的右 $D_Y$-模结构。这两个作用显然可交换，故 $D_{X\to Y}$ 是（左 $D_X$，右 $D_Y$）-双模。此外，对任意 $D_Y$-模 $N$ 有
$$F^*(N)=R(X)\otimes_{R(Y)}N=\bigl(R(X)\otimes_{R(Y)}D_Y\bigr)\otimes_{D_Y}N=D_{X\to Y}\otimes_{D_Y}N,$$
且 $D_X$ 在 $F^*(N)$ 上的作用由最后一式中第一个因子上的作用给出。

我们记此 $D_X$-模为 $F^+(N)$，称为 $D_Y$-模 $N$ 的\emph{逆像}。

显然逆像函子 $F^+$ 是从 $\mathcal{M}(D_Y)$ 到 $\mathcal{M}^L(D_X)$ 的右正合函子。它的左导出函子 $L^iF^+$ 对左 $D_Y$-模 $N$ 由
$$L^iF^+(N)=\operatorname{Tor}^{-i}_{D_Y}(D_{X\to Y},N)$$
给出。

设 $\operatorname{For}$ 为从 $D_X$-模范畴（相应地从 $D_Y$-模范畴）到 $R(X)$-模范畴（相应地 $R(Y)$-模范畴）的遗忘函子。则如下函子图交换：
$$\begin{tikzcd}
\mathcal{M}(D_Y) \arrow[r,"F^+"] \arrow[d,"\operatorname{For}"'] & \mathcal{M}(D_X) \arrow[d,"\operatorname{For}"]\\
\mathcal{M}(R(Y)) \arrow[r,"F^*"'] & \mathcal{M}(R(X))
\end{tikzcd}$$

我们断言对左导出函子也有类似陈述，即：

\begin{proposition}
对任意 $i\in\mathbb{Z}$，如下函子图交换：
$$\begin{tikzcd}
\mathcal{M}(D_Y) \arrow[r,"L^iF^+"] \arrow[d,"\operatorname{For}"'] & \mathcal{M}(D_X) \arrow[d,"\operatorname{For}"]\\
\mathcal{M}(R(Y)) \arrow[r,"L^iF^*"'] & \mathcal{M}(R(X))
\end{tikzcd}$$
\end{proposition}

\begin{proof}
设 $F^\cdot$ 是 $D_Y$-模 $N$ 由自由 $D_Y$-模构成的一个左分解。由于自由 $D_Y$-模由 5.10 也是自由 $R(Y)$-模，由上述注记，对任意 $i\in\mathbb{Z}$ 有
$$\operatorname{For}(L^iF^+(N))=\operatorname{For}(H^i(F^+(F^\cdot)))=H^i(\operatorname{For}(F^+(F^\cdot)))=H^i(F^*(\operatorname{For}F^\cdot))=L^iF^*(\operatorname{For}N).$$
\end{proof}

现在研究导出逆像对态射复合的行为。先需要一个无环性结果。

\begin{lemma}
设 $P$ 是射影左 $D_Y$-模。则 $F^*(P)$ 是射影 $R(X)$-模。
\end{lemma}

\begin{proof}
设 $P$ 是射影 $D_Y$-模。则它是自由 $D_Y$-模 $(D_Y)^{(I)}$ 的直和因子。这意味着 $F^+(P)$ 是 $F^+(D_Y^{(I)})$ 的直和因子。由于 $D_Y$ 是自由 $R(Y)$-模，$\operatorname{For}(F^+(D_Y^{(I)}))=R(X)\otimes_{R(Y)}D_Y^{(I)}$ 是自由 $R(X)$-模。
\end{proof}

\begin{theorem}
设 $X=k^n$，$Y=k^m$，$Z=k^p$，$F:X\to Y$ 和 $G:Y\to Z$ 是多项式映射。则
\begin{enumerate}[label=(\roman*)]
  \item 逆像函子 $(G\circ F)^+$（从 $\mathcal{M}^L(D_Z)$ 到 $\mathcal{M}^L(D_X)$）同构于 $F^+\circ G^+$；
  \item 对任意左 $D_Z$-模 $N$，存在谱序列，其 $E_2$ 项为 $E_2^{pq}=L^pF^+(L^qG^+(N))$，收敛于 $L^{p+q}(G\circ F)^+(N)$。
\end{enumerate}
\end{theorem}

\begin{proof}
(i) 先考虑多项式环结构。对任意 $D_Z$-模 $N$，
$$(G\circ F)^*(N)=R(X)\otimes_{R(Z)}N=R(X)\otimes_{R(Y)}(R(Y)\otimes_{R(Z)}N)=F^*(G^*(N)).$$
另一方面，对任意 $P\in R(X)$ 和 $v\in N$，
\begin{align*}
\frac{\partial}{\partial x_i}(P\otimes v)&=\frac{\partial}{\partial x_i}\bigl(P\otimes(1\otimes v)\bigr)\\
&=\frac{\partial P}{\partial x_i}\otimes(1\otimes v)+\sum_{j=1}^m P\frac{\partial F_j}{\partial x_i}\otimes\frac{\partial}{\partial y_j}(1\otimes v)\\
&=\frac{\partial P}{\partial x_i}\otimes v+\sum_{j=1}^m P\frac{\partial F_j}{\partial x_i}\otimes\Bigl(\sum_{k=1}^p\frac{\partial G_k}{\partial y_j}\otimes\frac{\partial}{\partial z_k}v\Bigr)\\
&=\frac{\partial P}{\partial x_i}\otimes v+\sum_{k=1}^p\sum_{j=1}^m P\Bigl(\frac{\partial F_j}{\partial x_i}\frac{\partial G_k}{\partial y_j}\Bigr)\circ F\otimes\frac{\partial}{\partial z_k}v\\
&=\frac{\partial P}{\partial x_i}\otimes v+\sum_{k=1}^p P\frac{\partial(G_k\circ F)}{\partial x_i}\otimes\frac{\partial}{\partial z_k}v.
\end{align*}
故 $D_X$ 作用一致。

(ii) 由 10.1 和 10.2，对任意射影 $D_Z$-模 $P$，逆像 $G^+(P)$ 是 $F^+$-无环的。因此结论由 Grothendieck 谱序列得出。
\end{proof}

这立即有如下推论。

\begin{corollary}
设 $X=k^n$，$Y=k^m$，$Z=k^p$，$F:X\to Y$ 和 $G:Y\to Z$ 是多项式映射。则
\begin{enumerate}[label=(\roman*)]
  \item $D_{X\to Z}=D_{X\to Y}\otimes_{D_Y}D_{Y\to Z}$；
  \item $\operatorname{Tor}^{D_Y}_j(D_{X\to Y},D_{Y\to Z})=0$ 对 $j\in\mathbb{N}$。
\end{enumerate}
\end{corollary}

\begin{proof}
(i) 由 10.3.(i) 有
$$D_{X\to Z}=(G\circ F)^+(D_Z)=F^+(G^+(D_Z))=F^+(D_{Y\to Z})=D_{X\to Y}\otimes_{D_Y}D_{Y\to Z}.$$

(ii) 如在 10.3.(ii) 的证明中所述，由 10.1 和 10.2 可见 $D_{Y\to Z}=G^+(D_Z)$ 是 $F^+$-无环的。故对 $j>0$，
$$0=L^{-j}F^+(G^+(D_Z))=\operatorname{Tor}^{-j}_{D_Y}(D_{X\to Y},D_{Y\to Z}).$$
\end{proof}

现在考虑两个简单的例子。首先设 $p$ 是 $X\times Y$ 到 $Y$ 的投影，由 $p(x,y)=y$（$x\in X$，$y\in Y$）定义。则众所周知，$R(X\times Y)=R(X)\boxtimesext R(Y)$。因此对 $R(Y)$-模 $N$ 有
$$p^*(N)=R(X\times Y)\otimes_{R(Y)}N=\bigl(R(X)\boxtimesext R(Y)\bigr)\otimes_{R(Y)}N=R(X)\boxtimesext N$$
作为 $R(X\times Y)=R(X)\boxtimesext R(Y)$ 上的模。另一方面，若 $N$ 是 $D_Y$-模，则立即看出 $\frac{\partial}{\partial x_i}$ 和 $\frac{\partial}{\partial y_j}$ 的作用也一致，即
$$p^+(N)=R(X)\boxtimesext N.$$
由 9.3 和 9.4 立即得到如下结果。

\begin{proposition}
设 $p:X\times Y\longrightarrow Y$ 是典范投影。则
\begin{enumerate}[label=(\roman*)]
  \item $p^+$ 是从 $\mathcal{M}^L(D_Y)$ 到 $\mathcal{M}^L(D_{X\times Y})$ 的正合函子；
  \item 对任意左 $D_Y$-模 $N$，$p^+(N)=R(X)\boxtimesext N$；
  \item 若 $N$ 是有限生成左 $D_Y$-模，则 $p^+(N)$ 是有限生成 $D_{X\times Y}$-模；
  \item 对任意有限生成左 $D_Y$-模 $N$，$d(p^+(N))=d(N)+n$。
\end{enumerate}
特别地，有限生成左 $D_Y$-模 $N$ 是完整的当且仅当 $p^+(N)$ 是完整的。
\end{proposition}

现在考虑另一个例子。设 $i$ 是 $X$ 到 $X\times Y$ 的典范注入，由 $i(x)=(x,0)$（$x\in X$）给出。则
$$D_{X\to X\times Y}=i^*(D_{X\times Y})=R(X)\otimes_{R(X)\boxtimesext R(Y)}(D_X\boxtimesext D_Y)=D_X\boxtimesext D_Y/\bigl((y_1,y_2,\dots,y_m)D_Y\bigr),$$
其中 $D_X$ 显然通过第一因子上的左乘作用，而 $D_{X\times Y}=D_X\boxtimesext D_Y$ 通过右乘作用。

此外设 $m=1$。则有正合列
$$0\longrightarrow D_Y\xrightarrow{\ y_1\ }D_Y\longrightarrow D_Y/y_1D_Y\longrightarrow 0,$$
其中第二个箭头由左乘 $y_1$ 给出。与 $D_X$ 作张量，得到左 $D_X$-模（对左乘）和右 $D_{X\times Y}$-模（对右乘）的短正合列
$$0\longrightarrow D_{X\times Y}\xrightarrow{\ y_1\ }D_{X\times Y}\longrightarrow D_{X\to X\times Y}\longrightarrow 0.$$
因此可把此正合列的前两项看作 $D_{X\to X\times Y}$ 的、由作为左 $D_{X\times Y}$-模自由的（左 $D_X$，右 $D_{X\times Y}$）-双模构成的左分解。故对 $D_{X\times Y}$-模 $N$，复形
$$0\longrightarrow N\xrightarrow{\ y_1\ }N\longrightarrow 0$$
的上同调计算导出逆像。特别地，有如下引理。

\begin{lemma}
设 $\dim Y=1$，$i$ 是 $X$ 到 $X\times Y$ 的典范注入。则对任意 $D_{X\times Y}$-模 $N$ 有
\begin{enumerate}[label=(\roman*)]
  \item $i^+(N)=\operatorname{coker}y_1$；
  \item $L^{-1}i^+(N)=\ker y_1$；
  \item 当 $p\neq 0,-1$ 时 $L^pi^+(N)=0$。
\end{enumerate}
特别地，$i^+$ 的左同调维数 $\le 1$。
\end{lemma}

最后陈述对任意 $Y$ 有显然的推广。

\begin{lemma}
设 $i$ 是 $X$ 到 $X\times Y$ 的典范注入。则 $i^+$ 的左同调维数 $\le\dim Y$。
\end{lemma}

\begin{proof}
对 $\dim Y$ 归纳。$\dim Y=1$ 时结论已证。可把 $Y$ 表为 $Y=Y'\times Y''$，其中 $Y'=k^{m-1}$，$Y''=k$。以 $i'$ 记 $X$ 到 $X\times Y'$ 的典范包含，以 $j$ 记 $X\times Y'$ 到 $X\times Y'\times Y''=X\times Y$ 的典范包含。则 $i=j\circ i'$。此外由 10.6，$j^+$ 的左同调维数 $\le 1$，而由归纳假设，$i'^+$ 的左同调维数 $\le\dim Y'$。因此由 10.3.(ii) 的 Grothendieck 谱序列断定，导出逆像 $L^{-p}i^+$ 对 $p\ge\dim Y'+1=\dim Y$ 为零。
\end{proof}

设 $F:X\longrightarrow X$ 是 $X$ 的同构，$G$ 是其逆。则由 $\alpha(f)=f\circ F$ 定义的映射 $\alpha:R(X)\to R(X)$ 是环 $R(X)$ 的自同构。其逆为 $\beta$：$\beta(f)=f\circ G$（$f\in R(X)$）。若 $M$ 是 $R(X)$-模，则作为 $k$ 上的线性空间，$F^*(M)$ 通过映射 $\varphi:m\longmapsto1\otimes m$ 同构于 $M$。另一方面，对 $f\in R(X)$ 有
$$f\varphi(m)=f\otimes m=f\circ G\circ F\otimes m=1\otimes(f\circ G)m=\varphi(\beta(f)m)$$
（$m\in M$），即 $R(X)$-模 $F^*(M)$ 同构于 $M$，其 $R(X)$-模结构由 $(f,m)\longmapsto\beta(f)m$（$f\in R(X)$，$m\in M$）给出。

现在想给出 $F^+(M)$ 的类似描述。先要把自同构 $\beta$ 扩充到 $D_X$。

设 $T$ 是 $X$ 上的微分算子，令 $\tilde\beta(T)(f)=\beta(T\alpha(f))$（$f\in R(X)$）。显然 $\tilde\beta(T)$ 是 $R(X)$ 的 $k$-线性自同态。此外 $T\longmapsto\tilde\beta(T)$ 是线性映射。又对 $D_X$ 中两个微分算子 $T$ 和 $S$，对一切 $f\in R(X)$ 有
\begin{align*}
\tilde\beta(TS)(f)&=\beta(TS\alpha(f))=\beta\bigl(T\alpha(\beta(S\alpha(f)))\bigr)\\
&=\beta\bigl(T\alpha(\tilde\beta(S)(f))\bigr)=\tilde\beta(T)\bigl(\tilde\beta(S)(f)\bigr),
\end{align*}
即 $\tilde\beta$ 是从 $k$-代数 $D_X$ 到 $R(X)$ 的 $k$-线性自同态代数的同态。由于对 $g\in R(X)$ 有
$$\tilde\beta(g)f=\beta(g\alpha(f))=\beta(g)f$$
（$f\in R(X)$），可见 $\tilde\beta$ 扩充了 $R(X)$ 的自同构 $\beta$。这进而蕴含 $\tilde\beta(T)\in D_X$（$T\in D_X$），即 $\tilde\beta$ 是 $D_X$ 的自同构，它扩充了 $R(X)$ 的自同构 $\beta$。因此可简单地仍记为 $\beta$。

设 $1\le i\le n$。则
\begin{align*}
\beta(\partial_i)(f)&=\beta(\partial_i\alpha(f))=\beta\bigl(\partial_i(f\circ F)\bigr)\\
&=\beta\Bigl(\sum_{j=1}^n\bigl((\partial_jf)\circ F\bigr)\partial_iF_j\Bigr)=\sum_{j=1}^n\bigl((\partial_iF_j)\circ G\bigr)\partial_jf=\Bigl(\sum_{j=1}^n\beta(\partial_iF_j)\partial_j\Bigr)(f).
\end{align*}

现在考虑与映射 $F$ 相伴的双模 $D_{X\to X}$。线性映射 $\varphi:f\otimes T\longmapsto\beta(f)T$ 把它与 $D_X$ 等同。由右乘给出的 $D_X$-模结构相同。另一方面，对任意 $T\in D_X$ 和 $1\le i\le n$ 有
$$\varphi\bigl(\partial_i(1\otimes T)\bigr)=\varphi\Bigl(\sum_{j=1}^n\partial_iF_j\otimes\partial_jT\Bigr)=\sum_{j=1}^n\beta(\partial_iF_j)\partial_jT=\beta(\partial_i)\varphi(1\otimes T).$$
因此双模 $D_{X\to X}$ 同构于 $D_X$，其右作用为右乘，左作用为 $\beta$ 与左乘的复合。这进而蕴含 $F^+(M)$ 同构于 $M$，其 $D_X$-模结构由 $(T,m)\longmapsto\beta(T)m$（$T\in D_X$，$m\in M$）给出。

因此由 3.10 得到如下结果。

\begin{lemma}
设 $F:X\longrightarrow X$ 是 $X$ 的同构。
\begin{enumerate}[label=(\roman*)]
  \item 设 $M$ 是 $D_X$-模。则作为线性空间 $F^+(M)$ 等于 $M$，其 $D_X$-作用由 $(T,m)\longmapsto\beta(T)m$（$T\in D_X$，$m\in M$）给出。
  \item 函子 $F^+:\mathcal{M}^L(D_X)\to\mathcal{M}^L(D_X)$ 是正合的。
  \item 函子 $F^+$ 把有限生成 $D_X$-模映为有限生成 $D_X$-模。若 $M$ 是有限生成 $D_X$-模，则 $d(F^+(M))=d(M)$。
\end{enumerate}
特别地，$F^+$ 把完整模映为完整模。
\end{lemma}

可通过描述有限生成 $D_X$-模 $M$ 的特征簇 $\operatorname{Ch}(F^+(M))$ 使上述陈述更精确。首先，由上述计算可见 $D_X$ 的自同构 $\beta$ 诱导 $\Gr D_X=k[X_1,X_2,\dots,X_n,\xi_1,\xi_2,\dots,\xi_n]$ 的自同构 $\Gr\beta$，它由
$$X_i\longmapsto\beta(X_i)=G_i,\qquad \xi_i\longmapsto\sum_{j=1}^n\beta(\partial_iF_j)\xi_j=\sum_{j=1}^n\bigl((\partial_iF_j)\circ G\bigr)\xi_j$$
（$1\le i\le n$）定义。

现在想用更几何的语言描述此构造。若 $x=(x_1,x_2,\dots,x_n)$ 是 $X=k^n$ 中的点，则通过 $df(x)\longmapsto\bigl((\partial_1f)(x),(\partial_2f)(x),\dots,(\partial_nf)(x)\bigr)$ 把 $x$ 处的余切空间 $T_x^*(X)$ 与 $k^n$ 等同。因此通过 $(x,df(x))\longmapsto\bigl(x_1,\dots,x_n,(\partial_1f)(x),\dots,(\partial_nf)(x)\bigr)$（$x\in X$）可把 $X$ 的余切丛 $T^*(X)$ 与 $k^{2n}$ 等同。设 $F:X\longrightarrow X$ 是 $X$ 的同构，$G$ 是其逆。则映射 $G$ 把 $X$ 中的点 $x$ 映为 $G(x)$，而 $F$ 把 $G(x)$ 映为 $x$。它们的微分 $T_x(G)$ 和 $T_{G(x)}(F)$ 是切空间 $T_x(X)$ 与 $T_{G(x)}(X)$ 之间互逆的线性同构。因此它们的伴随 $T_x(G)^*:T_{G(x)}^*(X)\to T_x^*(X)$ 和 $T_{G(x)}(F)^*:T_x^*(X)\to T_{G(x)}^*(X)$ 是互逆的线性同构。这进而蕴含可定义 $X$ 的余切丛 $T^*(X)$ 的同构 $\gamma$，由 $(x,\xi)\longmapsto\bigl(G(x),T_{G(x)}(F)^*\xi\bigr)$（$\xi\in T_x^*(X)$，$x\in X$）给出。若把 $T^*(X)$ 与 $k^{2n}$ 等同，则由考察上述公式可见 $(\Gr\beta)(P)=P\circ\gamma$（$P\in k[X_1,X_2,\dots,X_n,\xi_1,\xi_2,\dots,\xi_n]$）。

设 $M$ 是带好滤过 $F M$ 的有限生成 $D_X$-模。则可把 $F^+(M)$ 实现为按上述方式作用的 $M$。显然 $F M$ 是如上实现的 $F^+(M)$ 的好滤过。因此可把 $\Gr F^+(M)$ 与 $\Gr M$ 等同，其作用由 $(Q,m)\longmapsto(\Gr\beta)(Q)m$（$Q\in k[X_1,X_2,\dots,X_n,\xi_1,\xi_2,\dots,\xi_n]$，$m\in\Gr M$）给出。故 $Q$ 属于 $\Gr F^+(M)$ 的零化子当且仅当 $(\Gr\beta)(Q)$ 属于 $\Gr M$ 的零化子。若 $I$ 是 $\Gr F^+(M)$ 的零化子，则 $(\Gr\beta)I$ 是 $\Gr M$ 的零化子。因此 $(x,\xi)\in\operatorname{Ch}(F^+(M))$ 当且仅当 $\gamma^{-1}(x,\xi)\in\operatorname{Ch}(M)$。

\begin{lemma}
设 $M$ 是有限生成 $D_X$-模。则
$$\operatorname{Ch}(F^+(M))=\gamma(\operatorname{Ch}(M)).$$
\end{lemma}

最后，这允许给出逆像函子左同调维数的一个估计。

\begin{theorem}
设 $X=k^n$，$Y=k^m$，$F:X\longrightarrow Y$ 是多项式映射。则 $F^+$ 的左同调维数 $\le\dim Y$。
\end{theorem}

\begin{proof}
为证明此陈述，我们使用图像构造。设 $i:X\to X\times Y$ 是由 $i(x)=(x,0)$（$x\in X$）给出的映射。设 $\Phi:X\times Y\to X\times Y$ 是由 $\Phi(x,y)=(x,y+F(x))$（$x\in X$，$y\in Y$）给出的映射。最后，设 $p:X\times Y\to Y$ 是由 $p(x,y)=y$（$x\in X$，$y\in Y$）给出的投影。则 $F=p\circ\Phi\circ i$。此外，$\Phi$ 是 $X\times Y$ 的同构，其逆为 $(x,y)\longmapsto(x,y-F(x))$。

由 10.3，$F^+=i^+\circ\Phi^+\circ p^+$。此外由 10.5 和 10.8，函子 $p^+$ 和 $\Phi^+$ 正合。因此对一切 $q\in\mathbb{Z}$ 有 $L^qF^+=L^qi^+\circ\Phi^+\circ p^+$。由 10.7 断定，当 $q<-\dim Y$ 时 $L^qF^+=0$。
\end{proof}
\begin{annote}
本节以下六条（投影空间例子、$i_+$ 的同调维数、自同构情形的 $\beta$、特征簇的像、以及左同调维数 $\le\dim Y$ 的图像构造证明）在初版译本中整段缺失，本次校订据英文原书 p49–p52 补译。若与旧版页边编号不合，以此版为准。[来源：原书 §10.5–10.10, p49–p52；校订日志 chap1.md]
\end{annote}



\section{直像}

设 $X=k^n$，$Y=k^n$，$F:X\to Y$ 是如上节的多项式映射。与 $F$ 的复合定义自然环同态 $\widehat F:R(Y)\to R(X)$。此同态进而定义从 $R(X)$-模范畴到 $R(Y)$-模范畴的函子 $F_*$。对任意 $R(X)$-模 $M$，定义 $F_*(M)$ 为作为 $k$-线性空间等于 $M$ 的模，$R(Y)$ 的作用由 $(f,m)\longmapsto\widehat F(f)\cdot m$ 给出（$f\in R(Y)$，$m\in M$）。函子 $F_*:\mathcal{M}(R(X))\to\mathcal{M}(R(Y))$ 称为\emph{直像函子}。显然 $F_*$ 是正合函子。

遗憾的是，若 $M$ 是 $D_X$-模，直像 $F_*(M)$ 一般不允许 $D_Y$-模结构。例如考虑 $X=\{0\}$ 到 $Y=k$ 的包含 $i$，$D_X=R(X)=k$，$D_Y$ 是一变量多项式系数微分算子代数。$D_X$-模范畴就是 $k$-线性空间范畴。由 6.2，非零有限维 $D_X$-模 $M$ 的逆像不能有 $D_Y$-模结构。因此 $D$-模的直像不会像逆像那样与正则函数环上模的直像相关。

若对 $D_{X\to Y}$ 的两个作用都取转置，得到（左 $D_Y$，右 $D_X$）-双模 $D_{Y\leftarrow X}$。这允许对任意左 $D_X$-模 $M$ 定义左 $D_Y$-模
$$F_+(M)=D_{Y\leftarrow X}\otimes_{D_X}M.$$
显然 $F_+$ 是从 $\mathcal{M}^L(D_X)$ 到 $\mathcal{M}^L(D_Y)$ 的右正合函子。称之为\emph{直像函子}。$F_+$ 的左导出函子 $L^iF_+$ 对左 $D_X$-模 $M$ 由
$$L^{-j}F_+(M)=\operatorname{Tor}^{D_X}_j(D_{Y\leftarrow X},M)$$
给出。

设 $X=k^n$，$Y=k^m$，$Z=k^p$，$F:X\to Y$ 和 $G:Y\to Z$ 是多项式映射。若转置作用，10.4 蕴含如下陈述：
$$D_{Z\leftarrow X}=D_{Z\leftarrow Y}\otimes_{D_Y}D_{Y\leftarrow X},$$
以及
$$\operatorname{Tor}^{D_Y}_j(D_{Z\leftarrow Y},D_{Y\leftarrow X})=0\quad\text{对 }j\in\mathbb{N}.$$

若 $P$ 是射影左 $D_X$-模，$P\oplus Q=D_X^{(I)}$（对某左 $D_X$-模 $Q$ 和某 $I$）。因此 $F_+(P)\oplus F_+(Q)=F_+(D_X^{(I)})=(D_{Y\leftarrow X})^{(I)}$。这蕴含如下结果。

\begin{lemma}
设 $P$ 是射影左 $D_X$-模。则
$$\operatorname{Tor}^{D_Y}_j(D_{Z\leftarrow Y},F_+(P))=0\quad\text{对 }j\in\mathbb{N}.$$
\end{lemma}

\begin{theorem}
设 $X=k^n$，$Y=k^m$，$Z=k^p$，$F:X\to Y$ 和 $G:Y\to Z$ 是多项式映射。则
\begin{enumerate}[label=(\roman*)]
  \item 直像函子 $(G\circ F)_+$（从 $\mathcal{M}^L(D_X)$ 到 $\mathcal{M}^L(D_Z)$）同构于 $G_+\circ F_+$；
  \item 对任意左 $D_X$-模 $M$，存在谱序列，其 $E_2$ 项为 $E_2^{pq}=L^pG_+(L^qF_+(M))$，收敛于 $L^{p+q}(G\circ F)_+(M)$。
\end{enumerate}
\end{theorem}

\begin{proof}
(i) 对任意左 $D_X$-模 $M$，由 10.4.(i) 有
\begin{align*}
(G\circ F)_+(M)&=D_{Z\leftarrow X}\otimes_{D_X}M
=(D_{Z\leftarrow Y}\otimes_{D_Y}D_{Y\leftarrow X})\otimes_{D_X}M\\
&=D_{Z\leftarrow Y}\otimes_{D_Y}(D_{Y\leftarrow X}\otimes_{D_X}M)
=D_{Z\leftarrow Y}\otimes_{D_Y}F_+(M)=G_+(F_+(M)).
\end{align*}

(ii) 由 11.1，对任意射影 $D_X$-模 $P$，直像 $F_+(P)$ 是 $G_+$-无环的。因此结论由 Grothendieck 谱序列得出。
\end{proof}

现在考虑一个简单的例子。设 $i$ 为由 $i(x)=(x,0)$（$x\in X$）给出的 $X$ 到 $X\times Y$ 的典范注入。则
$$D_{X\to X\times Y}=i_+(D_{X\times Y})=i_+(D_X\boxtimesext D_Y)=D_X\boxtimesext D_Y/\bigl((y_1,y_2,\dots,y_m)D_Y\bigr)$$
以及
$$D_{X\times Y\leftarrow X}=D_X\boxtimesext D_Y/\bigl(D_Y(y_1,y_2,\dots,y_m)\bigr).$$
这意味着对任意左 $D_X$-模 $M$，
$$i_+(M)=M\boxtimesext D_Y/\bigl(D_Y(y_1,y_2,\dots,y_m)\bigr).$$
此外，模 $D_Y/(D_Y(y_1,y_2,\dots,y_m))$ 同构于 8.3 中讨论的 $\Delta_m$。

\begin{proposition}
设 $i:X\to X\times Y$ 是由 $i(x)=(x,0)$（$x\in X$）定义的注入。则
\begin{enumerate}[label=(\roman*)]
  \item $i_+$ 是从 $\mathcal{M}^L(D_X)$ 到 $\mathcal{M}^L(D_{X\times Y})$ 的正合函子；
  \item 对任意左 $D_X$-模 $M$，$i_+(M)=M\boxtimesext D_Y/\bigl(D_Y(y_1,y_2,\dots,y_m)\bigr)$；
  \item 若 $M$ 是有限生成 $D_X$-模，则 $i_+(M)$ 是有限生成 $D_{X\times Y}$-模；
  \item 对任意有限生成左 $D_X$-模 $M$，$d(i_+(M))=d(M)+m$。
\end{enumerate}
特别地，有限生成 $D_X$-模 $M$ 是完整的当且仅当 $i_+(M)$ 是完整的。
\end{proposition}

\begin{proof}
已证明 (ii)，它立即蕴含 (i)。如 8.3 中所述，$D_Y/(D_Y(y_1,y_2,\dots,y_m))$ 是不可约完整 $D_Y$-模，故 (iii) 由 9.2 得出。为证 (iv)，先注意由 9.3 有
$$d(i_+(M))=d(M)+d\bigl(D_Y/(D_Y(y_1,y_2,\dots,y_m))\bigr).$$
由于 $D_Y/(D_Y(y_1,y_2,\dots,y_m))$ 是完整的，其维数等于 $m$。
\end{proof}

现在研究投影 $p:X\times Y\to Y$（$p(x,y)=y$，$x\in X$，$y\in Y$）的直像。

先考虑 $\dim X=1$ 的情形。则
$$D_{X\times Y\to Y}=p_+(D_Y)=D_X/D_X(\partial_1)\boxtimesext D_Y.$$
故 $D_{Y\leftarrow X\times Y}=D_X/((\partial_1)D_X)\boxtimesext D_Y$。有正合列
$$0\longrightarrow D_{X\times Y}\xrightarrow{\ \partial_1\ }D_{X\times Y}\longrightarrow D_{Y\leftarrow X\times Y}\longrightarrow 0$$
（（左 $D_Y$，右 $D_{X\times Y}$）-双模），其中第二个箭头表示左乘 $\partial_1$。显然这是 $D_{Y\leftarrow X\times Y}$ 由自由右 $D_{X\times Y}$-模构成的左分解，故对任意 $D_{X\times Y}$-模 $M$，复形
$$\cdots\longrightarrow 0\longrightarrow M\xrightarrow{\ \partial_1\ }M\longrightarrow 0\longrightarrow\cdots$$
的上同调是 $\operatorname{Tor}^{D_{X\times Y}}_\cdot(D_{Y\leftarrow X\times Y},M)=L_\cdot p_+(M)$。由此可知 $L^q p_+(M)=0$ 对 $q\notin\{0,-1\}$。

因此建立了如下结果。

\begin{lemma}
设 $\dim X=1$，$p$ 是 $X\times Y$ 到 $Y$ 的典范投影。则对任意 $D_{X\times Y}$-模 $M$ 有
\begin{enumerate}[label=(\roman*)]
  \item $p_+(M)=\operatorname{coker}\partial_1$；
  \item $L^{-1}p_+(N)=\ker\partial_1$；
  \item $L^q p_+(N)=0$ 对不同于 $0$ 或 $-1$ 的 $q$。
\end{enumerate}
特别地，$p_+$ 的左同调维数 $\le 1$。
\end{lemma}

最后陈述对任意 $X$ 有如下推广。

\begin{lemma}
设 $p$ 是 $X\times Y$ 到 $Y$ 的典范投影。则 $p_+$ 的左同调维数 $\le\dim X$。
\end{lemma}

\begin{proof}
设 $X'=\{x_n=0\}\subset X$，以 $p'$ 记 $X'\times Y$ 到 $Y$ 的典范投影。又以 $p''$ 记 $X\times Y$ 到 $X'\times Y$ 的典范投影。则 $p=p'\circ p''$。故由 11.2.(ii)、11.4 和归纳假设得对 $q<-\dim X$ 有 $L^q p_+(M)=0$。
\end{proof}

设 $F:X\to X$ 是 $X$ 的同构，$G$ 是其逆。如 §10 中，定义 $D_X$ 的自同构 $\alpha$ 和 $\beta$。我们在那里把 $F$ 附带的双模 $D_{X\to X}$ 与 $D_X$ 等同，其作用是右乘和与 $\beta$ 复合的左乘。对它应用 $\alpha$，可见它也同构于 $D_X$，其作用是右乘与 $\alpha$ 复合和左乘。通过应用主反自同构，可见双模 $D_{X\leftarrow X}$ 同构于 $D_X$，其作用是与 $\alpha$ 复合的左乘和右乘。这进而蕴含对任意 $D_X$-模 $M$，直像 $F_+(M)$ 同构于 $M$，其作用由 $(T,m)\longmapsto\alpha(T)m$ 给出。特别地 $F_+(M)=G_+(M)$。

因此，由 10.8 立即推出如下结果。

\begin{lemma}
设 $F:X\to X$ 是 $X$ 的同构，$G:X\to X$ 是其逆。

(i) 设 $M$ 是 $D_X$-模。则 $F_+(M)$ 作为线性空间等于 $M$，$D_X$ 作用由 $(T,m)\longmapsto\alpha(T)m$ 给出（$T\in D_X$，$m\in M$）。

(ii) 函子 $F_+:\mathcal{M}^L(D_X)\to\mathcal{M}^L(D_X)$ 是正合的。

(iii) 函子 $F_+$ 把有限生成 $D_X$-模映为有限生成 $D_X$-模。若 $M$ 是有限生成 $D_X$-模，则 $d(F_+(M))=d(M)$。特别地 $F_+$ 把完整模映为完整模。
\end{lemma}

此外 $F_+=G_+$ 且 $F^+=G^+$，这些函子互为范畴的拟逆等价。

如上节，这允许给出直像函子左同调维数的估计。

\begin{theorem}
设 $X=k^n$，$Y=k^m$，$F:X\to Y$ 是多项式映射。则 $F_+$ 的左同调维数 $\le\dim X$。
\end{theorem}

\begin{proof}
如 10.10 的证明，使用图像构造。设 $i:X\to X\times Y$ 为由 $i(x)=(x,0)$（$x\in X$）给出的态射。设 $\Phi:X\times Y\to X\times Y$ 为由 $\Phi(x,y)=(x,y+F(x))$（$x\in X$，$y\in Y$）给出的态射。最后设 $p:X\times Y\to Y$ 为由 $p(x,y)=y$（$x\in X$，$y\in Y$）给出的投影。则 $F=p\circ\Phi\circ i$。此外 $\Phi$ 是 $X\times Y$ 的同构，其逆为 $(x,y)\longmapsto(x,y-F(x))$。

由 11.2，$F_+=p_+\circ\Phi_+\circ i_+$。此外由 11.3 和 11.6，函子 $i_+$ 和 $\Phi_+$ 是正合的。因此对所有 $q\in\mathbb{Z}$ 有 $L^qF_+=L^qp_+\circ\Phi_+\circ i_+$。由 11.5，对 $q<-\dim X$ 有 $L^qF_+=0$。
\end{proof}

\section{Kashiwara 定理}

设 $X=k^n$，$Y=\{x_n=0\}\subset X$。再令 $Z=\{x_1=x_2=\cdots=x_{n-1}=0\}\cong k$。于是 $X=Y\times Z$，这也意味着 $D_X=D_Y\boxtimesext D_Z$。设 $M$ 是 $D_X$-模，令
$$\Gamma_{[Y]}(M)=\{m\in M\mid x_n^p m=0\ \text{对某 }p\in\mathbb{N}\}.$$

\begin{lemma}
设 $M$ 是 $D_X$-模。则
\begin{enumerate}[label=(\roman*)]
  \item $\Gamma_{[Y]}(M)$ 是 $M$ 的 $D_X$-子模；
  \item $\operatorname{supp}(\Gamma_{[Y]}(M))\subset Y$；
  \item 若 $N$ 是 $M$ 的 $D_X$-子模且 $\operatorname{supp}(N)\subset Y$，则 $N\subset\Gamma_{[Y]}(M)$。
\end{enumerate}
\end{lemma}

\begin{proof}
(i) 设 $m\in\Gamma_{[Y]}(M)$。则对 $1\le i\le n$ 和 $1\le j<n$ 有 $x_i m\in\Gamma_{[Y]}(M)$、$\partial_j m\in\Gamma_{[Y]}(M)$。只需检验 $\partial_n m\in\Gamma_{[Y]}(M)$。对任意 $j\in\mathbb{N}$ 有
$$x_n^{j+1}\partial_n m=[x_n,\partial_n]m+\partial_n x_n m=-(j+1)x_n^j m+\partial_n x_n^{j+1}m.$$
故若 $x_n^j m=0$，则 $x_n^{j+1}\partial_n m=0$。

(ii) 若 $x\notin Y$，则 $x_n\notin\mathfrak{m}_x$，局部化 $\Gamma_{[Y]}(M)_x=0$。

(iii) 设 $N$ 是 $M$ 的 $D_X$-子模且 $\operatorname{supp}(N)\subset Y$。取 $m\in N$，记 $N'$ 为由 $m$ 生成的 $R(X)$-子模。则 $\operatorname{supp}(N')\subset Y$。由于 $N'$ 有限生成，由 4.2 其支集等于其在 $R(X)$ 中零化子 $I$ 决定的簇。由零点定理 $r(I)\supset(x_n)$。这意味着对某 $j\in\mathbb{N}$，$x_n^j$ 零化 $N'$，即 $m\in\Gamma_{[Y]}(M)$。
\end{proof}

因此 $\Gamma_{[Y]}(M)$ 是 $M$ 中支集含于 $Y$ 的最大 $D_X$-子模。

乘以 $x_n$ 定义了 $M$ 作为 $D_Y$-模的自同态。令
$$M_0=\ker x_n\subset\Gamma_{[Y]}(M),\qquad M_1=\operatorname{coker} x_n=M/x_nM.$$
记 $i$ 为 $Y$ 到 $X$ 的自然包含。如 10.6 所述，$i^+(M)=M_1$，$L^{-1}i^+(M)=M_0$，其余逆像全为零。

考虑双加性映射 $D_X\times M_0\to M$。显然它通过 $D_X\otimes_{D_Y}M_0\to M$ 分解。此外由 $M_0$ 的定义，后一态射在 $D_Xx_n\otimes_{D_Y}M_0$ 在 $D_X\otimes_{D_Y}M_0$ 中的像上为零。如 §11 中所述，
$$D_{X\leftarrow Y}=D_Y\boxtimesext D_Z/D_Zx_n=\bigoplus_{j=0}^{\infty}\partial_n^jD_Y.$$
因此上述态射诱导自然的 $D_X$-模态射
$$i_+(M_0)=D_{X\leftarrow Y}\otimes_{D_Y}M_0\longrightarrow M.$$
显然其像含于 $\Gamma_{[Y]}(M)$。容易检验这实际上是函子 $i_+\circ L^{-1}i^+$ 到 $\Gamma_{[Y]}$ 的态射。

本节的关键结果是如下引理。

\begin{lemma}
态射 $i_+(M_0)\to\Gamma_{[Y]}(M)$ 是 $D_X$-模的同构。
\end{lemma}

\begin{proof}
先证此态射是满射。断言对任意 $p\in\mathbb{N}$，
$$\{m\in M\mid x_n^p m=0\}\subset D_X\cdot M_0.$$
$p=1$ 时显然。若 $p>1$ 且 $x_n^p m=0$，则
$$0=\partial_n(x_n^p m)=x_n^{p-1}(pm+x_n\partial_n m),$$
由归纳假设 $pm+x_n\partial_n m\in D_X\cdot M_0$。此外由归纳假设 $x_n m\in D_X\cdot M_0$。这意味着
$$(p-1)m=pm+[x_n,\partial_n]m=pm+x_n\partial_n m-\partial_n x_n m\in D_X\cdot M_0,$$
故 $m\in D_X\cdot M_0$。因此映射是满射。

现证单射。由上述讨论
$$i_+(M_0)=D_{X\leftarrow Y}\otimes_{D_Y}M_0=\bigoplus_{j=0}^{\infty}\partial_n^j M_0.$$
设 $(m_0,\partial_n m_1,\dots,\partial_n^q m_q,0,\dots)$ 是此直和中映为 $0$ 的非零元，且 $q$ 最小，即
$$m_0+\partial_n m_1+\cdots+\partial_n^q m_q=0.$$
则
$$0=x_n\sum_{j=0}^q\partial_n^j m_j=\sum_{j=1}^q[x_n,\partial_n^j]m_j=-\sum_{j=1}^q j\,\partial_n^{j-1}m_j,$$
矛盾。因此映射的核为零。
\end{proof}

\begin{corollary}
$x_n\Gamma_{[Y]}(M)=\Gamma_{[Y]}(M)$。
\end{corollary}

\begin{proof}
由 12.2，$\Gamma_{[Y]}(M)$ 的任意元素形如 $\sum_{j\in\mathbb{Z}_+}\partial_n^j m_j$，其中 $m_j\in M_0$。另一方面，
$$x_n\sum_{j\in\mathbb{Z}_+}\frac{1}{j+1}\partial_n^{j+1}m_j=-\sum_{j\in\mathbb{Z}_+}\partial_n^j m_j.$$
\end{proof}

\begin{corollary}
设 $M$ 是 $D_X$-模。则
\begin{enumerate}[label=(\roman*)]
  \item $\Gamma_{[Y]}(M)$ 是有限生成 $D_X$-模当且仅当 $M_0$ 是有限生成 $D_Y$-模；
  \item $d(\Gamma_{[Y]}(M))=d(M_0)+1$。
\end{enumerate}
特别地，$\Gamma_{[Y]}(M)$ 是完整的当且仅当 $L^{-1}i^+(M)=M_0$ 是完整的。
\end{corollary}

\begin{proof}
(i) 由 12.2 和 11.3.(iii)，若 $M_0$ 有限生成则 $\Gamma_{[Y]}(M)$ 有限生成。反之设 $\Gamma_{[Y]}(M)$ 是有限生成 $D_X$-模。设 $N_j$（$j\in\mathbb{N}$）是 $M_0$ 的 $D_Y$-子模的递增序列。则它们生成 $\Gamma_{[Y]}(M)$ 的 $D_X$-子模 $i_+(N_j)=\bigoplus_{p=0}^{\infty}\partial_n^p N_j$。由于 $\Gamma_{[Y]}(M)$ 有限生成，递增序列 $i_+(N_j)$ 稳定。此外 $N_j$ 是 $i_+(N_j)$ 中 $x_n$ 的核，故序列 $N_j$ 也必稳定。因此 $M_0$ 有限生成。

(ii) 由 12.2 和 11.3.(iv) 得出。
\end{proof}

\begin{corollary}
设 $M$ 是完整 $D_X$-模。则 $M_0$ 是完整 $D_Y$-模。
\end{corollary}

\begin{proof}
若 $M$ 是完整的，则 $\Gamma_{[Y]}(M)$ 也是完整的。故结论由 12.4 得出。
\end{proof}

设 $\mathcal{M}_Y(D_X)$ 为 $\mathcal{M}(D_X)$ 中由支集含于 $Y$ 的 $D_X$-模构成的全子范畴。以 $\mathcal{M}_{fg,Y}(D_X)$ 和 $\operatorname{Hol}_Y(D_X)$ 分别记支集含于 $Y$ 的有限生成、完整 $D_X$-模的相应子范畴。则由 12.1，对 $\mathcal{M}_Y(D_X)$ 中任意 $M$ 有 $M=\Gamma_{[Y]}(M)$。由 10.6 和 12.3 可见对 $\mathcal{M}_Y(D_X)$ 中任意 $M$ 有 $i^+(M)=0$，故 $L^{-1}i^+$ 是从 $\mathcal{M}_Y(D_X)$ 到 $\mathcal{M}(D_Y)$ 的正合函子。另一方面，$i_+$ 定义反方向的正合函子，且由 12.2 复合 $i_+\circ L^{-1}i^+$ 同构于 $\mathcal{M}_Y(D_X)$ 上的恒等函子。也显然 $L^{-1}i^+\circ i_+$ 同构于 $\mathcal{M}(D_Y)$ 上的恒等函子。

这导出如下基本结果。

\begin{theorem}[Kashiwara]
直像函子 $i_+$ 定义范畴 $\mathcal{M}(D_Y)$（相应地 $\mathcal{M}_{fg}(D_Y)$、$\operatorname{Hol}(D_Y)$）与范畴 $\mathcal{M}_Y(D_X)$（相应地 $\mathcal{M}_{fg,Y}(D_X)$、$\operatorname{Hol}_Y(D_X)$）的等价。其逆为函子 $L^{-1}i^+$。
\end{theorem}

\begin{proof}
只需证明括号中的陈述，它们由 12.4 立即得出。
\end{proof}
\begin{annote}
完整模经真态射的直像仍完整（Kashiwara），这是“完整”一词的稳定性来源：它使完整模在六种运算下有封闭性，从而可以谈论 $D^b_{rh}(X)$ 与 $D^b_c(X^{an})$ 的等价（Riemann–Hilbert）。1 维情形（$\dim X=1$）就是 Deligne 关于正则奇点线性微分方程有解的存在性定理。[来源：nLab 条目 Riemann-Hilbert+correspondence（Last revised 2024-08-09）、holonomic+D-module（2014-02-16）；原书 §12, p51–54]
\end{annote}



\section{完整性的保持}

本节证明直像和逆像保持完整模。先给出一个完整性的简单判别法。

设 $X=k^n$，$Y=k^m$，$F:X\to Y$ 是多项式映射。想研究完整模在逆像和直像函子作用下的行为。

首先再次使用图像构造把问题归约到特殊映射。如 10.10 和 11.7 的证明：
$$\begin{tikzcd}
X \arrow[r,"F"] \arrow[d,"i"' ] & Y \arrow[d,"p"]\\
X\times Y \arrow[r,"\Phi"'] & X\times Y
\end{tikzcd}$$
其中 $i(x)=(x,0)$（$x\in X$），$p(x,y)=y$（$x\in X$，$y\in Y$），$\Phi(x,y)=(x,y+F(x))$（$x\in X$，$y\in Y$）。

由 10.5 知 $p^+$ 是正合的且把完整模映为完整模。由 11.3 知 $i_+$ 是正合的且把完整模映为完整模。此外由 10.8 和 11.6 知 $\Phi^+$ 和 $\Phi_+$ 是正合的且把完整模映为完整模。

因此只需研究 $i^+$ 和 $p_+$ 的导出函子。先讨论浸入 $i:X\to X\times Y$。

\begin{lemma}
设 $N$ 是完整 $D_{X\times Y}$-模。则 $D_X$-模 $L^qi^+(N)$（$q\in\mathbb{Z}$）是完整的。
\end{lemma}

由于由 8.1.(ii)，完整模的子模、商模和扩张都是完整的，如 10.7 的证明用谱序列论证可把证明归约到 $\dim Y=1$ 的情形。此时若以 $y$ 记 $Y$ 上自然坐标，考虑 $D_X$-模态射 $N\xrightarrow{\ y\ }N$，则有 $i^+(N)=\operatorname{coker}y$，$L^{-1}i^+(N)=\ker y$，其余导出逆像全为零（如 10.6 所述）。此外若 $N$ 是完整的，由 12.4 知 $L^{-1}i_+(N)$ 是完整的。故只需处理 $i^+(N)$。

\begin{lemma}
设 $N$ 是完整 $D_{X\times Y}$-模。则 $i^+(N)$ 是完整的。
\end{lemma}

\begin{proof}
令 $\bar N=N/\Gamma_{[X]}(N)$。可考虑短正合列
$$0\longrightarrow\Gamma_{[X]}(N)\longrightarrow N\longrightarrow\bar N\longrightarrow 0.$$
由于 $i^+$ 是右正合的，这导出正合列
$$i^+(\Gamma_{[X]}(N))\longrightarrow i^+(N)\longrightarrow i^+(\bar N)\longrightarrow 0.$$
另一方面，由 12.3 知 $i^+(\Gamma_{[X]}(N))=0$。因此自然映射 $i^+(N)\to i^+(\bar N)$ 是同构。

设 $\bar v\in\Gamma_{[X]}(\bar N)\subset\bar N$，以 $v\in N$ 记 $\bar v$ 的代表。则对充分大的 $p\in\mathbb{Z}_+$ 有 $y^p\bar v=0$。因此 $y^p v\in\Gamma_{[X]}(N)$。这进而蕴含对充分大的 $q\in\mathbb{Z}_+$ 有 $y^{p+q}v=y^q(y^pv)=0$。故 $v\in\Gamma_{[X]}(N)$，$\bar v=0$。由此可知 $\Gamma_{[X]}(\bar N)=0$。

此外，若 $N$ 是完整 $D_{X\times Y}$-模，则 $\bar N$ 是完整 $D_{X\times Y}$-模。

因此可从一开始假设 $\Gamma_{[X]}(N)=0$。这意味着乘以 $y$ 在 $N$ 上是单射，$N$ 嵌入其局部化 $N_y$。考虑正合列
$$0\longrightarrow N\longrightarrow N_y\longrightarrow L\longrightarrow 0.$$
由于 $N$ 是完整 $D_{X\times Y}$-模，由 8.6 知 $N_y$ 是完整的。故 $L$ 是完整 $D_{X\times Y}$-模。由上述讨论，这意味着 $L^{-1}i^+(L)$ 是完整 $D_X$-模。

对短正合列应用 $i$ 的逆像长正合列，得到
$$\cdots\to L^{-1}i^+(N_y)\to L^{-1}i^+(L)\to i^+(N)\to i^+(N_y)\to i^+(L)\to 0.$$
由于乘以 $y$ 在 $N_y$ 上可逆，由 10.6 知
$$i^+(N_y)=L^{-1}i^+(N_y)=0.$$
故 $i^+(N)\cong L^{-1}i^+(L)$。由上述讨论断定 $i^+(N)$ 是完整 $D_X$-模。
\end{proof}

因此，由 10.3 得到如下结果。

\begin{theorem}
设 $F:X\to Y$ 是多项式映射，$M$ 是完整 $D_Y$-模。则 $L^qF^+(M)$（$q\in\mathbb{Z}$）是完整 $D_X$-模。
\end{theorem}

现在研究 $p$ 的直像。

\begin{lemma}
设 $M$ 是完整 $D_{X\times Y}$-模。则 $D_Y$-模 $L^qp_+(M)$（$q\in\mathbb{Z}$）是完整的。
\end{lemma}

\begin{proof}
由于由 8.1.(ii)，完整模的子模、商模和扩张都是完整的，如 11.5 的证明用谱序列论证可把证明归约到 $\dim X=1$ 的情形。此时若以 $\partial$ 记关于 $X$ 上坐标 $x$ 的导数，考虑 $D_Y$-模态射 $M\xrightarrow{\ \partial\ }M$，则有 $p_+(M)=\operatorname{coker}\partial$，$L^{-1}p_+(M)=\ker\partial$，其余导出逆像全为零（如 11.4 所述）。通过应用 Fourier 变换得到计算 $F(L_\cdot p_+(M))$ 的复形
$$\cdots\longrightarrow 0\longrightarrow F(M)\xrightarrow{\ x\ }F(M)\longrightarrow 0\longrightarrow\cdots.$$
由 13.2 证明中的论证，可见此复形计算典范包含 $j:Y\to X\times Y$（$j(y)=(0,y)$，$y\in Y$）的逆像。故由上同调由 13.2 是完整的。由 6.4，对所有 $q\in\mathbb{Z}$ 有 $L^q p_+(M)$ 是完整的。
\end{proof}

因此，由 11.2 得到如下结果。

\begin{theorem}
设 $F:X\to Y$ 是多项式映射，$M$ 是完整 $D_X$-模。则 $L^qF_+(M)$（$q\in\mathbb{Z}$）是完整 $D_Y$-模。
\end{theorem}

\begin{remark}
13.4 和 13.5 关于有限生成模的类似陈述不成立。例如令 $X=\{0\}$，$Y=k$，以 $i:X\to Y$ 记自然包含，则逆像 $i^+(D_Y)$ 是 $k$ 上无限维向量空间。类似地，若 $p$ 是 $Y$ 到一点的投影，则 $p_+(D_Y)$ 是 $k$ 上无限维向量空间。
\end{remark}
